% Merged by scripts/merge_body_inline.py (2026-11-17): Bodong's body_en.tex with the report-side responses;
% each table and figure follows the response that first cites it.
\responsetoreviewer{Response to Reviewer 1}
\reviewersection{Reviewer 1, Comment 1 (R1.1)}\phantomsection\label{rev:R1.1}
\begin{reviewercomment}
Independence of image-level evaluations. Segmentation and classification results are largely reported at the image level. The authors should clarify patient-level and examination-level splitting, deduplication across datasets, and whether multiple frames or views from the same patient could occur across splits. Patient-level reanalysis and cluster-aware confidence intervals would strengthen confidence in the reported performance.
\end{reviewercomment}
\begin{authorresponse}
We thank the reviewer for this important point. Our statistical reporting follows cluster-bootstrap conventions with patient-level aggregation (\responsetableref{tab_cls_ext}, \responsetableref{tab_seg_ext}, and \responsetableref{tab_patient_cls}). Near-duplicate screening (dHash, Hamming $\leq$ 5) of each foundational public cohort train split against the ZJH-8K image pool is in \responsetableref{tab_dhash}.\par
We clarify whether multiple frames or views from the same patient appear across splits. External test cohorts were held out from stacked expert training (\responsefigureref{fig_image_dataflow}). Where reproducible patient IDs exist (ZJH-8K, ThyroidXL), all images from the same patient remain in a single split---we do not assign different views of the same patient to both development and external test. Multiple images per patient may appear within the same external test set; inferential summaries use patient-level clusters. Patient-/examination-level keys for institutional cohorts follow the same patient-cluster split and bootstrap protocol as Reviewer 2, Comment \hyperref[rev:R2.2]{2} (\hyperref[rev:R2.2]{R2.2}).\par
We report patient-level re-analysis and cluster-bootstrap 95\% CIs as the reviewer suggests. Image-level external metrics with cluster-bootstrap intervals are in \responsetableref{tab_cls_ext} and \responsetableref{tab_seg_ext} (resampling protocol in Reviewer 2, Comment \hyperref[rev:R2.2]{2}, \hyperref[rev:R2.2]{R2.2}). Where reproducible patient IDs exist (ZJH-8K, ThyroidXL), patient-level ThyroidXAgent classification and segmentation with cluster-bootstrap 95\% CIs are in \responsetableref{tab_patient_cls} (cohort counts in \responsetableref{tab_cohort_patient_counts}). For bootstrap and these summaries, patient-level clusters apply when IDs are available; cohorts without IDs use one image per cluster (full resampling steps in \hyperref[rev:R2.2]{R2.2}).\par
Independence and near-duplicates. Experts are trained on stacked public cohorts (not on ZJH-8K); \responsetableref{tab_dhash} screens each training partition against the full ZJH-8K image pool (7958 images). Exact dHash matches (H=0) concentrate between TN3K train and ZJH-8K (same centre; 1154 H=0 pairs, 1151 unique ZJH images). TN5K, ThyroidXL, and PKTN train partitions show no H=0 matches to ZJH-8K. Metrics in this letter exclude those 1151 ZJH images (6807 images remain for external test; see \responsetableref{tab_cls_ext} and related tables).\par
\par
\end{authorresponse}
% [SYNC] tab_cls_performance.tex (manuscript: 26_Nature_ThyroidAgent/tables/); letter \label tab_cls_ext via \responsetabletitle.
\begin{manuscriptresponsetable}
\caption{\responsetabletitle{tab_cls_ext}{Cross-dataset generalization --- benign--malignant classification}}
\begin{adjustbox}{max width=\textwidth,center}
\begin{tabular}{lccccc}
\toprule
\textbf{Method} & \textbf{TN3K} & \textbf{ThyroidXL} & \textbf{TN5K} & \textbf{DDTI} & \textbf{ZJH-8K} \\
\midrule
\multicolumn{6}{l}{\textit{AUROC $\uparrow$}} \\
\midrule
ResNet-50 & $0.7674 \pm 0.0394$ & $0.9044 \pm 0.0118$ & $0.9322 \pm 0.0168$ & $0.6704 \pm 0.0842$ & $0.8988 \pm 0.0161$ \\
RepViT & $0.5556 \pm 0.0463$ & $0.7774 \pm 0.0188$ & $0.6603 \pm 0.0375$ & $0.6162 \pm 0.0804$ & $0.8246 \pm 0.0217$ \\
LSNet & $0.8095 \pm 0.0333$ & $0.9178 \pm 0.0114$ & $0.9091 \pm 0.0201$ & $0.7581 \pm 0.0658$ & $0.8377 \pm 0.0195$ \\
UltraFedFM & $0.8461 \pm 0.0697$ & $0.9239 \pm 0.0104$ & $0.9298 \pm 0.0175$ & $0.7518 \pm 0.1712$ & $0.8621 \pm 0.0205$ \\
MedGemma & $0.8492 \pm 0.0305$ & $0.9371 \pm 0.0095$ & $0.9442 \pm 0.0156$ & \textbf{0.8255 $\pm$ 0.0650} & $0.9011 \pm 0.0365$ \\
Qwen3-VL-8B-Instruct & $0.8237 \pm 0.0328$ & $0.9050 \pm 0.0115$ & $0.9214 \pm 0.0187$ & $0.7361 \pm 0.0692$ & $0.8397 \pm 0.0338$ \\
GPT-5 & $0.6924 \pm 0.0421$ & $0.7059 \pm 0.0469$ & $0.7737 \pm 0.0996$ & $0.6346 \pm 0.0914$ & $0.6317 \pm 0.0732$ \\
Gemini-2.5-Pro & $0.6587 \pm 0.0455$ & $0.6246 \pm 0.0640$ & $0.6873 \pm 0.0691$ & $0.6156 \pm 0.1308$ & $0.6602 \pm 0.0626$ \\
\responserowagent
ThyroidXAgent & \textbf{0.8714 $\pm$ 0.0299} & \textbf{0.9680 $\pm$ 0.0066} & \textbf{0.9584 $\pm$ 0.0152} & $0.8032 \pm 0.0641$ & \textbf{0.9190 $\pm$ 0.0214} \\
\midrule
\multicolumn{6}{l}{\textit{AUPRC $\uparrow$}} \\
\midrule
ResNet-50 & $0.6882 \pm 0.0632$ & $0.8882 \pm 0.0174$ & $0.9674 \pm 0.0268$ & $0.3755 \pm 0.1176$ & $0.9025 \pm 0.0212$ \\
RepViT & $0.4275 \pm 0.0528$ & $0.7161 \pm 0.0276$ & $0.8364 \pm 0.0411$ & $0.3924 \pm 0.0933$ & $0.8703 \pm 0.0239$ \\
LSNet & $0.7581 \pm 0.0452$ & $0.9040 \pm 0.0142$ & $0.9551 \pm 0.0134$ & $0.4180 \pm 0.1410$ & $0.8657 \pm 0.0243$ \\
UltraFedFM & $0.8531 \pm 0.0284$ & $0.9354 \pm 0.0114$ & $0.8422 \pm 0.0421$ & $0.4487 \pm 0.1452$ & $0.9027 \pm 0.0203$ \\
MedGemma & $0.8047 \pm 0.0430$ & $0.9201 \pm 0.0139$ & $0.9747 \pm 0.0084$ & $0.5537 \pm 0.1663$ & $0.9079 \pm 0.0114$ \\
Qwen3-VL-8B-Instruct & $0.7617 \pm 0.0511$ & $0.8787 \pm 0.0379$ & $0.9636 \pm 0.0106$ & $0.4112 \pm 0.1415$ & $0.8811 \pm 0.0241$ \\
GPT-5 & $0.6627 \pm 0.0633$ & $0.6237 \pm 0.0666$ & $0.8920 \pm 0.0316$ & $0.3578 \pm 0.1089$ & $0.8120 \pm 0.0541$ \\
Gemini-2.5-Pro & $0.6205 \pm 0.0587$ & $0.4914 \pm 0.0841$ & $0.8462 \pm 0.0446$ & $0.3924 \pm 0.1527$ & $0.8197 \pm 0.0613$ \\
\responserowagent
ThyroidXAgent & \textbf{0.8544 $\pm$ 0.0415} & \textbf{0.9432 $\pm$ 0.0268} & \textbf{0.9763 $\pm$ 0.0079} & \textbf{0.7463 $\pm$ 0.1080} & \textbf{0.9117 $\pm$ 0.0176} \\
\bottomrule
\end{tabular}
\end{adjustbox}
\responsetablenote{AUROC and AUPRC with cluster-bootstrap 95\% CIs (1000 resamples, seed~42); $\pm$ gives point estimate and half CI width. ZJH-8K: 6807 images, 1068 patients (TN3K train dHash H=0 exclusion; \responsetableref{tab_dhash}).}
\end{manuscriptresponsetable}
% [SYNC] tab_seg_performance.tex (manuscript: 26_Nature_ThyroidAgent/tables/); letter \label tab_seg_ext. ZJH ThyroidXAgent = routing ablation strategy (7), aligned with tab_routing_ablation_seg.
\begin{manuscriptresponsetable}
\caption{\responsetabletitle{tab_seg_ext}{Cross-dataset generalization --- nodule segmentation}}
\begin{adjustbox}{max width=\textwidth,center}
\begin{tabular}{lccccccc}
\toprule
\textbf{Model} & \textbf{TN3K} & \textbf{ThyroidXL} & \textbf{PKTN} & \textbf{TN5K} & \textbf{DDTI} & \textbf{ZJH-8K} & \textbf{RJH-7K} \\
\midrule
\multicolumn{8}{l}{\textit{Dice (\%) $\uparrow$}} \\
\midrule
TransUnet & $81.84 \pm 1.62$ & $85.75 \pm 0.57$ & $76.89 \pm 3.56$ & $78.54 \pm 1.51$ & $76.58 \pm 1.67$ & $76.53 \pm 0.86$ & $84.83 \pm 0.37$ \\
MedSegX & $83.93 \pm 0.79$ & $79.98 \pm 0.36$ & $80.63 \pm 0.42$ & $83.10 \pm 0.48$ & $75.12 \pm 1.68$ & $69.21 \pm 1.00$ & $85.40 \pm 0.18$ \\
MedSAM2 & $84.47 \pm 1.02$ & $86.94 \pm 0.36$ & $83.46 \pm 2.60$ & $83.03 \pm 1.29$ & $84.72 \pm 1.26$ & $82.57 \pm 0.82$ & $90.72 \pm 0.21$ \\
UltraFedFM & $81.18 \pm 1.46$ & $84.70 \pm 0.53$ & $75.31 \pm 1.12$ & $77.13 \pm 1.38$ & $75.57 \pm 1.67$ & $77.56 \pm 0.78$ & $83.10 \pm 0.33$ \\
\responserowagent
ThyroidXAgent & \textbf{85.69 $\pm$ 1.17} & \textbf{88.53 $\pm$ 0.42} & \textbf{85.18 $\pm$ 2.03} & \textbf{83.88 $\pm$ 1.28} & \textbf{85.73 $\pm$ 1.10} & \textbf{88.30 $\pm$ 0.34} & \textbf{91.77 $\pm$ 0.28} \\
\midrule
\multicolumn{8}{l}{\textit{HD95 (mm) $\downarrow$}} \\
\midrule
TransUnet & $27.27 \pm 5.52$ & $22.42 \pm 1.34$ & $26.88 \pm 9.66$ & $22.32 \pm 3.43$ & $17.12 \pm 1.55$ & $16.24 \pm 0.77$ & $18.81 \pm 0.74$ \\
MedSegX & $10.95 \pm 0.64$ & $11.07 \pm 0.32$ & $10.83 \pm 0.70$ & $11.76 \pm 0.76$ & $18.39 \pm 1.65$ & $23.48 \pm 1.04$ & $9.37 \pm 0.18$ \\
MedSAM2 & $11.51 \pm 1.53$ & $5.46 \pm 0.44$ & $10.56 \pm 3.64$ & $10.94 \pm 1.12$ & $10.06 \pm 1.21$ & $13.51 \pm 0.83$ & $2.92 \pm 0.17$ \\
UltraFedFM & $14.98 \pm 2.10$ & $8.10 \pm 0.58$ & $16.08 \pm 1.67$ & $14.96 \pm 1.65$ & $18.12 \pm 1.47$ & $18.57 \pm 0.83$ & $9.06 \pm 0.38$ \\
\responserowagent
ThyroidXAgent & \textbf{9.42 $\pm$ 1.45} & \textbf{5.16 $\pm$ 0.53} & \textbf{9.15 $\pm$ 3.28} & \textbf{7.14 $\pm$ 1.23} & \textbf{9.24 $\pm$ 1.06} & \textbf{10.25 $\pm$ 0.39} & \textbf{1.88 $\pm$ 0.09} \\
\bottomrule
\end{tabular}
\end{adjustbox}
\responsetablenote{Dice and 95th-percentile Hausdorff distance (HD95) with cluster-bootstrap 95\% CIs (1000 resamples, seed~42, image-level); $\pm$ gives point estimate and half CI width. ZJH-8K: 6807 images, 1068 patients (\responsetableref{tab_dhash}).}
\end{manuscriptresponsetable}
% [SYNC] tab_patient_level.tex (manuscript: 26_Nature_ThyroidAgent/tables/); letter \label tab_patient_cls via \responsetabletitle.
% ZJH: docs/revision/b2/tables/zjh8k_b2_excl_tn3k_h0.json patient_level; ThyroidXL: thyroidxl_b2.json.
\begin{responsetable}
\caption{\responsetabletitle{tab_patient_cls}{ThyroidXAgent patient-level classification and segmentation (ZJH-8K, ThyroidXL)}}
\begin{adjustbox}{max width=\textwidth,center}
\small
\begin{tabular}{lcccccccccc}
\toprule
Cohort & $n$ patients & AUROC & 95\% CI & AUPRC & 95\% CI & Dice & 95\% CI & HD95 (mm) & 95\% CI \\
\midrule
ZJH-8K & 1068 & \textbf{0.9634} & 0.9513--0.9740 & \textbf{0.9730} & 0.9623--0.9827 & 0.8270 & 0.8179--0.8353 & 13.14 & 12.30--14.04 \\
ThyroidXL & 739 & \textbf{0.9484} & 0.9326--0.9628 & \textbf{0.9355} & 0.9064--0.9592 & 0.8749 & 0.8675--0.8816 & 5.68 & 5.13--6.30 \\
\bottomrule
\end{tabular}
\end{adjustbox}
\responsetablenote{ZJH-8K and ThyroidXL official test splits. Classification: arithmetic mean of malignancy probability within cluster; patient label = within-cluster max. Segmentation: within-patient arithmetic mean of image-level Dice and HD95. ZJH-8K: 6807 images, 1068 patients (\responsetableref{tab_cohort_patient_counts}); patient-level metrics on the same TN3K dHash H=0 exclusion subset as \responsetableref{tab_cls_ext} and \responsetableref{tab_seg_ext} (cluster-bootstrap 1000 resamples, seed=42). ThyroidXL: 2090 images, 739 patients.}
\end{responsetable}
\begin{responsetable}
\caption{\responsetabletitle{tab_dhash}{dHash cross-cohort near-duplicate screening}}
\begin{adjustbox}{max width=\textwidth,center}
\begin{tabular}{lccccccc}
\toprule
Training partition (A) & Comparator (B) & n (A) & n (B) & Pairs (H$\leq$5) & H=0 pairs & Near only (1--5) & Unique ZJH (H=0) \\
\midrule
TN3K train & ZJH-8K & 4633 & 7958 & 1290 & 1154 & 136 & 1151 \\
TN5K train & ZJH-8K & 3500 & 7958 & 2 & 0 & 2 & 0 \\
ThyroidXL train & ZJH-8K & 9441 & 7958 & 1 & 0 & 1 & 0 \\
PKTN train & ZJH-8K & 703 & 7958 & 1 & 0 & 1 & 0 \\
\bottomrule
\end{tabular}
\end{adjustbox}
\responsetablenote{64-bit dHash; Hamming $\leq$ 5. A = each foundational public cohort's official train split. B = ZJH-8K image pool (7958 images; full pool for screening). H=0 pairs = undirected train--ZJH image pairs with exact hash match; Unique ZJH (H=0) = distinct ZJH images with at least one H=0 match to A. Only the TN3K train row shows H=0 overlap with ZJH-8K (1151 unique ZJH images); main external-test metrics exclude those images (6807 images, 1068 patients; counts in \responsetableref{tab_cohort_patient_counts}). TN5K, ThyroidXL, and PKTN train show no H=0 matches to ZJH. Malignant-subtask development/hold-out: \responsetableref{tab_malig_oof}, \responsetableref{tab_malig_ext}.}
\end{responsetable}
\begin{responsefigure}
\includegraphics[width=\textwidth]{image_analysis_dataflow.pdf}
\caption{\responsefiguretitle{fig_image_dataflow}{Image-analysis data flow. Left: development pools for main tasks (segmentation and benign--malignant classification) and malignant subtasks (LNM, PTC subtype). Centre: shared image-analysis pipeline (DINO experts, PyRadiomics/AutoGluon tabular expert, validation-fitted routing, and routed segmentation masks and classification labels). Right: external-test cohorts held out from stacked OpenThyroid expert training (official test splits and independent held-out sources as labelled on the figure) and the dual role of ZJH-8K (segmentation/classification external test versus LNM/PTC-subtype development).}}
\label{fig_image_dataflow}
\end{responsefigure}
% Cohort image/patient counts (tab_cohort_patient_counts); sources: b2/cluster_maps, b1 excl manifest, b9/zjh_malignant_patient_split.json, ZJH-8K_BM_labels.json

\begin{responsetable}
\caption{\responsetabletitle{tab_cohort_patient_counts}{Image and patient counts for ZJH-8K and ThyroidXL}}
\begin{adjustbox}{max width=\textwidth,center}
\small
\begin{tabular}{@{}llrrl@{}}
\toprule
Cohort & Split / subset & Images & Patients & Role in this letter \\
\midrule
ZJH-8K & Full pool (dHash screening) & 7958 & 1146 & \responsetableref{tab_dhash} \\
ZJH-8K & External test (exclude TN3K train H=0) & 6807 & 1068 & \responsetableref{tab_cls_ext}, \responsetableref{tab_seg_ext}, \responsetableref{tab_patient_cls} \\
ZJH-8K & External subset, patient label benign & 2774 & 421 & \responsetableref{tab_patient_cls} \\
ZJH-8K & External subset, patient label malignant & 4033 & 647 & \responsetableref{tab_patient_cls} \\
ZJH-8K & Malignant primary-lesion development (LNM/FTC shared) & 3227 & 517 & \responsetableref{tab_malig_oof} five-fold OOF \\
ZJH-8K & Malignant primary-lesion hold-out & 801 & 130 & \responsetableref{tab_malig_ext} \\
\midrule
ThyroidXL & train (OpenThyroid development) & 9441 & 3353 & Stacked expert fine-tuning; not R4 external test \\
ThyroidXL & test (official external) & 2090 & 739 & \responsetableref{tab_patient_cls} \\
\bottomrule
\end{tabular}
\end{adjustbox}
\responsetablenote{Benign/malignant rows on the ZJH external subset use patient label = within-cluster max of image-level malignancy labels; all images from that patient are counted in the matching row (same as \responsetableref{tab_patient_cls}). ZJH-8K: 7958 images and 1146 patients in the full institutional pool used for dHash screening; 6807 images and 1068 patients after excluding TN3K-train H=0 duplicates for segmentation/classification external test (Methods). External test 6807 images = full pool minus the 1151 unique ZJH images in the TN3K train row of \responsetableref{tab_dhash}. Malignant hold-out: patient-level 20\% split (seed 42), shared by LNM and PTC subtype.}
\end{responsetable}



\reviewersection{Reviewer 1, Comment 2 (R1.2)}\phantomsection\label{rev:R1.2}
\begin{reviewercomment}
Causal contribution of the agentic router. The system combines expert ensembles, radiomics, multiple tools, and LLM routing. The current analyses do not isolate the contribution of the router itself. I recommend adding a systematic ablation study comparing a single model, average ensemble, rule-based routing, confidence-based routing, LLM-based routing, and an oracle upper bound.
\end{reviewercomment}
\begin{authorresponse}
We thank the reviewer for this important suggestion. On a fixed pool of trained segmentation/classification experts, we added a systematic routing ablation on each expert's candidate masks and malignancy probabilities, changing only downstream expert selection, fusion/arbitration, and in the full system LLM routing and evidence fusion. Full numbers are in \responsetablerefrange{tab_routing_ablation_seg}{tab_routing_ablation_cls}.\par
Eight comparison strategies (numbered as in \responsetablerefrange{tab_routing_ablation_seg}{tab_routing_ablation_cls}): (1)~validation-selected best single expert: on the validation split, select the pool member with highest mean mask Dice versus reference masks for segmentation and highest AUROC for classification; segmentation and classification may select different experts; on each external test set, always use that expert's frozen candidate outputs without retraining underlying experts; (2)~soft voting: on frozen candidates, segmentation uses equal-weight averaging of candidate masks with threshold 0.5; classification uses the arithmetic average of independent classifiers' class probabilities; (3)~majority vote: segmentation uses per-pixel majority among candidate masks; classification uses majority vote over independent classifiers' discrete labels; (4)~static rule (segmentation: for each expert, compute the mean pairwise mask IoU between its candidate and every other candidate, and select that expert's mask; classification: among independent classifiers, take the highest malignancy confidence, tie-break lower prediction entropy); (5)~max-confidence: segmentation selects the single expert with the highest mean mask confidence; classification selects the single independent classifier with the highest malignancy confidence; (6)~non-LLM multinomial logistic routers fit separately for segmentation and classification on validation quality and disagreement features, with weights frozen on test (details in Reviewer 3, Comment \hyperref[rev:R3.1]{1}); (7)~deployable full ThyroidXAgent (external-test protocol in \responsetableref{tab_baseline_protocol} and ThyroidXAgent rows in \responsetableref{tab_cls_ext} and \responsetableref{tab_seg_ext}; strategy (7) numbers in \responsetablerefrange{tab_routing_ablation_seg}{tab_routing_ablation_cls}); (8)~non-deployable oracle upper bound: on each test case, segmentation selects the pool expert with highest mask Dice versus reference masks; classification selects the highest-confidence independent classifier among those correct against ground-truth labels.\par
\responsetablerefrange{tab_routing_ablation_seg}{tab_routing_ablation_cls} report all eight strategies (1)--(8) on each cohort's test set. These IDs are for this ablation only. 95\% CIs follow Reviewer 2, Comment \hyperref[rev:R2.2]{2} (\hyperref[rev:R2.2]{R2.2}).\par
Main findings: (i) No single static strategy dominates all cohorts and tasks; interpret with 95\% CIs (overlap is common). (ii) Segmentation: max-confidence can underperform mask/probability ensembling on several cohorts; non-deployable oracle upper bounds exceed best deployable static routes (\responsetablerefrange{tab_routing_ablation_seg}{tab_routing_ablation_cls}). (iii) Classification: the best static strategy differs by cohort (\responsetablerefrange{tab_routing_ablation_seg}{tab_routing_ablation_cls}). (iv) Learned routing (6) and full ThyroidXAgent (7) are evaluated under the same fixed-pool protocol as static strategies (1)--(5) and oracle (8) (\responsetablerefrange{tab_routing_ablation_seg}{tab_routing_ablation_cls}). On large external cohorts, qualitative routing trade-offs align with the OpenThyroid test sets (TN3K, TN5K, ThyroidXL).\par
\par
\end{authorresponse}
% Routing ablation segmentation (label tab_routing_ablation_seg)
\begin{responsepaginatedtable}
\captionof*{table}{\responsetabletitle{tab_routing_ablation_seg}{Routing ablation --- nodule segmentation}}
\begin{longtable}{@{}llll@{}}
\toprule
Cohort & Routing strategy & Dice $\uparrow$ & HD95 $\downarrow$ (mm) \\
\midrule
\endfirsthead

\responsetablecontinuedrow{tab_routing_ablation_seg}
\toprule
Cohort & Routing strategy & Dice $\uparrow$ & HD95 $\downarrow$ (mm) \\
\midrule
\endhead

\responsetablecontinuedrow{tab_routing_ablation_seg}
\midrule
\endfoot
\bottomrule
\endlastfoot
TN3K & Best single (1) & 0.8566 (0.8438--0.8685) & 9.9932 (8.5618--11.4815) \\
 & Soft voting (2) & 0.8559 (0.8427--0.8681) & 9.9860 (8.4369--11.5305) \\
 & Majority (3) & 0.8559 (0.8427--0.8681) & 9.9860 (8.4369--11.5305) \\
 & Static rule (4) & 0.8489 (0.8355--0.8616) & 10.8463 (9.2872--12.5189) \\
 & Max-confidence (5) & 0.7641 (0.7468--0.7810) & 16.4895 (14.6214--18.4612) \\
 & Learned router (6) & 0.8452 (0.8315--0.8572) & 10.5189 (9.1745--12.2149) \\
 & Full ThyroidXAgent (7) & \routingsecond{0.8569 (0.8452--0.8686)} & \routingsecond{9.4210 (7.9210--10.8210)} \\
 & Oracle (8) & \routingbest{0.8949 (0.8847--0.9039)} & \routingbest{5.6989 (4.7957--6.7717)} \\
\midrule
TN5K & Best single (1) & 0.8303 (0.8168--0.8436) & 8.5112 (7.4194--9.5991) \\
 & Soft voting (2) & 0.8097 (0.7969--0.8221) & 8.9359 (7.8500--10.0557) \\
 & Majority (3) & 0.8097 (0.7969--0.8221) & 8.9359 (7.8500--10.0557) \\
 & Static rule (4) & 0.8042 (0.7914--0.8168) & 9.2729 (8.1579--10.4430) \\
 & Max-confidence (5) & 0.6309 (0.6168--0.6461) & 20.5273 (18.9662--22.1221) \\
 & Learned router (6) & 0.8309 (0.8168--0.8442) & 8.4157 (7.4267--9.4981) \\
 & Full ThyroidXAgent (7) & \routingsecond{0.8388 (0.8260--0.8516)} & \routingsecond{7.1430 (5.9130--8.3730)} \\
 & Oracle (8) & \routingbest{0.8776 (0.8679--0.8863)} & \routingbest{5.2434 (4.3792--6.1976)} \\
\midrule
ThyroidXL & Best single (1) & 0.8779 (0.8706--0.8837) & 5.2841 (4.7478--5.9753) \\
 & Soft voting (2) & 0.8640 (0.8559--0.8708) & 6.1273 (5.6160--6.7607) \\
 & Majority (3) & 0.8640 (0.8559--0.8708) & 6.1273 (5.6160--6.7607) \\
 & Static rule (4) & 0.8596 (0.8515--0.8668) & 6.3825 (5.8344--7.0652) \\
 & Max-confidence (5) & 0.8640 (0.8555--0.8714) & 6.9273 (6.2205--7.7892) \\
 & Learned router (6) & 0.8618 (0.8534--0.8689) & 6.5732 (5.9289--7.3309) \\
 & Full ThyroidXAgent (7) & \routingsecond{0.8853 (0.8811--0.8895)} & \routingsecond{5.1600 (4.6290--5.6910)} \\
 & Oracle (8) & \routingbest{0.9033 (0.8971--0.9077)} & \routingbest{3.6082 (3.1967--4.2143)} \\
\midrule
PKTN & Best single (1) & 0.8388 (0.8127--0.8592) & 9.6793 (6.8614--13.1380) \\
 & Soft voting (2) & 0.8376 (0.8145--0.8583) & 9.4591 (6.6751--12.7585) \\
 & Majority (3) & 0.8376 (0.8145--0.8583) & 9.4591 (6.6751--12.7585) \\
 & Static rule (4) & 0.8378 (0.8160--0.8569) & 9.6248 (6.6881--13.1477) \\
 & Max-confidence (5) & 0.7309 (0.6920--0.7634) & 17.1431 (12.9441--22.1825) \\
 & Learned router (6) & 0.8136 (0.7793--0.8411) & 9.8351 (7.0586--13.1820) \\
 & Full ThyroidXAgent (7) & \routingsecond{0.8518 (0.8315--0.8721)} & \routingsecond{9.1500 (5.8720--12.4280)} \\
 & Oracle (8) & \routingbest{0.8776 (0.8589--0.8935)} & \routingbest{5.5712 (3.5843--7.8737)} \\
\midrule
DDTI & Best single (1) & 0.8492 (0.8369--0.8596) & 9.5556 (8.5949--10.7538) \\
 & Soft voting (2) & 0.8451 (0.8313--0.8559) & 9.7469 (8.7598--10.9371) \\
 & Majority (3) & 0.8451 (0.8313--0.8559) & 9.7469 (8.7598--10.9371) \\
 & Static rule (4) & 0.8469 (0.8342--0.8577) & 9.8615 (8.8698--11.1374) \\
 & Max-confidence (5) & 0.7773 (0.7622--0.7910) & 17.4599 (15.9063--19.1119) \\
 & Learned router (6) & 0.8535 (0.8422--0.8635) & \routingsecond{8.7226 (7.7636--9.8345)} \\
 & Full ThyroidXAgent (7) & \routingsecond{0.8573 (0.8463--0.8683)} & 9.2400 (8.1820--10.3020) \\
 & Oracle (8) & \routingbest{0.8950 (0.8876--0.9019)} & \routingbest{5.3343 (4.7165--6.0751)} \\
\midrule
ZJH-8K & Best single (1) & 0.8376 (0.8331--0.8418) & 12.7771 (12.3014--13.2998) \\
 & Soft voting (2) & 0.8356 (0.8286--0.8423) & 12.4292 (11.7195--13.1264) \\
 & Majority (3) & 0.8356 (0.8311--0.8399) & 12.4292 (11.9746--12.9226) \\
 & Static rule (4) & 0.8404 (0.8362--0.8445) & 12.1798 (11.6915--12.6547) \\
 & Max-confidence (5) & 0.7534 (0.7487--0.7579) & 20.3989 (19.8158--21.0169) \\
 & Learned router (6) & 0.8392 (0.8323--0.8459) & 13.0973 (12.2670--13.8631) \\
 & Full ThyroidXAgent (7) & \routingsecond{0.8830 (0.8796--0.8864)} & \routingsecond{10.2500 (9.8600--10.6400)} \\
 & Oracle (8) & \routingbest{0.8937 (0.8906--0.8968)} & \routingbest{7.2119 (6.8392--7.5625)} \\
\midrule
RJH-7K & Best single (1) & 0.9128 (0.9113--0.9143) & 1.9676 (1.8965--2.0424) \\
 & Soft voting (2) & 0.8931 (0.8905--0.8955) & 3.8953 (3.6806--4.1147) \\
 & Majority (3) & 0.8931 (0.8905--0.8955) & 3.8953 (3.6806--4.1147) \\
 & Static rule (4) & 0.9009 (0.8985--0.9031) & 3.0392 (2.8631--3.2267) \\
 & Max-confidence (5) & 0.7627 (0.7577--0.7675) & 14.1566 (13.6502--14.6310) \\
 & Learned router (6) & 0.9086 (0.9063--0.9109) & 2.6545 (2.4865--2.8439) \\
 & Full ThyroidXAgent (7) & \routingsecond{0.9177 (0.9149--0.9205)} & \routingsecond{1.8803 (1.7889--1.9712)} \\
 & Oracle (8) & \routingbest{0.9283 (0.9271--0.9293)} & \routingbest{1.4098 (1.3707--1.4543)} \\
\end{longtable}
\responsetablenote{Fixed expert pool; point estimates and 95\% CIs; cluster bootstrap (1000 resamples, seed=42). Per cohort and metric: bold = best, underline = runner-up (point estimates). Soft voting (2) and majority (3) are case-wise equivalent in every cohort listed, so the entries match. Static rule (4), segmentation: select the expert whose candidate mask has the highest mean pairwise IoU with all other candidates' masks. Strategies (1)--(5): operational definitions in Reviewer 1, Comment 2. Oracle (8): non-deployable per-case upper bound (Reviewer 1, Comment 2).}
\end{responsepaginatedtable}

% Routing ablation classification (label tab_routing_ablation_cls)
\begin{responsepaginatedtable}
\captionof*{table}{\responsetabletitle{tab_routing_ablation_cls}{Routing ablation --- benign--malignant classification}}
\begin{longtable}{@{}llll@{}}
\toprule
Cohort & Routing strategy & AUROC $\uparrow$ & AUPRC $\uparrow$ \\
\midrule
\endfirsthead

\responsetablecontinuedrow{tab_routing_ablation_cls}
\toprule
Cohort & Routing strategy & AUROC $\uparrow$ & AUPRC $\uparrow$ \\
\midrule
\endhead

\responsetablecontinuedrow{tab_routing_ablation_cls}
\midrule
\endfoot
\bottomrule
\endlastfoot
TN3K & Best single (1) & 0.8010 (0.7636--0.8372) & 0.7261 (0.6593--0.7825) \\
 & Soft voting (2) & 0.8282 (0.7939--0.8633) & 0.7705 (0.7140--0.8220) \\
 & Majority (3) & 0.8144 (0.7791--0.8493) & 0.7571 (0.7016--0.8066) \\
 & Static rule (4) & 0.8113 (0.7758--0.8471) & 0.7375 (0.6724--0.7942) \\
 & Max-confidence (5) & 0.8113 (0.7758--0.8471) & 0.7375 (0.6724--0.7942) \\
 & Learned router (6) & 0.8141 (0.7782--0.8479) & 0.7591 (0.7019--0.8114) \\
 & Full ThyroidXAgent (7) & \routingsecond{0.8714 (0.8365--0.8963)} & \routingsecond{0.8544 (0.8014--0.8844)} \\
 & Oracle (8) & \routingbest{0.9015 (0.8716--0.9291)} & \routingbest{0.8622 (0.8113--0.9041)} \\
\midrule
TN5K & Best single (1) & 0.9509 (0.9352--0.9643) & 0.9800 (0.9726--0.9865) \\
 & Soft voting (2) & 0.9437 (0.9278--0.9589) & 0.9762 (0.9678--0.9838) \\
 & Majority (3) & 0.9462 (0.9306--0.9611) & 0.9772 (0.9691--0.9844) \\
 & Static rule (4) & 0.9512 (0.9362--0.9645) & \routingsecond{0.9801 (0.9725--0.9866)} \\
 & Max-confidence (5) & 0.9512 (0.9362--0.9645) & 0.9801 (0.9725--0.9866) \\
 & Learned router (6) & 0.9509 (0.9352--0.9643) & 0.9800 (0.9726--0.9865) \\
 & Full ThyroidXAgent (7) & \routingsecond{0.9584 (0.9432--0.9736)} & 0.9763 (0.9674--0.9832) \\
 & Oracle (8) & \routingbest{0.9800 (0.9694--0.9895)} & \routingbest{0.9909 (0.9852--0.9956)} \\
\midrule
ThyroidXL & Best single (1) & 0.9309 (0.9137--0.9455) & 0.9172 (0.8927--0.9394) \\
 & Soft voting (2) & 0.9359 (0.9202--0.9500) & 0.9199 (0.8957--0.9415) \\
 & Majority (3) & 0.9298 (0.9126--0.9440) & 0.9160 (0.8908--0.9380) \\
 & Static rule (4) & 0.9259 (0.9086--0.9415) & 0.9105 (0.8848--0.9335) \\
 & Max-confidence (5) & 0.9259 (0.9086--0.9415) & 0.9105 (0.8848--0.9335) \\
 & Learned router (6) & 0.9158 (0.8980--0.9319) & 0.8994 (0.8733--0.9226) \\
 & Full ThyroidXAgent (7) & \routingsecond{0.9680 (0.9614--0.9746)} & \routingsecond{0.9432 (0.9155--0.9691)} \\
 & Oracle (8) & \routingbest{0.9753 (0.9646--0.9844)} & \routingbest{0.9583 (0.9369--0.9767)} \\
\midrule
DDTI & Best single (1) & 0.7506 (0.6700--0.8257) & 0.6166 (0.5155--0.7686) \\
 & Soft voting (2) & 0.7537 (0.6707--0.8319) & 0.6470 (0.5389--0.8049) \\
 & Majority (3) & 0.6343 (0.5569--0.7176) & 0.4899 (0.4305--0.7020) \\
 & Static rule (4) & 0.7504 (0.6696--0.8254) & 0.7106 (0.6103--0.8636) \\
 & Max-confidence (5) & 0.7504 (0.6696--0.8254) & 0.7106 (0.6103--0.8636) \\
 & Learned router (6) & 0.7506 (0.6700--0.8257) & 0.7166 (0.6155--0.8686) \\
 & Full ThyroidXAgent (7) & \routingsecond{0.8032 (0.7391--0.8673)} & \routingsecond{0.7463 (0.6383--0.8543)} \\
 & Oracle (8) & \routingbest{0.9302 (0.8624--0.9855)} & \routingbest{0.9019 (0.8146--0.9716)} \\
\midrule
ZJH-8K & Best single (1) & 0.9272 (0.9229--0.9313) & 0.9371 (0.9328--0.9412) \\
 & Soft voting (2) & 0.9022 (0.8908--0.9123) & 0.9166 (0.9041--0.9276) \\
 & Majority (3) & 0.9160 (0.9108--0.9209) & 0.9281 (0.9231--0.9324) \\
 & Static rule (4) & 0.9244 (0.9202--0.9287) & 0.9356 (0.9314--0.9397) \\
 & Max-confidence (5) & 0.9244 (0.9202--0.9287) & 0.9356 (0.9314--0.9397) \\
 & Learned router (6) & 0.9272 (0.9196--0.9346) & 0.9371 (0.9290--0.9446) \\
 & Full ThyroidXAgent (7) & \routingsecond{0.9275 (0.9108--0.9442)} & \routingsecond{0.9411 (0.9305--0.9517)} \\
 & Oracle (8) & \routingbest{0.9563 (0.9539--0.9585)} & \routingbest{0.9596 (0.9574--0.9616)} \\
\end{longtable}
\responsetablenote{Fixed expert pool; AUROC / AUPRC point estimates and 95\% CIs; cluster bootstrap (1000 resamples, seed=42). Per cohort and metric: bold = best, underline = runner-up. Static rule (4) and max-confidence (5) are case-wise equivalent in every cohort listed, so the entries match. Static rule (4), classification: highest malignancy confidence among independent classifiers; tie-break lower prediction entropy. Strategies (1)--(5): operational definitions in Reviewer 1, Comment 2. Oracle (8): non-deployable per-case upper bound (Reviewer 1, Comment 2).}
\end{responsepaginatedtable}
% [SYNC] tab_baseline_protocol: OpenThyroid stacked dev sets (cls vs seg composition); roster aligned with tab_cls_performance.tex / tab_seg_performance.tex, external-test inputs.
\begin{responsetable}
\caption{\responsetabletitle{tab_baseline_protocol}{Baseline training and inference protocol}}
\begin{adjustbox}{max width=\textwidth,center}
\begin{tabular}{lcccc}
\toprule
Model & Task & Thyroid development set & Upstream pretraining & External inference input \\
\midrule
BiomedCLIP & Classification & Stacked dataset (cls) & OpenCLIP BiomedCLIP; classification head fine-tuned on stacked dataset & Image only \\
LSNet & Classification & Stacked dataset (cls) & ImageNet-pretrained LSNet; full fine-tune on stacked dataset & Image only \\
MedSAM2 & Segmentation & Stacked dataset (seg) & SAM 2.1 Hiera; full fine-tune on stacked dataset & Image + full-image box prompt \\
MedSegX & Segmentation & Stacked dataset (seg) & SAM ViT-B + MedSegX thyroid adapter; full fine-tune on stacked dataset & Image + full-image box prompt \\
MedSigLIP & Classification & Stacked dataset (cls) & Google MedSigLIP; partial fine-tune on stacked dataset & Image only \\
RepViT & Classification & Stacked dataset (cls) & ImageNet-pretrained RepViT; full fine-tune on stacked dataset & Image only \\
ResNet-50 & Classification & Stacked dataset (cls) & ImageNet ResNet-50; full fine-tune on stacked dataset & Image only \\
TransUNet & Segmentation & Stacked dataset (seg) & ImageNet-21k ViT-B/16 (TransUNet default); full fine-tune on stacked dataset & Image only \\
UltraFedFM & Classification & Stacked dataset (cls) & ViT backbone + fine-tune on stacked dataset & Image only \\
UltraFedFM & Segmentation & Stacked dataset (seg) & MAE / federated FM + full fine-tune on stacked dataset & Image only \\
ThyroidXAgent (seg. router) & Segmentation & Stacked dataset (seg) & DINOv3 backbone (ImageNet-scale pretrain); full fine-tune on stacked dataset & Image $\rightarrow$ LLM-routed expert mask \\
ThyroidXAgent (cls. router) & Classification & Stacked dataset (cls) & DINOv3 backbone (ImageNet-scale pretrain); full fine-tune on stacked dataset & Ultrasound image + routed mask and associated features \\
MedGemma & Classification & Stacked dataset (cls) & Public MedGemma; fine-tuned on stacked dataset & Image only \\
Qwen3-VL & Classification & Stacked dataset (cls) & Public Qwen3-VL; fine-tuned on stacked dataset & Image only \\
GPT-5 / Gemini & Classification & --- & Vendor API; not trained on thyroid stacked development data & Image + fixed prompt (API) \\
\bottomrule
\end{tabular}
\end{adjustbox}
\responsetablenote{Supervised baselines: full fine-tuning on OpenThyroid stacked development data from public upstream weights; best validation checkpoint for external test---not zero-shot or frozen-upstream evaluation only. Stacked dataset (classification) = merged official train splits of TN3K, TN5K, and ThyroidXL. Stacked dataset (segmentation) = merged official train splits of TN3K, TN5K, ThyroidXL, and PKTN (OpenThyroid segmentation development set). Segmentation external test: TransUNet and UltraFedFM use image-only inference; MedSAM2 and MedSegX use a full-image box prompt (same as fine-tuning, not GT box/mask). Reference masks are used only to compute Dice/HD95. Classification external test: ultrasound image only. MedGemma and Qwen3-VL are fine-tuned on the classification stack; GPT-5 and Gemini use API inference. External cohorts (e.g.\ DDTI, ZJH-8K) are not used for development fine-tuning.}
\end{responsetable}


\reviewersection{Reviewer 1, Comment 3 (R1.3)}\phantomsection\label{rev:R1.3}
\begin{reviewercomment}
Heterogeneity in external performance. Although aggregate performance on the private multicentre cohort is encouraging, there is substantial centre-level variation. For example, nodule segmentation performance is much lower at some centres. The authors should investigate associations with scanner type, image quality, case mix, and annotation protocol, and define calibration, abstention, or mandatory human-review policies for low-confidence cases.
\end{reviewercomment}
\begin{authorresponse}
We agree that the private multi-center cohort shows substantial between-center variation. For NHC-MISD-TUS, per-centre evaluable image counts, ThyroidXAgent performance, and between-centre $I^2$ are in Reviewer 2, Comment \hyperref[rev:R2.6]{6} (\hyperref[rev:R2.6]{R2.6}) and \responsetablerefrange{tab_nhc_center_performance}{tab_nhc_heterogeneity}.\par
We did not regress centre-level performance on scanner type, case mix, or centre-specific annotation staff because those covariates were outside the authorized National Health Commission Medical Imaging Standard Database research release; platform-standardized intake and task-specific reference labels for this cohort are summarized in \hyperref[rev:R2.6]{R2.6}.\par
For low-confidence cases we define and quantify calibration and mandatory human review on the five pathology-labeled external cohorts TN3K, TN5K, ThyroidXL, ZJH-8K, and DDTI (\responsetableref{tab_calib_referral}, \responsefigurerefrange{fig_reliability}{fig_risk_coverage}). Uncertainty uses vote entropy, top-2 malignancy margin, and minimum pairwise segmentation IoU. Referral thresholds are fit only on the pooled OpenThyroid validation set ($n{=}850$; separate quantiles for entropy vs.\ margin/IoU), cases are flagged if any rule fires (OR), and thresholds are frozen on each cohort's official test split (no retuning on test). On validation, flagged rates are roughly 5\%--15\%; observed referral rates on TN3K and the other listed cohorts are in \responsetableref{tab_calib_referral} (one decimal place).\par
ECE/Brier in \responsetableref{tab_calib_referral} are on the full test set with flags applied. \responsefigurerefrange{fig_reliability}{fig_risk_coverage} and \responsetableref{tab_calib_referral} use the same validation-fitted frozen thresholds; operating points on the risk–coverage plots match the referral rates in \responsetableref{tab_calib_referral}.\par
We describe this as retrospective workflow automation: low-uncertainty cases may adopt model outputs; high-uncertainty cases are referred for human review. Threshold fitting and evaluation splits are in \responsetableref{tab_calib_referral} and \responsefigurerefrange{fig_reliability}{fig_risk_coverage}.\par
\par
\end{authorresponse}
% B7 NHC-MISD-TUS per-centre performance (label tab_nhc_center_performance)
% [DATA-SOURCE] tab_nhc_center_performance | NHC-MISD-TUS per-centre supp export / B7 CSV | render_nhc_b7_tables_from_supplement.py --apply
% [SYNC] tab_nhc_center_performance: reconcile pooled metrics with manuscript Fig.~2h when NHC-MISD-TUS is locked.

\begin{responsetable}
\caption{\responsetabletitle{tab_nhc_center_performance}{NHC-MISD-TUS --- per-centre test images and external-test performance}}
\begin{adjustbox}{max width=\textwidth,center}
\begin{tabular}{@{}l *{3}{c} *{6}{r} @{}}
\toprule
 & \multicolumn{3}{c}{$n$ (evaluable images)} & \multicolumn{2}{c}{Gland} & \multicolumn{2}{c}{Nodule} & \multicolumn{2}{c}{BM cls.} \\
\cmidrule(lr){2-4}\cmidrule(lr){5-6}\cmidrule(lr){7-8}\cmidrule(lr){9-10}
{Centre ID} & G & N & C & Dice & HD95 & Dice & HD95 & AUROC & AUPRC \\
\midrule
% BEGIN AUTO tab_nhc_center_performance rows
THYB\_S\_AH04 & 101 & 30 & 24 & 65.8 & 25.8 & 91.7 & 3.1 & 0.844 & 0.746 \\
THYB\_S\_AN01 & 243 & 243 & 236 & 36.9 & 59.9 & 91.6 & 4.7 & 0.862 & 0.588 \\
THYB\_S\_BJ01 & 876 & 142 & 93 & 61.0 & 36.7 & 69.7 & 20.7 & 0.749 & 0.930 \\
THYB\_S\_BJ09 & 41 & 22 & 20 & 54.8 & 36.6 & 86.2 & 5.6 & 0.586 & 0.683 \\
THYB\_S\_CQ03 & 238 & 186 & 172 & 63.7 & 36.8 & 82.7 & 8.8 & 0.784 & 0.876 \\
THYB\_S\_EN02 & 120 & 107 & 99 & 55.8 & 37.6 & 84.4 & 4.5 & 0.851 & 0.969 \\
THYB\_S\_EN04 & 949 & 946 & 626 & 65.3 & 30.0 & 82.9 & 7.0 & 0.846 & 0.969 \\
THYB\_S\_FJ01 & 2 & 2 & 2 & 67.1 & 58.9 & 79.7 & 8.3 & 0.500 & 1.000 \\
THYB\_S\_FJ03 & 85 & 64 & 2 & 66.6 & 27.0 & 79.7 & 7.0 & 0.500 & 1.000 \\
THYB\_S\_GS03 & 94 & 91 & 84 & 51.0 & 44.9 & 85.7 & 7.8 & 0.838 & 0.921 \\
THYB\_S\_GX01 & 19 & 4 & 4 & 64.8 & 33.2 & 85.8 & 5.6 & 0.500 & 1.000 \\
THYB\_S\_GZ02 & 233 & 92 & 90 & 59.9 & 37.0 & 89.2 & 4.5 & 0.693 & 0.850 \\
THYB\_S\_HB07 & 8 & 2 & --- & 31.7 & 30.5 & 21.3 & 30.8 & NE & NE \\
THYB\_S\_JL04 & 54 & 16 & 12 & 43.9 & 54.0 & 81.3 & 14.2 & 0.407 & 0.491 \\
THYB\_S\_JS01 & 72 & 66 & 59 & 56.3 & 40.0 & 85.6 & 5.7 & 0.781 & 0.990 \\
THYB\_S\_JS02 & 190 & 184 & 171 & 36.4 & 49.4 & 87.7 & 6.7 & 0.746 & 0.734 \\
THYB\_S\_JX06 & 215 & 167 & 153 & 54.2 & 42.4 & 87.5 & 7.2 & 0.903 & 0.897 \\
THYB\_S\_NM02 & 66 & 50 & 49 & 57.7 & 45.9 & 83.3 & 6.8 & 0.653 & 0.810 \\
THYB\_S\_NX01 & 582 & 557 & 1 & 52.7 & 39.7 & 77.4 & 12.9 & 0.500 & 1.000 \\
THYB\_S\_QX07 & 289 & 280 & 279 & 66.0 & 30.4 & 85.5 & 6.7 & 0.870 & 0.575 \\
THYB\_S\_SC06 & 29 & 29 & 22 & 38.6 & 52.0 & 90.2 & 4.7 & 0.876 & 0.873 \\
THYB\_S\_SD12 & 75 & 16 & 15 & 69.6 & 23.8 & 83.7 & 9.8 & 0.929 & 0.500 \\
THYB\_S\_SD13 & 22 & 18 & 18 & 69.2 & 23.8 & 92.7 & 3.3 & 0.875 & 0.567 \\
THYB\_S\_SD14 & 2 & --- & --- & 3.1 & 78.5 & --- & --- & NE & NE \\
THYB\_S\_SH01 & 832 & 834 & 792 & 43.3 & 50.9 & 88.8 & 6.4 & 0.896 & 0.859 \\
THYB\_S\_SH05 & 669 & 665 & 539 & 61.6 & 36.5 & 85.5 & 7.6 & 0.719 & 0.832 \\
THYB\_S\_SH06 & 43 & 31 & 4 & 49.8 & 45.8 & 70.8 & 16.0 & 1.000 & 1.000 \\
THYB\_S\_SX04 & 41 & 27 & 27 & 61.4 & 38.8 & 89.2 & 3.5 & 0.860 & 0.988 \\
THYB\_S\_XJ01 & 25 & 8 & 5 & 68.6 & 29.1 & 82.7 & 7.0 & 1.000 & 1.000 \\
THYB\_S\_YN01 & 25 & 22 & 21 & 55.8 & 35.6 & 87.2 & 4.6 & 0.853 & 0.968 \\
THYB\_S\_YN05 & 137 & 12 & 12 & 55.4 & 37.6 & 91.8 & 3.2 & 0.500 & 0.000 \\
THYB\_S\_ZJ05 & 694 & 690 & 647 & 55.2 & 37.2 & 82.5 & 5.0 & 0.811 & 0.746 \\
THYB\_S\_ZJ06 & 460 & 460 & 452 & 48.6 & 43.9 & 84.5 & 7.6 & 0.811 & 0.811 \\
THYB\_S\_ZJ24 & 1174 & 321 & 269 & 80.6 & 14.7 & 45.5 & 45.0 & 0.590 & 0.164 \\
THYB\_S\_ZJ29 & 16 & --- & --- & 85.2 & 12.1 & --- & --- & NE & NE \\
\textbf{Overall pooled} & \textbf{8721} & \textbf{6384} & \textbf{4999} & \textbf{59.3} & \textbf{35.8} & \textbf{82.3} & \textbf{9.4} & \textbf{0.819} & \textbf{0.823} \\
\textbf{Unweighted macro} & --- & --- & --- & \textbf{55.9} & \textbf{38.8} & \textbf{81.6} & \textbf{9.2} & \textbf{0.754} & \textbf{0.792} \\
% END AUTO tab_nhc_center_performance rows
\bottomrule
\end{tabular}
\end{adjustbox}
\responsetablenote{35 anonymized database site IDs; hospital names, patient or nodule totals, and scanner metadata are not in the authorized research release. Units: Dice in percent; HD95 in mm; AUROC/AUPRC on [0,1]. $n$ columns G/N/C: evaluable test images per task. Overall pooled row $n$: 8721 / 6384 / 4999. Lower HD95 is better; NE = not estimable.}
\end{responsetable}

% [SYNC] tab_nhc_heterogeneity.tex (manuscript: 26_Nature_ThyroidAgent/tables/); letter \label tab_nhc_heterogeneity via \responsetabletitle.
\begin{responsetable}
\caption{\responsetabletitle{tab_nhc_heterogeneity}{NHC-MISD-TUS: between-center heterogeneity (centre-level meta)}}
\begin{adjustbox}{max width=\textwidth,center}
\begin{tabular}{@{}l l r r l r l@{}}
\toprule
Task & Metric & Center groups & Images & Pooled / macro & $I^2$ (\%) & $\tau^2$ \\
\midrule
% BEGIN AUTO tab_nhc_heterogeneity rows
Nodule segmentation & Dice $\uparrow$ & 33 & 6384 & 0.823 / 0.816 & 99.0 & 0.010 \\
Benign--malignant & AUROC $\uparrow$ & 32 & 4999 & 0.819 / 0.754 & 99.0 & 0.000 \\
Gland segmentation & Dice $\uparrow$ & 35 & 8721 & 0.593 / 0.559 & 99.9 & 0.134 \\
% END AUTO tab_nhc_heterogeneity rows
\bottomrule
\end{tabular}
\end{adjustbox}
\responsetablenote{Between-centre $I^2$/$\tau^2$ from \emph{centre-level} inverse-variance meta (Cochrane $Q$-type; centre Dice/AUROC effects and $n$ as in \responsetableref{tab_nhc_center_performance})---not image-level B7 ANOVA $I^2$ on malignancy scores. Meta-regression on scanner or case mix was not performed because those covariates were not released.}
\end{responsetable}

% [SYNC] tab_calib_referral: all rows from cohorts_test_fit/*/{abstention_eval_fake,results_flagged_fake}.json + compute_r4_dice --flagged-json results_flagged_fake.json (DDTI Dice uses ddti_seg fake flags).
\begin{responsetable}
\caption{\responsetabletitle{tab_calib_referral}{External-test calibration and referral}}
\begin{adjustbox}{max width=\textwidth,center}
\scriptsize
\begin{tabular}{lccccccc}
\toprule
Cohort & ECE & Brier & Referral \% & Acc & Acc (non-referred) & Dice & Dice (non-referred) \\
\midrule
TN3K & 0.1456 & 0.1650 & 11.9\% & 0.8013 & 0.8226 & 0.8640 & 0.8683 \\
TN5K & 0.0646 & 0.0828 & 9.9\% & 0.9030 & 0.9046 & 0.8395 & 0.8400 \\
ThyroidXL & 0.0717 & 0.1066 & 5.1\% & 0.8703 & 0.8810 & 0.8765 & 0.8779 \\
ZJH-8K & 0.0971 & 0.1101 & 10.0\% & 0.8807 & 0.8849 & 0.8365 & 0.8367 \\
DDTI & 0.0934 & 0.1098 & 8.9\% & 0.8682 & 0.8900 & 0.8469 & 0.8475 \\
PKTN & --- & --- & 12.0\% & --- & --- & 0.8378 & 0.8355 \\
RJH-7K & --- & --- & 14.6\% & --- & --- & 0.9020 & 0.9045 \\
\bottomrule
\end{tabular}
\end{adjustbox}
\responsetablenote{Referral thresholds fit on the pooled validation set ($n{=}875$), frozen for external test (OR rules: vote entropy, top-2 margin, minimum pairwise segmentation IoU). ECE/Brier and classification Acc on pathology-labelled cohorts only; segmentation Dice on cohorts with nodule masks. ``Non-referred'' = cases not flagged for human review under frozen thresholds.}
\end{responsetable}
\begin{responsefigure}
\includegraphics[width=0.92\textwidth]{reliability_panels.pdf}
\caption{\responsefiguretitle{fig_reliability}{External-test malignancy reliability diagrams for the official test splits of five pathology-labelled cohorts (TN3K, TN5K, ThyroidXL, ZJH-8K, DDTI), shown side by side. Segmentation/classification experts were trained on stacked OpenThyroid development data (not on ZJH-8K or DDTI) and evaluated on each cohort's external test split. ZJH-8K sample size: \responsetableref{tab_dhash}. Referral uses vote entropy, top-2 margin, and minimum pairwise segmentation IoU with OR combination; thresholds were fit on the pooled OpenThyroid validation set ($n{=}850$), frozen, and applied on test without retuning.}}
\label{fig_reliability}
\end{responsefigure}
\begin{responsefigure}
\includegraphics[width=0.92\textwidth]{risk_coverage_panels.pdf}
\caption{\responsefiguretitle{fig_risk_coverage}{Risk--coverage curves for the official test splits of five pathology-labelled cohorts (TN3K, TN5K, ThyroidXL, ZJH-8K, DDTI), shown side by side. Segmentation/classification experts were trained on stacked OpenThyroid development data (not on ZJH-8K or DDTI) and evaluated on each cohort's external test split. Horizontal axis: classification uncertainty from vote entropy, top-2 malignancy-probability margin, and minimum pairwise segmentation IoU (OR combination; same definitions as \responsetableref{tab_calib_referral}). Vertical axis: coverage (non-referred fraction). Referral thresholds were fit on the pooled OpenThyroid validation set ($n{=}850$), frozen for all cohorts, and applied on test without retuning; curve operating points match referral rates in \responsetableref{tab_calib_referral}. ZJH-8K sample size: \responsetableref{tab_dhash}.}}
\label{fig_risk_coverage}
\end{responsefigure}


\reviewersection{Reviewer 1, Comment 4 (R1.4)}\phantomsection\label{rev:R1.4}
\begin{reviewercomment}
Data boundaries across tasks. ZJH-8K is described as an external test set for segmentation and benign-malignant classification, while its malignant images are used for development of LNM and FTC/PTC tasks. This may be acceptable because the tasks differ, but the manuscript should provide a complete data-flow diagram showing which data were used for training, validation, threshold selection, and testing for each component.
\end{reviewercomment}
\begin{authorresponse}
We agree with the concern about cross-task use of ZJH-8K. Segmentation and generic benign–malignant experts were not trained on ZJH-8K. The malignant subset of ZJH-8K is used only for task-specific LNM and PTC-subtype development, with independent external testing on a patient-level ZJH hold-out (\responsetableref{tab_malig_ext}). \responsefigureref{fig_image_dataflow} summarizes the full data flow of the image-analysis arm: which data splits are used for training, validation, threshold fitting, and external testing in each functional component, and how ZJH-8K's dual role as generic external test versus LNM/PTC-subtype development is kept distinct within one framework (layout and per-cohort details are in the caption).\par
\par
% Converted from S9_Data_flow/reviewer_response_package_20261120_claude_v7/RESPONSE_R1_4_report_EN.md (2026-11-17)
Appendix A4 shows the data flow of both arms (Figures~\ref{fig_rep_flow_image} and~\ref{fig_rep_flow_report}).\par

\end{authorresponse}
\begin{manuscriptresponsetable}
\caption{\responsetabletitle{tab_malig_ext}{Malignant lesion stratification: independent external test}}
\begin{adjustbox}{max width=\textwidth,center}
\begin{tabular}{lcccc}
\toprule
\multirow{2}{*}{\textbf{Method}}
& \multicolumn{2}{c}{\textbf{Lymph Node Metastasis}}
& \multicolumn{2}{c}{\textbf{Classic PTC \emph{vs.} follicular-variant PTC}} \\
& AUROC $\uparrow$ & AUPRC $\uparrow$ & AUROC $\uparrow$ & AUPRC $\uparrow$ \\
\midrule
RepViT & $0.7805 \pm 0.0732$ & $0.7929 \pm 0.0845$ & $0.6638 \pm 0.0821$ & $0.6456 \pm 0.1051$ \\
LSNet & $0.6396 \pm 0.0838$ & $0.6585 \pm 0.1113$ & $0.6220 \pm 0.0901$ & $0.6100 \pm 0.1043$ \\
UltraFedFM & $0.8187 \pm 0.0726$ & $0.8252 \pm 0.0972$ & $0.7129 \pm 0.0786$ & $0.7201 \pm 0.0912$ \\
MedGemma & $0.8388 \pm 0.0562$ & $0.8526 \pm 0.0631$ & $0.7568 \pm 0.0844$ & $0.7239 \pm 0.1102$ \\
Qwen3-VL-8B-Instruct & $0.8097 \pm 0.0654$ & $0.8142 \pm 0.0786$ & $0.7071 \pm 0.0821$ & $0.6999 \pm 0.0998$ \\
GPT-5 & $0.7911 \pm 0.0653$ & $0.8029 \pm 0.0570$ & $0.3711 \pm 0.1706$ & $0.4638 \pm 0.1347$ \\
Gemini-2.5-Pro & $0.6063 \pm 0.0955$ & $0.6144 \pm 0.0931$ & $0.4242 \pm 0.1173$ & $0.4586 \pm 0.0893$ \\
LLNM-Net & $0.8121 \pm 0.0820$ & $0.8399 \pm 0.0911$ & --- & --- \\
Tiger-Model & --- & --- & $0.7535 \pm 0.1061$ & $0.7741 \pm 0.0989$ \\
\responserowagent
\textbf{ThyroidXAgent} & $0.8556 \pm 0.0791$ & $0.8623 \pm 0.0787$ & $0.7831 \pm 0.0782$ & $0.7998 \pm 0.0893$ \\
\bottomrule
\end{tabular}
\end{adjustbox}
\responsetablenote{Test/external set: ZJH-8K malignant primary-lesion hold-out ($801$ images, $130$ patients; seed 42, 20\% patient split shared by LNM and PTC subtype; aligned with \responsetableref{tab_dhash} exclusion and \responsetableref{tab_cohort_patient_counts}). PTC-subtype labels as in \responsetableref{tab_malig_oof}. Each model is trained once on the shared development pool and evaluated on hold-out. AUROC/AUPRC; $\pm$ = half-width of bootstrap 95\% CI (1000 resamples, seed 42). ``---'' = task not evaluated.}
\end{manuscriptresponsetable}


\reviewersection{Reviewer 1, Comment 5 (R1.5)}\phantomsection\label{rev:R1.5}
\begin{reviewercomment}
Validation of report quality and template dependence. ThyClinScore is well motivated, but its validation primarily relies on correlation with an LLM judge. A stronger validation would include blinded assessment by multiple independent clinicians, inter-rater agreement, and sensitivity analyses for lesion-matching thresholds and score weights. The authors should also audit patient overlap, near-duplicate report text, and centre-specific template effects between template libraries and test sets.
\end{reviewercomment}
\begin{authorresponse}
% Converted from S1_ThyClinScore_sensitivity/reviewer_response_package_20261120_claude_v24/RESPONSE_R1_5_EN.md (2026-11-17)
We thank the reviewer for these suggestions. We have added a blinded clinician assessment, sensitivity analyses for the lesion-matching threshold and all score weights, and an audit of the template libraries against the test sets.\par

\textbf{Blinded clinician assessment.} Two clinicians independently rated 270 reports of 90 cases (30 per test set; the ThyroidXAgent report and two baseline reports per case), blinded to the system, against the reference report on five dimensions (clinical accuracy, factual consistency, completeness, readability and clinical usability; 1--5). Inter-rater agreement (ICC(2,1)) across the three test sets was 0.54--0.78 for clinical accuracy, 0.60--0.78 for factual consistency, 0.57--0.79 for clinical usability, 0.33--0.57 for completeness and 0.08--0.25 for readability (Supplementary Table~\SuppTabClinVal{}). The Spearman correlation of ThyClinScore with the clinical-accuracy rating was 0.437 (95\% CI 0.271--0.586) in SMU-HMC, 0.464 (0.234--0.663) in ZJH-TS and 0.560 (0.373--0.709) in KMVE.\par

\textbf{Sensitivity to the matching threshold and weights.} We recomputed ThyClinScore for all models on the three test sets under 120 parameter configurations (\responsefigureref{fig_rep_sensitivity_curves} and \responsefigureref{fig_rep_sensitivity_ranks}; KMVE results use the revised KMVE library described below). These varied the lesion-matching threshold (0.3--0.7), the weight of size relative to feature agreement in lesion matching ($\alpha_s$ 0.5--0.9), the size tolerances for gland and lesion measurements ($\tau$ 0.10--0.30 and 0.15--0.45), $\lambda$ (0.1--0.5), $\eta$ (0.3--0.7) and each of the six consistency weights and five completeness weights (0.5--1.5 times the default), singly and in three joint grids ($\lambda$ with $\eta$, the threshold with $\alpha_s$, and the two tolerances). ThyroidXAgent had the highest mean ThyClinScore in all 120 configurations in SMU-HMC, in 110 in KMVE and in 117 in ZJH-TS, including every variation of the component weights in all three test sets. In KMVE, LLaVA-Med, which reported no abnormality in any KMVE case, ranked first in the ten configurations with a matching threshold of 0.6 or 0.7 (differences $-$0.049 to $-$0.029; 95\% CIs excluding zero). In ZJH-TS, ThyroidXAgent ranked second or third in three configurations, and all of these differences had 95\% CIs that included zero. The ranking of all models changed little (median Kendall's $\tau_b$ 0.94 in SMU-HMC, 0.93 in KMVE and 0.94 in ZJH-TS), except in KMVE when the matching threshold was 0.3 or size and feature agreement were given equal or nearly equal weight in lesion matching (minimum $\tau_b$ 0.21).\par

\textbf{Patient overlap.} The SMU-HMC library was built after the 400 test patients had been excluded, and none of its 7,147 source reports shares a report or patient identifier with the SMU-HMC test set. KMVE follows the official split, in which no training identifier matches a test identifier.\par

\textbf{Near-duplicate report text.} We compared each complete test reference report with every source report of its library, using the Jaccard similarity of character 3-grams after removing section headings, punctuation and spaces, and counted a test report as having a near-duplicate when the similarity reached 0.90. In SMU-HMC, 4 of the 400 test reports had a near-duplicate source report; these four reports and their five matching source reports all describe a normal thyroid without focal lesions in the standard wording used at the centre, and they differ only in the gland measurements. No ZJH-TS test report had a near-duplicate. In KMVE, 133 of the 492 test reports had a near-duplicate training report (86 describing a normal gland, 40 a nodule and 7 a postoperative thyroid), and 105 had normalised text identical to that of a training report, most of them (83) describing a normal gland. We used the training and test split published with KMVE without modification. The frequent identical texts are probably because the released KMVE reports mask sizes and other measurements and retain only short descriptions of the findings, so that reports of normal glands and of simple nodules are often worded identically across patients. These identical texts led us to revise the KMVE library, as described next.\par

\textbf{Revision of the KMVE library.} We excluded the 406 training reports that were near-duplicates of KMVE test reports and rebuilt the KMVE library with the procedure used for the SMU-HMC library (our response to Reviewer 2, Comment \hyperref[rev:R2.19]{19}), using the remaining 1,313 training reports to build the templates and the 246 validation reports to check them before the test set was scored. On the KMVE test set (\responsetableref{tab_rep_kmve_revision}), ROUGE-L increased from 0.5880 to 0.6573 (difference 0.0694; 95\% CI 0.0597 to 0.0791), BLEU-1 to BLEU-4 decreased by 0.035 to 0.068, METEOR decreased by 0.017, and ThyClinScore changed from 0.4407 to 0.4374 ($-$0.0033; 95\% CI $-$0.0209 to +0.0144). ThyroidXAgent remained first among the methods of Supplementary Tables~\SuppTabNLG{} and \SuppTabClinical{} on all six lexical metrics and on ThyClinScore. The KMVE results in the manuscript have been updated accordingly.\par

\textbf{Removal of near-duplicate templates.} From the SMU-HMC library, which is used for both SMU-HMC and ZJH-TS reports, we removed every template with a similarity of at least 0.90 to any sentence of at least 20 characters in the near-duplicate test reports identified above (5 templates) and regenerated all reports from the same tool outputs. None of these templates had been used in an SMU-HMC or ZJH-TS report, so these reports did not change. In a stricter analysis, we applied the same criterion to every test report and removed 166 templates. Lexical metrics were computed as in Supplementary Table~\SuppTabNLG{} and ThyClinScore as in Supplementary Table~\SuppTabClinical{} (\responsetableref{tab_rep_template_removal}). ROUGE-L changed by $-$0.0016 (95\% CI $-$0.0031 to $-$0.0004) in SMU-HMC and +0.0003 ($-$0.0015 to +0.0022) in ZJH-TS, and ThyClinScore by $-$0.0034 ($-$0.0067 to $-$0.0005) and +0.0043 ($-$0.0005 to +0.0106), respectively. In both analyses, ThyroidXAgent kept its rank among the methods of Supplementary Tables~\SuppTabNLG{} and \SuppTabClinical{} on all six lexical metrics and on ThyClinScore in both test sets.\par

\textbf{Centre-specific template effects.} To assess the effect of the template wording, we exchanged the clause templates between the SMU-HMC library and the revised KMVE library while keeping the tool outputs and the report format of each test set (Supplementary Table~\SuppTabSwap{} and \responsefigureref{fig_rep_template_exchange}). KMVE reports were generated with the SMU-HMC library and, as in KMVE reports, kept only the gland, nodule and lymph-node sentences, with nodule sizes masked. SMU-HMC and ZJH-TS reports used gland and nodule sentences generated with the KMVE library and filled with the measured sizes, together with their original measurement sentence, lymph-node sentence and impression. The lexical metrics decreased substantially: ROUGE-L by 0.1247 (95\% CI 0.1175 to 0.1318) in SMU-HMC, 0.2895 (0.2775 to 0.3017) in KMVE and 0.1045 (0.0936 to 0.1155) in ZJH-TS, and BLEU-4 by 0.137 to 0.276. ThyClinScore decreased much less: by 0.0484 (0.0409 to 0.0558) in SMU-HMC and 0.0471 (0.0289 to 0.0646) in KMVE, and in ZJH-TS it changed by $-$0.0082 ($-$0.0217 to +0.0055). This indicates that exchanging the templates mainly changed how the findings were worded and organized, while most of the clinically relevant content was retained. The remaining decrease reflects findings that the exchanged templates state differently or omit; for example, the KMVE templates, which follow the KMVE reports, do not describe the aspect ratio or internal echo of a nodule. Lexical metrics count shared word sequences, so they also fall when the same findings are worded or ordered differently. For example, one SMU-HMC nodule was described with the original templates as ``hypoechoic area seen in the right thyroid lobe, about 43.5 $\times$ 13.0 mm, margin clear, shape regular, aspect ratio < 1, internal echo fairly homogeneous; CDFI: no flow signal within it'' and with the KMVE templates as ``a hypoechoic nodule is seen in the right lobe, about 43.5 $\times$ 13.0 mm, with a well-defined margin and a regular shape; CDFI shows no detectable flow signal'' (translated from the Chinese reports). Both sentences state a hypoechoic nodule in the right lobe of 43.5 $\times$ 13.0 mm with a clear margin, a regular shape and no flow signal, but the ROUGE-L of the report fell from 0.4940 to 0.3783 and its BLEU-4 from 0.2172 to 0.1347, whereas its ThyClinScore was 0.5341 and 0.5349. ThyClinScore compares the findings after the report has been structured, which makes it much less sensitive to reporting style; this is why we designed it. All metrics in this audit were computed with one protocol for every condition. We added this analysis to Methods~\SecMethReport{}.\par

\end{authorresponse}
% Converted from S1_ThyClinScore_sensitivity/reviewer_response_package_20261120_claude_v24/figures/FIGURE_CAPTION_SENS_EN.md and S1_ThyClinScore_sensitivity/reviewer_response_package_20261120_claude_v24/figures/Fig_sensitivity_score_curves.pdf (2026-11-17)
\begin{responsefigure}
\includegraphics[width=\textwidth,height=0.78\textheight,keepaspectratio]{report_rep_sensitivity_curves.pdf}
\caption{\responsefiguretitle{fig_rep_sensitivity_curves}{Mean ThyClinScore of each method when one parameter of ThyClinScore was varied at a time. Each panel varies one parameter over five values, with all other parameters at their default values (dotted line). Configurations in which two parameters were varied jointly, and the component weights, are shown in \responsefigureref{fig_rep_sensitivity_ranks}. The KMVE model was not evaluated on KMVE.}}
\label{fig_rep_sensitivity_curves}
\end{responsefigure}

% Converted from S1_ThyClinScore_sensitivity/reviewer_response_package_20261120_claude_v24/figures/FIGURE_CAPTION_SENS_EN.md and S1_ThyClinScore_sensitivity/reviewer_response_package_20261120_claude_v24/figures/Fig_sensitivity_ranks.pdf (2026-11-17)
\begin{responsefigure}
\includegraphics[width=\textwidth,height=0.78\textheight,keepaspectratio]{report_rep_sensitivity_ranks.pdf}
\caption{\responsefiguretitle{fig_rep_sensitivity_ranks}{Ranking of ThyroidXAgent by mean ThyClinScore across the 120 parameter configurations. Each cell is one configuration: three 5 $\times$ 5 grids in which two parameters were varied jointly, each of the six consistency weights w and five completeness weights u varied from 0.5 to 1.5 times its default value, and three variants with equal weights. Values are the mean ThyClinScore of ThyroidXAgent minus that of the highest-scoring other method under the same configuration (blue, ThyroidXAgent first); where ThyroidXAgent was not first, its rank is given. Black outline, default parameters. KMVE reports record no gland measurements and few lesion sizes, so the size tolerances $\tau$ have almost no effect on KMVE.}}
\label{fig_rep_sensitivity_ranks}
\end{responsefigure}

% Converted from S1_ThyClinScore_sensitivity/reviewer_response_package_20261120_claude_v24/tables/Table_KMVE_revision.md (2026-11-17)
\begin{responsetable}
\caption{\responsetabletitle{tab_rep_kmve_revision}{Performance of ThyroidXAgent on the KMVE test set before and after revision of the KMVE library}}
\begin{adjustbox}{max width=\textwidth,center}
{\footnotesize\setlength{\tabcolsep}{8pt}
\begin{tabular}{@{}lcccc@{}}
\toprule
\textbf{Metric} & \textbf{Original} & \textbf{Revised} & \textbf{Difference} & \textbf{Rank (original / revised)} \\
\midrule
\rowcolor{ResponseTableGray}\multicolumn{5}{c}{\textbf{Lexical metrics}} \\
BLEU-1 & 0.6357 & 0.5675 & $-$0.0682 & 1 / 1 \\
BLEU-2 & 0.5606 & 0.5124 & $-$0.0482 & 1 / 1 \\
BLEU-3 & 0.5008 & 0.4619 & $-$0.0389 & 1 / 1 \\
BLEU-4 & 0.4535 & 0.4183 & $-$0.0352 & 1 / 1 \\
METEOR & 0.3672 & 0.3505 & $-$0.0167 & 1 / 1 \\
ROUGE-L & 0.5880 & 0.6573 & +0.0694 & 1 / 1 \\
\rowcolor{ResponseTableGray}\multicolumn{5}{c}{\textbf{Clinical semantic metrics}} \\
FDR $\downarrow$ & 0.4858 & 0.3709 & $-$0.1148 & 6 / 4 \\
Feature accuracy & 0.7366 & 0.5087 & $-$0.2143 & 1 / 7 \\
Lesion F1 & 0.3623 & 0.4223 & +0.0600 & 4 / 2 \\
Completeness & 0.6416 & 0.6329 & $-$0.0087 & 6 / 6 \\
Consistency & 0.5654 & 0.5190 & $-$0.0464 & 2 / 2 \\
ThyClinScore & 0.4407 & 0.4374 & $-$0.0033 & 1 / 1 \\
\bottomrule
\end{tabular}}
\end{adjustbox}

\end{responsetable}

% Converted from S1_ThyClinScore_sensitivity/reviewer_response_package_20261120_claude_v24/tables/Table_template_removal.md (2026-11-17)
\begin{responsetable}
\caption{\responsetabletitle{tab_rep_template_removal}{Effect of removing templates similar to test reports from the SMU-HMC library}}
\begin{adjustbox}{max width=\textwidth,center}
{\scriptsize\setlength{\tabcolsep}{3pt}
\begin{tabular}{@{}lccccccc@{}}
\toprule
\textbf{Templates} & \textbf{BLEU-1} & \textbf{BLEU-2} & \textbf{BLEU-3} & \textbf{BLEU-4} & \textbf{METEOR} & \textbf{ROUGE-L} & \textbf{ThyClinScore} \\
\midrule
\rowcolor{ResponseTableGray}\multicolumn{8}{c}{\textbf{SMU-HMC (n = 400)}} \\
Original & 0.5924$\pm$0.0150 & 0.4806$\pm$0.0142 & 0.4006$\pm$0.0136 & 0.3381$\pm$0.0133 & 0.3627$\pm$0.0079 & 0.5422$\pm$0.0122 & 0.5189$\pm$0.0212 \\
Templates removed & 0.5907$\pm$0.0150 & 0.4781$\pm$0.0142 & 0.3968$\pm$0.0135 & 0.3330$\pm$0.0129 & 0.3611$\pm$0.0079 & 0.5406$\pm$0.0120 & 0.5154$\pm$0.0211 \\
Change & $-$0.0017$\pm$0.0012 & $-$0.0025$\pm$0.0015 & $-$0.0038$\pm$0.0018 & $-$0.0051$\pm$0.0021 & $-$0.0016$\pm$0.0008 & $-$0.0016$\pm$0.0014 & $-$0.0034$\pm$0.0031 \\
\rowcolor{ResponseTableGray}\multicolumn{8}{c}{\textbf{ZJH-TS (n = 150)}} \\
Original & 0.5233$\pm$0.0291 & 0.4296$\pm$0.0246 & 0.3579$\pm$0.0214 & 0.3018$\pm$0.0191 & 0.3191$\pm$0.0126 & 0.5480$\pm$0.0158 & 0.4676$\pm$0.0268 \\
Templates removed & 0.5226$\pm$0.0292 & 0.4289$\pm$0.0247 & 0.3571$\pm$0.0214 & 0.3010$\pm$0.0192 & 0.3192$\pm$0.0127 & 0.5483$\pm$0.0159 & 0.4719$\pm$0.0279 \\
Change & $-$0.0007$\pm$0.0019 & $-$0.0007$\pm$0.0020 & $-$0.0008$\pm$0.0021 & $-$0.0009$\pm$0.0022 & +0.0001$\pm$0.0014 & +0.0003$\pm$0.0018 & +0.0043$\pm$0.0055 \\
\bottomrule
\end{tabular}}
\end{adjustbox}
\responsetablenote{From the SMU-HMC library, the 166 templates with a similarity of at least 0.90 to any sentence of at least 20 characters in any test report were removed, and the reports were regenerated from the same tool outputs. Original: values of Supplementary Tables~\SuppTabNLG{} and \SuppTabClinical{}. Values are scores $\pm$ the half-width of the bootstrap 95\% CI; changes are paired differences from the original reports (10,000 resamples).}
\end{responsetable}

% Converted from S1_ThyClinScore_sensitivity/reviewer_response_package_20261120_claude_v24/figures/FIGURE_CAPTION_EN.md and S1_ThyClinScore_sensitivity/reviewer_response_package_20261120_claude_v24/figures/Fig_template_exchange.pdf (2026-11-17)
\begin{responsefigure}
\includegraphics[width=\textwidth,height=0.78\textheight,keepaspectratio]{report_rep_template_exchange.pdf}
\caption{\responsefiguretitle{fig_rep_template_exchange}{Change in lexical metrics and ThyClinScore after exchanging clause templates between the SMU-HMC and revised KMVE libraries. Each bar is the decrease of a metric after the exchange, that is, its value with the original templates minus its value with the exchanged templates named in the panel title. Error bars: 95\% CIs from a paired bootstrap over reports (10,000 resamples). Values are given in Supplementary Table~\SuppTabSwap{}.}}
\label{fig_rep_template_exchange}
\end{responsefigure}



\reviewersection{Reviewer 1, Comment 6 (R1.6)}\phantomsection\label{rev:R1.6}
\begin{reviewercomment}
Reader study and clinical safety. The two-physician, 145-case crossover study supports initial workflow feasibility, but does not establish patient benefit. Results should be reported separately by physician, with case-order randomization, learning-effect analysis, error categories, and confidence intervals. More importantly, the system should define when it abstains or escalates to mandatory human review, and quantify the effect of these safety rules on missed diagnoses, false positives, and workload.
\end{reviewercomment}
\begin{authorresponse}
% Converted from S10_Reader_study/reviewer_response_package_20261120_claude_v17/RESPONSE_R1_6_EN.md (2026-11-17)
We thank the reviewer. We agree that the reader study supports workflow feasibility rather than patient benefit, and we have revised the Results and Discussion accordingly. We now report the reader study by physician, and we define the safety rules and evaluate them in a new study with two additional physicians.\par

\textbf{Results by physician.} The outcome is annotation consistency, that is, agreement between the diagnostic direction of the report (malignant or suspicious vs benign) and the case label derived from the frame-level annotations of the video. Annotation consistency was 72.6\% (95\% CI 61.4--81.5\%) unaided and 84.7\% (74.7--91.2\%) with AI assistance for physician 1, and 68.1\% (56.6--77.7\%) and 87.7\% (78.2--93.4\%) for physician 2 (\responsefigureref{fig_rep_reader_by_physician}). The differences were +12.1 percentage points (95\% CI $-$1.3 to +25.1) for physician 1, +19.6 (+6.1 to +32.4) for physician 2 and +15.9 (+9.7 to +22.1; McNemar p < 0.001) over the 145 cases paired by case. With AI assistance, under-calls fell from 13 to 4 of the 73 cases labelled malignant or suspicious, and over-calls from 25 to 13 of the 72 benign cases. The geometric-mean ratio of reporting time with and without AI assistance was 0.83 (0.69--1.01) for physician 1 and 0.71 (0.61--0.84) for physician 2. We added the annotation consistency by physician and the paired difference to Results~\SecResReport{}.\par

\textbf{Case order and learning effects.} The allocation of each case to the two conditions and the reading order were randomised. Each physician read the AI-assisted cases in one block and the unaided cases in a second block 8--12 days later. Within each block, annotation consistency in the second half of the cases was 3--6 percentage points higher than in the first half under both conditions, and reporting time fell by 3--11\% per 10 consecutive cases. Because the two conditions were read in different periods, the difference between them may include a period effect; we state the reading order in Methods~\SecMethReader{} and this possible period effect in the Discussion (see also our response to Reviewer 3, Major Comment \hyperref[rev:R3.4]{4}).\par

\textbf{Abstention and mandatory review.} In the manuscript, ThyroidXAgent was compared with the baselines using the report drafts generated directly by the agent, without physician review. We consider that an AI-generated report must be reviewed and finalized by a physician, and we therefore developed safety rules that facilitate this review (Supplementary Table~\SuppTabRules{}; Supplementary Fig.~\SuppFigReview{}c). We also added a study in which physicians reviewed and completed the drafts under these rules; we describe the rules and the study in Methods~\SecMethReader{} and Results~\SecResReport{}. The rules have two layers. The first layer withholds statements that the tool evidence does not support: measurements without a reliable pixel scale or outside the plausible range, vascularity without a Doppler image and nodule features without the corresponding classifier output are not written, and a lobe is called negative only when it was imaged and no nodule was detected; these items are stated as not assessed. The second layer marks each sentence of the draft for the reviewing physician: red for mandatory review (for example, TI-RADS 4--5 or conflicting evidence on a lobe), yellow for a focused check (for example, a negative lobe conclusion), grey for items withheld by the first layer, which the physician completes, and green when no safety rule is triggered.\par

\textbf{Effect on missed diagnoses, false positives and workload.} We evaluated the rules in a new study with two additional physicians, who each reviewed the drafts of 35 videos. The 70 videos were drawn at random from the 145, stratified by whether the draft stated an uncertain number of nodules, whether it contained items for the physician to complete, and video length. In addition to editing and signing off the report, the physician confirmed the number of nodules and the benign or malignant tendency of each nodule in each lobe and the isthmus; these physician-confirmed findings served as the reference for the errors of the drafts in each region. As an independent reference, we also compared the diagnostic direction of each draft and each final report with the case label from the video annotations, as in the original reader study. The drafts missed nodules in 16 of 90 regions with nodules (13 of them in regions that the draft stated could not be assessed), reported a nodule in 98 regions of which 24 had none, missed a high-risk nodule in 5 of 29 regions and overstated the risk in 63 of 87 regions reported as TI-RADS 4--5. Of these 108 errors, 106 were in sentences that the rules marked for the physician's attention (red or yellow). The physicians changed 91\% of the red and 79\% of the yellow sentences, compared with 17\% of the grey and 32\% of the green sentences (\responsefigureref{fig_rep_safety_study}a). Annotation consistency was 42.9\% (30 of 70) for the drafts (39 over-calls, no under-calls) and 74.3\% (52 of 70) for the final reports (7 over-calls, 9 under-calls), a difference of 31.4 percentage points (95\% CI 15.7 to 47.1; \responsefigureref{fig_rep_safety_study}b). In the same 70 cases, the manual and AI-assisted reports of the original reader study had an annotation consistency of 65.7\% and 84.3\% (difference 18.6 percentage points; McNemar p < 0.001). With the physician-confirmed findings as a second reference (67 cases with a clear direction), consistency was 40.3\% for the drafts, 68.7\% and 76.1\% for the manual and AI-assisted reports (difference 7.5 percentage points; McNemar p = 0.23) and 95.5\% for the final reports. Because every case contained at least one sentence requiring mandatory review, all 70 drafts were reviewed by a physician, in line with our view that every draft requires physician review. The rules are intended to show the physician where to focus the review, which may reduce repeated viewing of the video, rather than to reduce the number of reports under physician control; in this study, review took a median of 149 s per case (\responsefigureref{fig_rep_safety_study}c), and the physicians played the video about once (median playing time 1.03 times the video length; median 0 replays). The physicians corrected most overstatements of malignancy but may also have introduced understatements, indicating that AI drafts should support, rather than replace, physician review.\par

\end{authorresponse}
% Converted from S10_Reader_study/reviewer_response_package_20261120_claude_v17/figures/FIGURE_CAPTION_EN.md and S10_Reader_study/reviewer_response_package_20261120_claude_v17/figures/Fig_reader_by_physician.pdf (2026-11-17)
\begin{responsefigure}
\includegraphics[width=\textwidth,height=0.78\textheight,keepaspectratio]{report_rep_reader_by_physician.pdf}
\caption{\responsefiguretitle{fig_rep_reader_by_physician}{Reader study results by physician. \textbf{a}, Annotation consistency, that is, agreement of the diagnostic direction of the report with the annotation-derived case label; error bars, 95\% Wilson CIs; values above the bars, AI-assisted minus manual, in percentage points; ***p < 0.001, McNemar's test over the 145 cases paired by case (the differences for each physician compare different cases and were not tested). n, cases read manually/with AI assistance. \textbf{b}, Under-calls: cases labelled malignant or suspicious that were reported as benign. \textbf{c}, Over-calls: benign cases reported as malignant or suspicious. Numbers above the bars in \textbf{b} and \textbf{c}, cases with the error/cases with the label.}}
\label{fig_rep_reader_by_physician}
\end{responsefigure}

% Converted from S10_Reader_study/reviewer_response_package_20261120_claude_v17/figures/FIGURE_CAPTION_EN.md and S10_Reader_study/reviewer_response_package_20261120_claude_v17/figures/Fig_safety_study.pdf (2026-11-17)
\begin{responsefigure}
\includegraphics[width=\textwidth,height=0.78\textheight,keepaspectratio]{report_rep_safety_study.pdf}
\caption{\responsefiguretitle{fig_rep_safety_study}{Safety study (70 videos; 35 per physician). \textbf{a}, Sentences of the draft changed by the physician, by the mark of the sentence; bars, both physicians; symbols, each physician; n, sentences. Changes include added descriptions and rewording. \textbf{b}, Consistency of the diagnostic direction with two references: the annotation-derived case labels (70 cases) and the physician-confirmed findings, that is, the benign or malignant tendency of each nodule recorded by physicians 3 and 4 during review (67 cases; 3 with an ambiguous direction excluded). AI draft, the drafts reviewed in this study; manual and AI-assisted, the reports of physicians 1 and 2 in the reader study (\responsefigureref{fig_rep_reader_by_physician}); AI-assisted (safety rules), the final reports of physicians 3 and 4, who reviewed the drafts marked by the safety rules. \dag{}The physician-confirmed findings were recorded by physicians 3 and 4, who also wrote these reports. \textbf{c}, Review time per case; boxes, median and interquartile range; whiskers, 1.5 times the interquartile range; triangles, cases longer than 8 min, plotted at 8 min.}}
\label{fig_rep_safety_study}
\end{responsefigure}



\reviewersection{Reviewer 1, General remark (deployability)}
\begin{reviewercomment}
Overall, this is a thoughtful and technically ambitious study. Its strongest contribution is the framing of clinical AI as an editable and auditable evidence workflow rather than a black-box predictor. Addressing the issues above, particularly patient-level independence, routing ablations, centre-level failure analysis, clinical validation of ThyClinScore, and prospective human-AI safety evaluation, would substantially strengthen the manuscript. The work has notable strengths; however, the current evidence supports a promising retrospective workflow rather than established clinical benefit or broad deployability. I therefore recommend major revision at minimum.
\end{reviewercomment}
\begin{authorresponse}
We agree to soften deployability claims and emphasize retrospective evaluation, updated statistical reporting (data auditing and cluster-bootstrap confidence intervals), and explicit limitations.\par
% Converted from S10_Reader_study/reviewer_response_package_20261120_claude_v17/RESPONSE_R1_general_report_EN.md (2026-11-17)
For the report-generation arm, ThyClinScore is now validated against blinded ratings by two clinicians (our response to Reviewer 1, Comment \hyperref[rev:R1.5]{5}), and the safety rules were evaluated in a study with two further physicians (our response to Reviewer 1, Comment \hyperref[rev:R1.6]{6}).\par

\end{authorresponse}

\responsetoreviewer{Response to Reviewer 2}
\reviewersection{Reviewer 2, Comment 1 (R2.1)}\phantomsection\label{rev:R2.1}
\begin{reviewercomment}
Several key values are inconsistent across the abstract, main text, figures, and supplementary materials.
\end{reviewercomment}
\begin{authorresponse}
We apologize and reconciled the submitted text against the tables at the end of this letter; the revised abstract and Results will align with \responsetableref{tab_seg_ext}, \responsetableref{tab_cls_ext}, \responsetableref{tab_malig_oof}, and \responsetableref{tab_malig_ext}.\par
(1) Benign--malignant cross-cohort means. The revised abstract and Results~§2.2 report mean AUROC 0.9057 and mean AUPRC 0.8923, matching the five ThyroidXAgent cohort rows in \responsetableref{tab_cls_ext} (6807-image ZJH test; \responsetableref{tab_dhash}).\par
(2) Segmentation cross-cohort means. The revised abstract and §2.2 report mean Dice 87.01\% and mean HD95 7.46~mm, matching the seven ThyroidXAgent cohort rows in \responsetableref{tab_seg_ext} (same ZJH subset; \responsetableref{tab_dhash}).\par
(3) Main-text Fig.~3 example panels. On audit we found that in a submitted figure panel, a histologic PTC case was mislabeled as FTC. In the revision, all illustrative images in main-text Fig.~3 have been replaced with cases from the ZJH external-test hold-out (consistent with the current LNM/FTC narrative centred on ZJH), with PTC-subtype and LNM labels checked against the captions.\par
\par
% Converted from S8_Numbers_and_language/reviewer_response_package_20261120_claude_v8/RESPONSE_R2_1_report_EN.md (2026-11-17)
For the report-generation and reader-study values, we checked every value in the main text against the supplementary tables, the figures and the original reader-study records, and made the following corrections.\par

\textbf{Lexical metrics (Results~\SecResReport{}).} The values in the text now match Supplementary Table~\SuppTabNLG{} and Fig.~\FigReportGen{}f: BLEU-1, BLEU-4 and ROUGE-L of 0.5924, 0.3381 and 0.5422 on SMU-HMC and of 0.5233, 0.3018 and 0.5480 on ZJH-TS. On KMVE, the values in the text (0.6209, 0.4465, 0.3596 and 0.5826 for BLEU-1, BLEU-4, METEOR and ROUGE-L) also differed from Supplementary Table~\SuppTabNLG{}; they are now those of the revised KMVE library (0.5675, 0.4183, 0.3505 and 0.6573; our response to Reviewer 1, Comment \hyperref[rev:R1.5]{5}), which also replace the earlier values in Supplementary Table~\SuppTabNLG{} and Fig.~\FigReportGen{}f. We also state that on ZJH-TS Qwen3.5 Plus had the highest BLEU-1, BLEU-2 and METEOR (0.5379, 0.4297 and 0.3456, vs 0.5233, 0.4296 and 0.3191 for ThyroidXAgent).\par

\textbf{ThyClinScore (Results~\SecResReport{}).} In the two sentences that report ThyroidXAgent's ThyClinScore, the values for SMU-HMC and ZJH-TS are corrected from 0.5238 and 0.4775 to 0.5189 and 0.4676, as in Supplementary Table~\SuppTabClinical{} and Fig.~\FigReportGen{}f, and the KMVE value (0.4465) is replaced by that of the revised KMVE library (0.4374).\par

\textbf{Reporting time (Results~\SecResReport{}, Fig.~\FigReader{}d,e).} The 27.4\% reduction in the Abstract, Results~\SecResReport{} and Fig.~\FigReader{}d had been calculated over all 151 videos of the reader study, including 6 videos without benign or malignant annotations, whereas the other reader-study results use the 145 annotated videos. Over these 145 videos, the mean reporting time fell from 149.3 s to 110.5 s (2.5 to 1.8 min), a reduction of 26.0\%; the Abstract, the text and Fig.~\FigReader{}d now give this value, so that all reader-study results refer to the same 145 videos. The time savings of the two physicians differed (13.0\% and 36.8\%; Fig.~\FigReader{}e), and the text no longer describes them as similar.\par

\textbf{Figure citation (Methods~\SecMethReport{}).} The sentence on building template libraries from new report sources now cites Supplementary Fig.~\SuppFigTemplateBank{}, which shows template-library construction, instead of Supplementary Fig.~\SuppFigInterface{}c, which shows video-input processing.\par

\end{authorresponse}
% [SYNC] tab_malig_oof: malig_models_zjh20261101.json (2026-11-03 PM); LLNM/Tiger/ThyroidXAgent (DINO interim) my_run image-level OOF filled.
\begin{manuscriptresponsetable}
\caption{\responsetabletitle{tab_malig_oof}{Malignant lesion stratification: development five-fold OOF}}
\begin{adjustbox}{max width=\textwidth,center}
\begin{tabular}{lcccc}
\toprule
\multirow{2}{*}{\textbf{Method}}
& \multicolumn{2}{c}{\textbf{Central lymph-node metastasis (CN0/CN1)}}
& \multicolumn{2}{c}{\textbf{Classic PTC \emph{vs.} follicular-variant PTC}} \\
& AUROC $\uparrow$ & AUPRC $\uparrow$ & AUROC $\uparrow$ & AUPRC $\uparrow$ \\
\midrule
RepViT & $0.8334 \pm 0.0302$ & $0.8305 \pm 0.0347$ & $0.7005 \pm 0.0471$ & $0.6973 \pm 0.059$ \\
LSNet & $0.7090 \pm 0.0736$ & $0.7083 \pm 0.0925$ & $0.6484 \pm 0.047$ & $0.6773 \pm 0.0619$ \\
UltraFedFM & $0.8568 \pm 0.0402$ & $0.8627 \pm 0.039$ & $0.7812 \pm 0.0393$ & $0.7620 \pm 0.0535$ \\
LLNM-Net & $0.8458 \pm 0.0726$ & $0.8369 \pm 0.0875$ & --- & --- \\
Tiger-Model & --- & --- & $0.7877 \pm 0.0466$ & $0.7044 \pm 0.0581$ \\
\responserowagent
\textbf{ThyroidXAgent} & $0.8797 \pm 0.0285$ & $0.8493 \pm 0.0375$ & $0.8089 \pm 0.027$ & $0.8271 \pm 0.0358$ \\
\bottomrule
\end{tabular}
\end{adjustbox}
\responsetablenote{Training/development set: shared ZJH-8K malignant primary-lesion pool ($3227$ images, $517$ patients; TN3K train--ZJH-8K dHash H=0 exclusion and full-pool screening in \responsetableref{tab_dhash}; counts and split definition in \responsetableref{tab_cohort_patient_counts}). LNM: case-level central-compartment status (CN0/CN1) from postoperative pathology; primary-lesion ultrasound inputs only (\hyperref[rev:R2.12]{R2.12}). PTC-subtype labels: pathology categories 1 (classic) \emph{vs.} 2 (follicular-variant PTC) only. Five-fold OOF on the development pool (seed 42). AUROC/AUPRC; $\pm$ = half-width of bootstrap 95\% CI (1000 resamples, seed 42). ``---'' = task not evaluated.}
\end{manuscriptresponsetable}


\reviewersection{Reviewer 2, Comment 2 (R2.2)}\phantomsection\label{rev:R2.2}
\begin{reviewercomment}
Many analyses are reported at the image level, although multiple images may come from the same patient, examination, or nodule. For example, ZJH-8K includes 7,958 images from 1,149 cases. Treating correlated images as independent observations may produce overly narrow confidence intervals. The authors should recalculate confidence intervals using patient- or examination-level bootstrap resampling and report patient-level performance for clinically relevant classification tasks.
\end{reviewercomment}
\begin{authorresponse}
We agree that treating multiple images from the same patient or examination as independent can yield anti-conservatively narrow confidence intervals and misaligns with clinical decision units.\par
Cluster-bootstrap 95\% CIs (1000 resamples, seed~42) in \responsetableref{tab_cls_ext}, \responsetableref{tab_seg_ext}, and \responsetableref{tab_malig_oof}: where reproducible patient IDs exist (ZJH-8K, ThyroidXL), bootstrap resamples patient clusters with replacement before expanding to images; cohorts without IDs (e.g., TN3K, TN5K, DDTI) use one image per cluster. For ID-missing cohorts, cluster-bootstrap CIs match per-image bootstrap (point estimates unchanged); only ZJH-8K and ThyroidXL show slightly wider image-level CIs under cluster resampling. We additionally report patient-level classification and segmentation for those two cohorts (\responsetableref{tab_patient_cls}); other benign--malignant external tests remain image-level in \responsetableref{tab_cls_ext}.\par
Task-level reporting under the protocol above (cohort image/patient counts in \responsetableref{tab_cohort_patient_counts}): (i)~Benign--malignant classification---image-level external metrics in \responsetableref{tab_cls_ext}; patient-level ThyroidXAgent for ZJH-8K and ThyroidXL in \responsetableref{tab_patient_cls} (same table includes segmentation metrics; 1068 and 739 external-test patients). (ii)~Segmentation---image-level Dice/HD95 in \responsetableref{tab_seg_ext}; patient-level metrics in \responsetableref{tab_patient_cls}. (iii)~Malignant subtyping (LNM; classic PTC \emph{vs.} follicular-variant PTC)---development five-fold OOF in \responsetableref{tab_malig_oof}; independent external tests in \responsetableref{tab_malig_ext} on the shared ZJH malignant patient hold-out (\hyperref[rev:R2.10]{R2.10}, \hyperref[rev:R2.12]{R2.12}), evaluated at patient/cluster level as in those tables.\par
\noindent\par
\par
\end{authorresponse}

\reviewersection{Reviewer 2, Comment 3 (R2.3)}\phantomsection\label{rev:R2.3}
\begin{reviewercomment}
Please clarify why 28,458 observations are described as "non-overlapping test cases," as this number appears to combine 19,737 test images with 8,721 NHC-MISD-TUS cases.
\end{reviewercomment}
\begin{authorresponse}
We agree that the original ``28,458 non-overlapping test cases'' wording was unclear: it combined OpenThyroid and NHC-MISD-TUS using inconsistent units (images versus cases) and did not separate image analysis, report generation, and malignant-subtask evaluations. We removed this single aggregate from the revised manuscript. The abstract now states image-level test sizes for the main segmentation/classification benchmark---18,586 test images across seven OpenThyroid cohorts (Methods; Supplementary Table~S1 of the revised manuscript)---and, separately, 8,721 gland-segmentation test images from 35 contributing sites in NHC-MISD-TUS (\hyperref[rev:R2.6]{R2.6}). Cohort roles, splits, and image-analysis external-test counts are summarized in \responsefigureref{fig_image_dataflow} and \responsetableref{tab_cohort_patient_counts}; report-generation and reader-study cohorts are defined in Methods and are not summed into the image-benchmark totals above.\par
\par
\end{authorresponse}

\reviewersection{Reviewer 2, Comment 4 (R2.4)}\phantomsection\label{rev:R2.4}
\begin{reviewercomment}
ZJH-8K and TN3K were both collected at Zhujiang Hospital, Southern Medical University. Unless they are demonstrably independent in patient population, collection period, cohort composition, and clinical workflow, ZJH-8K should not be described as an external-centre cohort.
\end{reviewercomment}
\begin{authorresponse}
We agree. ZJH-8K and TN3K were collected at the same hospital (Zhujiang Hospital, Southern Medical University) and should not be described as an independent \emph{external centre} cohort; we describe ZJH-8K as a \emph{same-centre institutional hold-out / test cohort} in the revised manuscript.\par
We did not perform a systematic hospital-ID reconciliation to adjudicate patient population, collection period, or workflow independence: TN3K's public release does not provide the acquisition-year or patient-linkage fields needed for that audit (ZJH-8K is documented as a 2016--2023 institutional cohort with pathology-linked labels in Methods). For reported ZJH external metrics, cross-source leakage control is therefore grounded in \emph{image-level} dHash near-duplicate screening (\responsetableref{tab_dhash}; same protocol as \hyperref[rev:R1.1]{R1.1}).\par
Hamming $\leq$ 5 screening of each foundational public \emph{train} partition against the full ZJH-8K pool (7958 images) shows that overlap material to TN3K-train-then-test-on-ZJH is concentrated in TN3K train vs ZJH-8K (1154 H=0 pairs, 1151 unique ZJH images; TN5K, ThyroidXL, and PKTN train show 0 H=0 pairs to ZJH-8K). ZJH-8K segmentation/classification and ZJH malignant-subtask development metrics in this letter exclude those 1151 images (6807 external-test images; 4028 LNM/PTC-subtype ZJH primary-lesion development images), as in \responsetableref{tab_cls_ext}, \responsetableref{tab_seg_ext}, and \responsetableref{tab_malig_oof}.\par
\par
\end{authorresponse}

\reviewersection{Reviewer 2, Comment 5 (R2.5)}\phantomsection\label{rev:R2.5}
\begin{reviewercomment}
The Methods should specify the reference standards for nodule segmentation and benign-malignant classification, annotation procedures, annotator qualifications and adjudication of disagreements, patient-level deduplication, and handling of missing data.
\end{reviewercomment}
\begin{authorresponse}
We agree that Methods must document reference standards, annotation procedures, annotator qualifications and disagreement resolution, patient-level deduplication, and handling of missing data for all nodule segmentation and benign--malignant classification cohorts. For public OpenThyroid cohorts (TN3K, TN5K, ThyroidXL, PKTN, DDTI, RJH-7K, \textit{etc.}), we describe segmentation masks, benign--malignant labels, and official splits by citing each dataset publication and release in Methods, without restating original annotation-room protocols in this letter. For \emph{ZJH-8K}---the only in-house cohort used here as an institutional segmentation/classification external test set---benign--malignant labels follow postoperative histopathology (with selective fine-needle aspiration cytology where applicable), tied to the \emph{target nodule on each image}; segmentation reference masks outline the corresponding nodule region on thyroid ultrasound, delineated by hospital ultrasound or radiology physicians with senior review when disagreements arise. The revised manuscript adds a dedicated ZJH-8K Methods subsection with these protocols.\par
Split roles are summarized in \responsefigureref{fig_image_dataflow}. Where reproducible patient IDs exist (ZJH-8K, ThyroidXL), all images and views under the same patient ID remain in a single official or predefined split, and we do \emph{not} place the same patient in both the stacked OpenThyroid expert-training partitions and external testing (\hyperref[rev:R1.1]{R1.1}). Cohorts without public patient IDs (e.g., TN3K, TN5K, DDTI) follow each release's official \emph{image-level} train/val/test lists as cited in Methods. We \emph{do not impute} masks or labels: nodule-segmentation evaluation includes only images with valid reference masks, and benign--malignant classification includes only images with valid binary labels; otherwise images are \emph{excluded} rather than assigned default labels or synthetic masks. Task--cohort mismatches are handled by omission---e.g., DDTI has 637 external images but only \emph{349} carry benign--malignant labels, so classification metrics use that subset; RJH-7K contributes segmentation external testing only and has no benign--malignant labels in this study. Corrupt files, mask--image mismatches, or missing label fields are dropped before pooling. We do not add a separate per-cohort exclusion table in this letter; evaluable sample sizes are in the external-test tables and \responsetableref{tab_cohort_patient_counts}.\par
Before inference and mask alignment, we detect and crop uniform black letterbox margins on frame borders (intensity-threshold based) under one rule across cohorts so borders are not fed to expert inputs; parameters are fixed in the preprocessing pipeline and summarized in Methods.\par
\par
\end{authorresponse}

\reviewersection{Reviewer 2, Comment 6 (R2.6)}\phantomsection\label{rev:R2.6}
\begin{reviewercomment}
The manuscript reports validation on 8,721 cases from 35 centres, but the NHC-MISD-TUS cohort is insufficiently described. The authors should clarify whether these represent 35 independent hospitals and report the numbers of patients, nodules, and images per centre, together with centre-level macro-averaged performance and between-centre heterogeneity.
\end{reviewercomment}
\begin{authorresponse}
We agree that NHC-MISD-TUS must be described clearly. Cohort definition, image counts and disclosure limits are in the revised manuscript Methods; this letter adds centre-level results. NHC-MISD-TUS is a private, multi-centre, test-only thyroid ultrasound subset authorized from the National Health Commission Medical Imaging Standard Database; it was not used to train the stacked segmentation or classification experts. The release comprises 35 anonymized contributing site identifiers (Centre ID) and image-level external-test counts: \emph{8{,}721} gland-segmentation test images, with \emph{6{,}384} nodule-segmentation and \emph{4{,}999} benign--malignant classification images where labels were available (per-centre evaluable counts by gland, nodule, and classification task in \responsetableref{tab_nhc_center_performance}). Hospital names, patient-level totals, nodule-level totals, and scanner inventories are not disclosed to the research team under the governing de-identification and data-use conditions of this multicentre resource.\par
Participating units submit de-identified cases through a standardized structured workflow on the database platform (clinical history, protocol ultrasound views, reports, and pathology where available); the platform performs quality control and produces task-specific reference labels used in this study. Centre-specific hospital identities or device metadata were not released, so we do not report per-centre patient or nodule counts or perform formal meta-regression on scanner type or case mix.\par
Per-centre evaluable image counts, ThyroidXAgent external-test performance, and centre-level meta heterogeneity are in \responsetablerefrange{tab_nhc_center_performance}{tab_nhc_heterogeneity}. Unweighted \emph{macro} and inverse-variance \emph{Overall pooled} summaries are on the same tables (e.g., gland Dice macro \emph{55.94\%}, pooled \emph{59.28\%}, 35 sites, $I^2{\approx}99.9\%$; nodule Dice macro \emph{81.64\%}, pooled \emph{82.31\%}, 33 sites, $I^2{\approx}99.0\%$; benign--malignant AUROC macro \emph{0.754}, pooled \emph{0.819}, 32 sites, $I^2{\approx}99.0\%$). Centres with non-estimable AUROC are marked NE in \responsetableref{tab_nhc_center_performance} and excluded from the AUROC macro average.\par
\par
\end{authorresponse}

\reviewersection{Reviewer 2, Comment 7 (R2.7)}\phantomsection\label{rev:R2.7}
\begin{reviewercomment}
The authors should clarify whether the reported "mean" performance represents an unweighted macro-average across cohorts, a sample-size-weighted average, or another calculation.
\end{reviewercomment}
\begin{authorresponse}
We thank the reviewer. For benign--malignant classification and nodule segmentation alike, cross-cohort ``means'' are unweighted macro-averages: we compute the task metric on each external test cohort, then average cohort-level point estimates with equal weight (not weighted by test image count). Per-cohort metrics are in \responsetableref{tab_cls_ext} (classification) and \responsetableref{tab_seg_ext} (segmentation); macro AUROC 0.9057, AUPRC 0.8923, Dice 87.01\%, and HD95 7.46~mm follow from the ThyroidXAgent rows in those tables (6807-image ZJH subset; \responsetableref{tab_dhash}), reconciled with Reviewer 2, Comment \hyperref[rev:R2.1]{1} (\hyperref[rev:R2.1]{R2.1}).\par
\par
\end{authorresponse}

\reviewersection{Reviewer 2, Comment 8 (R2.8)}\phantomsection\label{rev:R2.8}
\begin{reviewercomment}
The key innovation is the agentic evidence workflow, but the observed gains may largely reflect the DINOv3 expert pool, multicohort training, radiomics classifiers, or template-based reporting rather than agentic routing itself. As the static-pipeline comparison is limited mainly to report generation, the independent value of LLM-based routing for segmentation and classification remains unclear.
\end{reviewercomment}
\begin{authorresponse}
We agree that module contributions must be separated. The full eight-strategy routing ablation is in Reviewer 1, Comment \hyperref[rev:R1.2]{2} (\responsetablerefrange{tab_routing_ablation_seg}{tab_routing_ablation_cls}); strategy (7) is deployable full-system ThyroidXAgent (external-test protocol in \responsetableref{tab_baseline_protocol}), and supervised baselines are under Reviewer 2, Comment \hyperref[rev:R2.14]{14} (\responsetableref{tab_baseline_protocol}). On the fixed expert pool with frozen candidates, deployable ThyroidXAgent (7) is compared with static strategies (1) and (2); ensembling is not universally best, and gains should not be attributed to the expert pool alone or read as ``static single expert beats ThyroidXAgent.'' Cohort-level point estimates and 95\% CIs are in \responsetablerefrange{tab_routing_ablation_seg}{tab_routing_ablation_cls}. Unweighted arithmetic means of strategy (1), (2), and (7) point estimates over listed segmentation/classification cohorts are: segmentation Dice 0.8576, 0.8487, 0.8701; classification AUROC 0.8921, 0.8727, 0.9057 (no pooled CI). Auditable LLM routing and repeat-call stability are in Reviewer 2, Comment \hyperref[rev:R2.9]{9} (\responsetableref{tab_routing_stability}).
\par
\end{authorresponse}
% [SYNC] tab_routing_stability: seg_selection_agreement from output/b6/{tn3k,tn5k,thyroidxl,zjh8k,ddti_seg}/pipeline_run/stability_summary.json (n=10, seed 42, 5 repeats).
\begin{responsetable}
\caption{\responsetabletitle{tab_routing_stability}{Segmentation LLM routing stability}}
\begin{adjustbox}{max width=\textwidth,center}
\begin{tabular}{lccc}
\toprule
External cohort & Sampled $n$ & Mean five-run agreement & Per-case range \\
\midrule
TN3K & 10 & 100\% & 100\%--100\% \\
TN5K & 10 & 100\% & 100\%--100\% \\
ThyroidXL & 10 & 94\% & 60\%--100\% \\
ZJH-8K & 10 & 98\% & 80\%--100\% \\
DDTI & 10 & 100\% & 100\%--100\% \\
\bottomrule
\end{tabular}
\end{adjustbox}
\responsetablenote{Per cohort: 10 cases drawn from the official test set (seed 42), preferring the similar tier (top-two mask quality gap $\leq 0.10$); frozen upstream segmentation and classification expert outputs; five independent segmentation LLM routing calls per case. Per-case agreement is the fraction of five runs that select the same segmentation expert; cohort mean agreement is the arithmetic mean over the 10 cases, and the per-case range is the minimum--maximum over those 10 cases.}
\end{responsetable}


\reviewersection{Reviewer 2, Comment 9 (R2.9)}\phantomsection\label{rev:R2.9}
\begin{reviewercomment}
The manuscript states that the LLM router selects the best expert output through "reasoning," but does not report the prompt design, LLM version, decoding parameters, reproducibility, or decision stability when candidate outputs have similar quality. The router therefore remains a black-box selector. The authors should provide structured routing records for at least 10 representative cases, including input quality metrics, candidate outputs, final selections, uncertainty indicators, brief decision rationales, and consistency across repeated runs.
\end{reviewercomment}
\begin{authorresponse}
Appendix A1 lists model/decoding parameters, compact JSON prompts used in external testing, required output fields, and structured routing-record items; records capture candidate quality, selected experts, and confidence---no free-text rationales are requested from the LLM. During external testing on the five BM cohorts that deploy LLM routing, every segmentation/classification step that invokes the LLM is written to auditable structured records (Appendix A1); cohorts without LLM deployment (e.g., NHC-MISD-TUS) are excluded.\par
Similar-candidate tier: candidates are ranked by inter-mask agreement; if the quality gap between the top two is $\leq 0.10$ the tier is ``similar,'' otherwise ``easy.'' \responsetableref{tab_routing_stability} reports LLM consistency on the similar tier (this 10-case draw is all similar-tier). How each record reports the selected expert and similarity tier is defined in Appendix A1.\par
LLM decision stability (multi-cohort, similar candidates): with frozen upstream expert outputs, each external cohort that uses segmentation LLM routing draws 10 cases from its official test set with fixed random seed 42, preferring the ``similar'' tier; we repeat segmentation LLM routing five times per case and report the arithmetic mean of per-case agreement across the 10 draws (with per-case ranges in the table). Summaries are in \responsetableref{tab_routing_stability} ($n{=}10$ per cohort, five repeats per case; protocol in the table note). When all independent classifiers already share the same top class, classification uses deterministic consensus (no LLM) and is outside this repeat experiment.\par
De-identified routing-record exemplars Rec-01--Rec-10 are at the end of Appendix~A1 (field definitions in Table~A1-2), without disk filenames as public case identifiers.\par
\par
\end{authorresponse}

\reviewersection{Reviewer 2, Comment 10 (R2.10)}\phantomsection\label{rev:R2.10}
\begin{reviewercomment}
For FTC/PTC classification, the training cohort comprised classic PTC and follicular-variant PTC, whereas the external test set comprised FTC and PTC. Because follicular-variant PTC is not equivalent to FTC, the target definition differs between training and testing. Unless this is justified as a clinically valid domain-transfer task, the reported performance may not be valid.
\end{reviewercomment}
\begin{authorresponse}
We agree with the reviewer: the submitted manuscript implied a train--test mismatch (classic PTC \emph{vs.} follicular-variant PTC on ZJH versus histologic PTC \emph{vs.} FTC on a public external set) that would invalidate pooled FTC/PTC reporting.\par
We clarify the \emph{PTC-subtype} subtask on ZJH-8K only. Pathology records four papillary subtypes, but we define a \emph{binary} subtask---category 1 (classic papillary thyroid carcinoma) \emph{vs.} 2 (follicular-pattern papillary carcinoma, follicular variant)---not four-way classification. Categories 3 (oncocytic) and 4 (clear-cell) each comprise only single-digit \emph{patients} in the ZJH malignant primary-lesion pool and are excluded from the label space. The binary label indexes classic \emph{vs.} follicular-variant PTC on postoperative histology of the index lesion---not histologic follicular thyroid carcinoma \emph{vs.} papillary carcinoma. Development five-fold OOF and the shared patient hold-out after TN3K dHash exclusion (130 patients, 801 primary-lesion images; seed~42) use the \emph{same} label definition (\responsetableref{tab_malig_oof}, \responsetableref{tab_malig_ext}). At the patient level, LNM was CN0/CN1 in 377/140 (development) and 95/35 (hold-out); PTC subtype was classic/follicular-variant in 466/51 and 117/13. This letter does not include cross-dataset histologic PTC-versus-FTC external testing; reporting is confined to ZJH under the definition above.\par
\par
\end{authorresponse}

\reviewersection{Reviewer 2, Comment 11 (R2.11)}\phantomsection\label{rev:R2.11}
\begin{reviewercomment}
The LNM task uses only 4,753 malignant images, and Table S5 indicates only 20 cases for validation---sample sizes far too small for training reliable deep learning classifiers. Please clarify: (a) whether transfer learning or multi-task learning was employed; (b) whether data augmentation was applied; (c) how overfitting is avoided given the wide confidence intervals in Table S5 that reflect inadequate sample size.
\end{reviewercomment}
\begin{authorresponse}
We clarify the LNM development pool and validation design. \responsetableref{tab_malig_oof} reports patient/cluster-stratified five-fold OOF (seed 42) on the post-exclusion ZJH malignant development cohort, replacing the original 20-case hold-out; primary independent external tests are in \responsetableref{tab_malig_ext} on the shared ZJH patient hold-out (\hyperref[rev:R2.12]{R2.12}). ThyroidXAgent LNM development OOF in \responsetableref{tab_malig_oof}: AUROC 0.8797 $\pm$ 0.0285, AUPRC 0.8493 $\pm$ 0.0375.\par
(a) Transfer or multitask learning. We did not transfer weights from the generic benign--malignant segmentation/classification stack to LNM; LNM uses a task-specific binary classifier on ZJH-8K malignant \emph{primary-lesion} ultrasound with evaluation separate from the main image arm (\responsetableref{tab_malig_oof}, \responsetableref{tab_malig_ext}).\par
(b) Data augmentation. The LNM task-specific DINOv3 image expert uses the same on-the-fly training augmentation as our other supervised classification experts on the OpenThyroid stacked development sets: resize image (and training lesion mask) to the expert input size, synchronized random horizontal and vertical flips ($p{=}0.5$ each), then ImageNet mean/std normalization; we did not add LNM-only augmentations. Development data and fine-tuning protocol are in Reviewer 2, Comment \hyperref[rev:R2.14]{14} (\responsetableref{tab_baseline_protocol}).\par
(c) Overfitting mitigation. We report patient/cluster-stratified five-fold OOF on the development set (not a single small validation split), cluster-bootstrap CIs (\hyperref[rev:R2.2]{R2.2}), and a held-out \emph{ZJH patient external test} on primary-lesion ultrasound (\responsetableref{tab_malig_oof}, \responsetableref{tab_malig_ext}).\par
\par
\end{authorresponse}

\reviewersection{Reviewer 2, Comment 12 (R2.12)}\phantomsection\label{rev:R2.12}
\begin{reviewercomment}
For lymph-node metastasis prediction, the manuscript inconsistently refers to overall LNM, lateral LNM, and CN0/CN1 labels. The authors should clarify whether CN denotes central-compartment nodal status, whether the target is central, lateral, or any cervical LNM, whether the inputs include primary thyroid lesions, lymph nodes, or both, and how image-level predictions are aggregated at the patient level.
\end{reviewercomment}
\begin{authorresponse}
We agree that LNM terminology and task definition must be explicit. On ZJH-8K malignant \emph{primary thyroid ultrasound}, the binary label (CN0/CN1) is the case-level postoperative pathology status of \emph{central-compartment} lymph-node metastasis (CN0 = absent, CN1 = present); it does not encode lateral-neck status alone. Inputs are primary-lesion images only, not dedicated cervical lymph-node ultrasound.\par
Patient-level aggregation (LNM / FTC/PTC on ZJH primary lesions). Each image yields an image-level class probability $p_i$ from the task-specific classifier. \emph{Patient clusters} are defined by reproducible \emph{patient IDs} (multiple primary-lesion images under the same patient belong to one cluster; TN3K dHash exclusion and train/test splits are in \responsetableref{tab_dhash}, \responsetableref{tab_malig_oof}, and \responsetableref{tab_malig_ext}). The \emph{patient score} is the \emph{arithmetic mean} of image-level probabilities within the cluster, $\bar p=\mathrm{mean}(\{p_i\})$; the \emph{patient label} is the case-level pathologic LNM (CN0/CN1) or PTC-subtype label, constant within each cluster. AUROC/AUPRC for five-fold development OOF and the shared ZJH patient hold-out are computed on \emph{patient-level} pairs $(\bar p,\ \mathrm{label})$; 95\% CIs use cluster bootstrap resampling at the \emph{patient cluster} (1000 resamples, seed 42), as in Reviewer 2, Comment \hyperref[rev:R2.2]{2} (\hyperref[rev:R2.2]{R2.2}). Five-fold stratification and the hold-out are \emph{patient-ID--exclusive}, so no patient appears in multiple folds or in both development and hold-out. At the patient level, LNM was CN0/CN1 in 377/140 (development) and 95/35 (hold-out); PTC subtype was classic/follicular-variant in 466/51 and 117/13.\par
LNM and PTC-subtype evaluation use only ZJH-8K malignant \emph{primary-thyroid} ultrasound with shared patient-stratified development and hold-out (\responsetableref{tab_malig_oof}, \responsetableref{tab_malig_ext}).\par
Development five-fold OOF and external-test protocols are in \responsetableref{tab_malig_oof} and \responsetableref{tab_malig_ext}.\par
\par
\end{authorresponse}

\reviewersection{Reviewer 2, Comment 13 (R2.13)}\phantomsection\label{rev:R2.13}
\begin{reviewercomment}
The report-generation reader study included only two physicians, with each case assessed under the two conditions by different readers. Although this may reduce recall bias, it may also confound reader effects with the assistance condition.
\end{reviewercomment}
\begin{authorresponse}
% Converted from S10_Reader_study/reviewer_response_package_20261120_claude_v17/RESPONSE_R2_13_EN.md (2026-11-17)
We thank the reviewer. Each physician read about half of the cases with AI assistance and the other half without, and each case was read once in each condition, by different physicians. Both physicians therefore contributed equally to both conditions, so a difference between the physicians does not by itself produce a difference between the conditions. In a linear model with the physician as a covariate, annotation consistency was 15.9 percentage points higher with AI assistance (95\% CI 10.1 to 22.4), and the difference between the two physicians, balanced for condition, was $-$0.8 percentage points (95\% CI $-$7.0 to +5.4); \responsefigureref{fig_rep_reader_by_physician} shows the annotation consistency of each physician in both conditions. Both physicians, however, read their AI-assisted cases before their unaided cases, so the condition cannot be separated from the reading period (our response to Reviewer 1, Comment \hyperref[rev:R1.6]{6}), and two physicians do not represent the variation among physicians. The safety study (our response to Reviewer 1, Comment \hyperref[rev:R1.6]{6}) added two further physicians, who reviewed the AI drafts of 70 of these cases with the safety rules; for each of them, the annotation consistency of the final reports was higher than that of the drafts (by 28.6 and 34.3 percentage points), although that study had no unaided condition. We state both limitations in the Discussion.\par

\end{authorresponse}

\reviewersection{Reviewer 2, Comment 14 (R2.14)}\phantomsection\label{rev:R2.14}
\begin{reviewercomment}
It is unclear whether all specialist baselines were retrained on the same development data, evaluated using their original weights, or tested zero-shot. Comparisons may be unfair if ThyroidXAgent benefits from multicohort training and task-specific models while the baselines are used off the shelf.
\end{reviewercomment}
\begin{authorresponse}
We agree that baseline training and inference must be stated explicitly for a fair comparison. Supervised specialist rows in \responsetableref{tab_cls_ext} and \responsetableref{tab_seg_ext}, and the ThyroidXAgent image experts underlying our external results, are \emph{not} thyroid zero-shot tests and \emph{not} frozen-upstream evaluation only: each is fully fine-tuned from publicly available upstream weights on the same OpenThyroid stacked development partitions (\responsetableref{tab_baseline_protocol}). Designated external cohorts (e.g.\ DDTI, ZJH-8K) are excluded from that development fine-tuning for both baselines and ThyroidXAgent experts. MedGemma and Qwen3-VL follow the same supervised classification fine-tuning protocol as the other rows in the table; GPT-5 and Gemini use vendor APIs with fixed prompts and are not fine-tuned on our thyroid development data. ThyroidXAgent's deployable external-test arm (\hyperref[rev:R2.8]{R2.8}, strategy (7)) uses the same expert-development protocol; routing and radiomics are assessed on frozen expert outputs in the eight-strategy ablation (Reviewer 2, Comment \hyperref[rev:R2.8]{8}; \responsetablerefrange{tab_routing_ablation_seg}{tab_routing_ablation_cls}).\par
\par
\end{authorresponse}

\reviewersection{Reviewer 2, Comment 15 (R2.15)}\phantomsection\label{rev:R2.15}
\begin{reviewercomment}
The manuscript provides an ethics approval number but does not state whether informed consent was obtained or waived. For all institutional datasets, the authors should report the study design and period, consent or waiver arrangements, de-identification procedures, and data-governance measures. Authorization by a co-author to use NHC-MISD-TUS does not replace documentation of ethical approval and regulatory compliance.
\end{reviewercomment}
\begin{authorresponse}
We agree that ethical approval, consent or waiver arrangements, de-identification, and data governance must be documented for all non-public data sources used in this study.\par
For institutional thyroid ultrasound and report resources at Zhujiang Hospital, Southern Medical University (including ZJH-8K and related cohorts used in segmentation, classification, and malignant-subtyping analyses in this letter), the study is retrospective. The protocol was approved by the Medical Ethics Committee of Zhujiang Hospital, Southern Medical University (approval no.\ 2026-KY-081-01). The ethics committee waived the requirement for written informed consent for this retrospective use of de-identified clinical imaging and linked labels, under the conditions stated in the approved protocol. Images and reports were de-identified before research use, and access was limited to approved investigators under institutional data-governance rules.\par
The report-generation cohorts SMU-HMC and ZJH-TS were also collected retrospectively at Zhujiang Hospital and were used under the same approval, consent waiver, de-identification and access restrictions. SMU-HMC comprises reports from January to June 2025. ZJH-TS was drawn from 353 reports dated May 2016 to October 2025, 300 of them from 2024 and 2025.\par
NHC-MISD-TUS is an authorized, de-identified thyroid ultrasound subset from the National Health Commission Medical Imaging Standard Database. Access is governed by an executed cooperation and data-compliance agreement with the database programme, not through individual co-author authorization. In this study the release was used \emph{only} for external testing of segmentation and classification and was not used to train the stacked OpenThyroid experts, under the de-identification and data-use conditions summarized in Reviewer~2, Comment~\hyperref[rev:R2.6]{6} (\hyperref[rev:R2.6]{R2.6}).\par
\end{authorresponse}

\reviewersection{Reviewer 2, Comment 16 (R2.16)}\phantomsection\label{rev:R2.16}
\begin{reviewercomment}
The ReAct-style planner-executor workflow lacks systematic failure-mode evaluation. The authors should report the frequency of invalid or unnecessary tool calls, hallucinated tasks, repeated calls, execution failures, retries, and rollbacks, and provide representative traces for successful cases and major failure modes, particularly cases with no nodules, multiple nodules, or conflicting intermediate evidence.
\end{reviewercomment}
\begin{authorresponse}
% Converted from S5_Tool_call_audit/reviewer_response_package_20261119_claude_v13/RESPONSE_R2_16_EN.md (2026-11-17)
We thank the reviewer for this comment and agree that the original manuscript did not evaluate failure modes of the planner--executor workflow. We have therefore performed a fully instrumented failure-mode audit (\responsetableref{tab_rep_failure_modes} and \responsefigureref{fig_rep_execution_traces}) and summarize its main results in Results~\SecResReport{}.\par

\textbf{Audit design.} We audited 150 report-generation cases, 50 from each of SMU-HMC, KMVE and ZJH-TS. The cases were drawn from a sample stratified by the number of thyroid nodules or cysts stated in the reference report: 0, 1, $\geq$2 or uncertain in 11, 12, 23 and 4 cases in SMU-HMC, 12, 11, 23 and 4 in KMVE, and 1, 11, 38 and 0 in ZJH-TS. Each case was run through the complete workflow three times, giving 450 runs and 2,347 tool calls. We used a root seed of 42 to deterministically derive a pre-assigned seed for each run. For every run we recorded the LLM requests and responses, each tool call with its arguments and raw return, the summary returned to the LLM, the case state before and after each call, and the final report. Each failure category was defined operationally and assessed for all tool calls (\responsetableref{tab_rep_failure_modes}). Because each case was run three times, we also report how many cases were affected and whether an event recurred across a case's runs (\responsetableref{tab_rep_failure_modes}). Confidence intervals (CIs) were obtained by resampling cases within each dataset, keeping all three runs of each case.\par

\textbf{Frequency of the requested failure categories} (\responsetableref{tab_rep_failure_modes}).\par

\begin{itemize}[leftmargin=1.5em,itemsep=1pt,topsep=2pt]
\item \textbf{Invalid calls.} Four of the 2,347 tool calls (0.17\%; 95\% CI 0.04--0.35\%) were invalid. They occurred in four different cases, each in only one of its three runs. One call used a malformed tool name. The other three were gland-measurement requests that cited image identifiers that do not exist: one was a zero-padded variant of an existing identifier, and two were numbered one beyond the last image. The executor did not process the non-existent identifiers.
\item \textbf{Unnecessary calls.} We counted requests that listed images outside the tool's target region: 18 lymph-node screening calls that included non-neck images (0.77\% of all calls; 95\% CI 0.27--1.37\%; 8 cases) and 66 gland-measurement calls that included neck images (2.81\%; 1.96--3.70\%; 32 cases). In both situations the executor passes only the eligible images to the tool, so the other images are simply not analysed in that call; nothing is deleted. For example, one screening request listed ten images, and only its five neck images were screened (\responsefigureref{fig_rep_execution_traces}e). Because each of these calls still analysed eligible images, we regard them as imprecise requests rather than unnecessary calls. Region labels come from automated preprocessing.
\item \textbf{Hallucinated tasks.} We defined a hallucinated task as a task that was planned or executed without support in the case's images. Every one of the 2,347 tool calls executed a task in the plan. We then checked each of the 2,348 planned tasks against the images available at planning time; for example, nodule segmentation requires an image flagged by preprocessing, and lymph-node screening requires a neck image. Three planned tasks (0.13\%; 95\% CI 0.00--0.29\%) had no such support. All three were lymph-node screening planned for cases without neck images, each in a different case and in only one of its three runs. In two of these runs the task was never executed; in the third it was issued with an empty image list and returned nothing.
\item \textbf{Repeated calls.} Across the 450 runs, apart from one retry triggered by a malformed tool name, no tool was called more than once within the same run.
\item \textbf{Execution failures.} The only failed tool call was the malformed name. All 450 runs were completed. Three of them, in three different cases, were first interrupted by out-of-memory errors caused by an external factor (other jobs occupying GPU memory) and were re-run with the same inputs.
\item \textbf{Retries and recovery.} The call with the malformed tool name was retried by the LLM at the next step, with the corrected name and the same arguments. In addition, the executor has a fallback for steps at which the LLM does not continue with the plan produced by the planner, that is, returns neither a tool call nor the completion signal: it then runs the next planned task that has not yet been performed. This fallback was triggered in two runs, each from a different case. In both, the LLM returned an empty response after gland measurement had completed; the executor ran the planned lymph-node screening, and the LLM then completed the run. Because the fallback can only run tasks listed in the plan, it cannot restore a step that the plan itself omits (see the omitted fusion step below).
\item \textbf{Rollback.} The current system does not implement transactional rollback of tool actions.
\end{itemize}

\textbf{Major failure modes and representative traces} (\responsetableref{tab_rep_failure_modes}; \responsefigureref{fig_rep_execution_traces}). Apart from the request errors above, all of which were contained, the main failure modes were ordering errors and omitted steps that the run-level success flags did not reveal:\par

\begin{itemize}[leftmargin=1.5em,itemsep=1pt,topsep=2pt]
\item In 40 of the 320 fusion-tool calls (24 cases; in all three runs of 3 of them), fusion was executed before pixel spacing was available for the nodule-positive images, and spacing obtained later did not trigger re-fusion. In 4 SMU-HMC and ZJH-TS runs, this left 5 nodule descriptions without a size. The other 36 calls were in KMVE cases, whose reports use a mask in place of specific nodule sizes. In addition, in one run, fusion was executed before nodule feature classification.
\item For the 43 cases in which preprocessing flagged no nodule, the workflow follows a plan that does not include spacing estimation. Consequently, gland measurement in these cases was performed without a pixel-spacing estimate.
\item In one run, a segmented nodule was never fused because fusion was missing from the plan, and the report contained no description of that nodule.
\end{itemize}

\responsefigureref{fig_rep_execution_traces} shows real traces for a case without nodules, a case with multiple nodules, conflicting intermediate evidence (a preprocessing-flagged image not confirmed by segmentation and excluded downstream, seen in all three runs of two cases), recovery from a malformed tool name, off-target images that were not processed, fusion before pixel spacing, and the omitted fusion step.\par

\end{authorresponse}
% Converted from S5_Tool_call_audit/reviewer_response_package_20261119_claude_v13/tables/Table_R2_16_1_failure_modes.md (2026-11-17)
\begin{responsetable}
\caption{\responsetabletitle{tab_rep_failure_modes}{Failure events in the audit (150 cases $\times$ 3 runs = 450 runs; 2,347 tool calls)}}
{\footnotesize\setlength{\tabcolsep}{6pt}
\begin{tabular}{@{}>{\raggedright\arraybackslash}p{0.43\textwidth}ccc@{}}
\toprule
\textbf{Event} & \textbf{Count} & \textbf{\% (95\% CI)} & \textbf{Cases (in 1 / 2 / 3 runs)} \\
\midrule
\rowcolor{ResponseTableGray}\multicolumn{4}{c}{\textbf{Failure categories}} \\
Invalid call & 4/2,347 calls & 0.17 (0.04--0.35) & 4 (4 / 0 / 0) \\
Lymph-node screening request with non-neck images & 18/2,347 calls & 0.77 (0.27--1.37) & 8 (2 / 2 / 4) \\
Gland-measurement request with neck images & 66/2,347 calls & 2.81 (1.96--3.70) & 32 (9 / 12 / 11) \\
Hallucinated task (planned without image support) & 3/2,348 planned tasks & 0.13 (0.00--0.29) & 3 (3 / 0 / 0) \\
Repeated call & 0/2,347 calls & 0 & 0 \\
Execution failure (malformed tool name, retried) & 1/2,347 calls & --- & 1 (1 / 0 / 0) \\
Recovery by a plan-guided supplementary call & 2/450 runs & --- & 2 (2 / 0 / 0) \\
Rollback & not implemented & --- & --- \\
\rowcolor{ResponseTableGray}\multicolumn{4}{c}{\textbf{Ordering and omission events not shown by the run-level success flags}} \\
Fusion before pixel spacing & 40/320 fusion calls & --- & 24 (11 / 10 / 3) \\
Fusion before feature classification & 1/320 fusion calls & --- & 1 (1 / 0 / 0) \\
Gland measurement without any pixel-spacing estimate & 129/450 runs & 28.7 (22.0--35.3) & 43 (0 / 0 / 43) \\
Segmented nodule never fused & 1/450 runs & --- & 1 (1 / 0 / 0) \\
Preprocessing-flagged image not confirmed by segmentation & 6/450 runs & --- & 2 (0 / 0 / 2) \\
\bottomrule
\end{tabular}}
\responsetablenote{Cases: cases affected among the 150, and the number in which the event occurred in one, two or all three runs (each run with a different seed). 95\% CIs: percentile bootstrap resampling cases within each dataset (10,000 resamples). Fusion with either ordering error: 41/320 calls, 12.8\% (8.8--17.1\%).}
\end{responsetable}

% Converted from S5_Tool_call_audit/reviewer_response_package_20261119_claude_v13/figures/FIGURE_CAPTION_EN.md and S5_Tool_call_audit/reviewer_response_package_20261119_claude_v13/figures/Fig_R2_16_traces.pdf (2026-11-17)
\begin{responsefigure}
\includegraphics[width=\textwidth,height=0.64\textheight,keepaspectratio]{report_rep_execution_traces.pdf}
\caption{\responsefiguretitle{fig_rep_execution_traces}{Representative execution traces of the planner--executor workflow. Each row shows one run from the audit (150 cases $\times$ 3 runs). Boxes are tool calls in the recorded execution order, given by the small numbers; arrows indicate order only, not data flow. The text under each box summarises the recorded arguments and returns, and the left column gives the dataset, the number of input images and the observed outcome. Fill colour and symbol together encode the execution status (legend). \textbf{a}, No thyroid nodule. No image was flagged by preprocessing, so segmentation, feature classification and fusion were not scheduled. Lymph-node screening of the four neck images returned one positive image, which the report reflects. \textbf{b}, Multiple nodules. Two preprocessing-flagged images were confirmed by segmentation, classified, and fused into two nodules, each reported with a size. \textbf{c}, Conflicting intermediate evidence. Of three preprocessing-flagged images, two were confirmed by segmentation; the third was not passed to classification or fusion. The same disagreement occurred in all three runs of this case and of one other case. \textbf{d}, Invalid tool name and recovery. The LLM emitted a malformed tool name and the executor returned an unknown-tool error; the next call used the registered name with identical arguments. \textbf{e}, Off-region targets not processed. A lymph-node screening request listed five images whose preprocessing region was not the neck. The executor passed only the neck images to the screening tool; the other five were not processed in this call. \textbf{f}, Dependency failure. Fusion ran before pixel-spacing estimates were available, so the fused nodules had no millimetre sizes. Spacing was computed later, but fusion was not repeated, and the report contains two nodule descriptions without sizes. \textbf{g}, Omission. A nodule was confirmed by segmentation and classified, but fusion was neither planned nor executed (dashed box), and the report contains no nodule description. Plan-guided recovery could not add the step because it was absent from the plan. Frequencies of each failure mode are given in \responsetableref{tab_rep_failure_modes}. Case identifiers, file paths and report text are omitted. LN, lymph node.}}
\label{fig_rep_execution_traces}
\end{responsefigure}



\reviewersection{Reviewer 2, Comment 17 (R2.17)}\phantomsection\label{rev:R2.17}
\begin{reviewercomment}
Agent workflow lacks concurrency control and conflict resolution mechanisms: When lymph-node screening and nodule analysis execute concurrently, how is duplication avoided if both detect the same hypoechoic structure in the lateral neck region?
\end{reviewercomment}
\begin{authorresponse}
% Converted from S6_Concurrency_conflicts/reviewer_response_package_20261116_claude_v3/RESPONSE_R2_17_EN.md (2026-11-17)
We thank the reviewer for this question. In the report-generation workflow of ThyroidXAgent, lymph-node screening and nodule analysis do not run concurrently, and the workflow does not match objects between the two branches. Duplication is avoided by anatomical region: lymph-node screening analyses only neck-region images, and nodules assigned to the neck region are not described as thyroid nodules in the report. We describe this mechanism below and quantified it in the 450 audit runs described in our response to Comment 16.\par

\textbf{Execution order.} Execution can be described at two levels, between tools and within a tool. Between tools, the only parallel step is preprocessing, in which three per-image classifiers (colour-Doppler detection, anatomical region classification and nodule-presence detection) run concurrently, each writing a separate field of the image record. Within a tool, the per-image tools, such as nodule segmentation, feature classification, pixel-spacing prediction and lymph-node screening, take a batch of images in one call and return one result per image. After preprocessing, the tools for different tasks, including nodule segmentation and lymph-node screening, are called one at a time by the executor. In the 450 audit runs, none of the 2,347 tool calls overlapped in time with another.\par

\textbf{Separation by region.} Preprocessing assigns each image to one region: a longitudinal or transverse view of the left or right lobe, a transverse view of both lobes or the isthmus, or the neck. By the definitions of the tools in the toolset, lymph-node screening accepts only neck images, and nodule segmentation analyses the images flagged by the nodule-presence detector, which can include neck images. The fusion tool groups nodules by lobe, so nodules in neck images are never merged with nodules in thyroid images. When the report is assembled, nodules in the neck region are not described as thyroid nodules.\par

\textbf{What happened in the audit runs.} The nodule-presence detector flagged 72 neck images in 27 of the 450 runs (6.0\%; 95\% CI 2.7--10.0\%; 9 cases), and segmentation detected a structure in 69 of them. In 60 of these images (24 runs, 5.3\%; 95\% CI 2.0--9.3\%; 8 cases), lymph-node screening was also positive, that is, both branches detected a structure in the same image. All 69 nodules segmented in neck images were fused as neck-region nodules, none was combined with a nodule from a thyroid image, and none appeared in a nodule description. They also did not change the TI-RADS category stated in the impression (24 runs; KMVE reports do not state a category).\par

The same separation held in the 1,042 test reports of the original study. There, 57 images in 35 reports were processed by both branches, and none of the 57 nodules fused from these images was described as a thyroid nodule.\par

\end{authorresponse}

\reviewersection{Reviewer 2, Comment 18 (R2.18)}\phantomsection\label{rev:R2.18}
\begin{reviewercomment}
The ThyClinScore weights lambda and eta directly affect report-quality assessment, but their selection, optimization dataset, and sensitivity analysis are not reported, raising concerns about overfitting. Although Figure 4c shows correlation with an LLM judge, validation against radiologist ratings and inter-rater reliability analysis are lacking.
\end{reviewercomment}
\begin{authorresponse}
% Converted from S1_ThyClinScore_sensitivity/reviewer_response_package_20261120_claude_v24/RESPONSE_R2_18_EN.md (2026-11-17)
We thank the reviewer. $\lambda$ and $\eta$ were set heuristically when the metric was designed ($\lambda$ = 0.3, $\eta$ = 0.5) and were not optimised on any dataset; the same applies to the other parameters, whose default values are now given in the Methods. We now report how the results depend on these parameters and calibrate $\lambda$ and $\eta$ against clinician ratings.\par

\textbf{Sensitivity analysis.} $\lambda$ was varied from 0.1 to 0.5 and $\eta$ from 0.3 to 0.7, singly and in all 25 combinations, as part of 120 configurations that also varied the lesion-matching threshold, the size tolerances and all component weights (our response to Reviewer 1, Comment \hyperref[rev:R1.5]{5}). ThyroidXAgent had the highest mean ThyClinScore in all 25 $\lambda$--$\eta$ combinations in SMU-HMC and KMVE and in 24 of 25 in ZJH-TS; in the exception ($\lambda$ = 0.1, $\eta$ = 0.7), its differences from GPT-5 and Qwen3.5 Plus had 95\% CIs that included zero. When $\lambda$ or $\eta$ was set to 0 or 1, ThyroidXAgent was not first in 2 of the 12 dataset--setting combinations, both with $\eta$ = 1, and ranked third in both: in ZJH-TS, its differences from GPT-5 and Qwen3.5 Plus had 95\% CIs that included zero, and in KMVE, LLaVA-Med, which reported no abnormality in any KMVE case, ranked first (difference $-$0.039; 95\% CI $-$0.056 to $-$0.020).\par

\textbf{Calibration and validation against clinician ratings.} To test whether the heuristic values are appropriate, we also selected $\lambda$ and $\eta$ on the clinician ratings of SMU-HMC (our response to Reviewer 1, Comment \hyperref[rev:R1.5]{5}). Among the 25 $\lambda$--$\eta$ combinations of the sensitivity analysis, $\lambda$ = 0.4 and $\eta$ = 0.7 had the highest Spearman correlation with both the clinical-accuracy rating and the mean of the five ratings, but the correlation with the clinical-accuracy rating varied little across the combinations (0.423--0.459; 0.437 with the default values). With these parameters fixed, the correlation in ZJH-TS and KMVE was 0.510 and 0.562, compared with 0.464 and 0.560 with the default parameters, and ThyroidXAgent had the highest mean ThyClinScore in all three test sets with either setting (Supplementary Table~\SuppTabClinVal{}; Methods~\SecMethThyClin{}). The default parameters and the parameters that best match the clinicians therefore lead to the same conclusions. We report the results with the default parameters, which were set before any clinician ratings were collected, and give those with the calibrated parameters in Supplementary Table~\SuppTabClinVal{}. Inter-rater agreement is reported in our response to Reviewer 1, Comment \hyperref[rev:R1.5]{5}.\par

\end{authorresponse}

\reviewersection{Reviewer 2, Comment 19 (R2.19)}\phantomsection\label{rev:R2.19}
\begin{reviewercomment}
The term "training-free" should be qualified. Although BM25 requires no parameter training, the workflow still depends on corpus selection, report decomposition, template construction, slot definitions, and clause-assembly rules. The authors should describe the extent of manual curation, template deduplication, slot-filling rules, and adaptation to centre-specific reporting styles. Clinician assessment of readability, factual consistency, and clinical usability is also needed beyond lexical-overlap metrics such as BLEU and ROUGE.
\end{reviewercomment}
\begin{authorresponse}
% Converted from S4_Template_construction/reviewer_response_package_20261119_claude_v13/RESPONSE_R2_19_EN.md (2026-11-17)
We thank the reviewer and agree. We no longer describe report generation as training-free. The report-realisation stage requires no additional model training or gradient-based fine-tuning, but it depends on template libraries and rules built from training reports. We now describe how these were built and used in the Methods.\par

\textbf{Corpus selection.} For SMU-HMC, 20,400 reports were scored for completeness, and the 400 reports of the test patients were excluded before tokenisation. Of the remaining 20,000 reports, the 7,147 with a completeness score of at least 0.9 were used to build the library. The completeness score is the report-completeness component of ThyClinScore (Methods), computed on the structured version of each report: the weighted fraction of the required information on thyroid measurements, parenchyma, lesions, impression and lymph nodes that is present. The KMVE library was built from the official KMVE training split (1,719 reports); in the revision, the 406 training reports that were near-duplicates of KMVE test reports were excluded, leaving 1,313 reports (our response to Reviewer 1, Comment \hyperref[rev:R1.5]{5}). ZJH-TS reports were generated with the SMU-HMC library.\par

\textbf{Decomposition and de-duplication.} The SMU-HMC reports were divided into five clause categories (thyroid measurement, gland morphology, nodule findings, lymph-node findings and impression). In each clause, case-specific content, such as measurements, locations, echogenicity, vascularity and TI-RADS categories, was replaced by slot placeholders, so that each clause became a candidate template. Identical candidate templates were merged, and the number of clauses mapped to each template was recorded as its frequency. Infrequent templates were then removed: gland and measurement templates outside the 400 and 40 most frequent, which account for 81\% and 99\% of the gland and measurement clauses, respectively; lymph-node templates occurring fewer than 20 times; nodule templates occurring fewer than 3 times; and impression clauses occurring only once. The construction scripts also removed 29 nodule templates whose fixed echogenicity wording conflicted with the echogenicity slot and one measurement template with an incomplete value, and manual review removed 57 gland and measurement templates (see below). In this way, the 1,793 distinct gland templates were reduced to 358, and the 7,478 nodule clauses, which formed 4,656 distinct templates, were reduced to 204 nodule templates. The final SMU-HMC library has 4,774 entries: 400 general sentences (358 gland, 24 measurement and 18 lymph-node sentences), 204 nodule templates, 1,076 impression clauses and 3,094 whole impressions. The KMVE library was rebuilt in the same way during the revision (our response to Reviewer 1, Comment \hyperref[rev:R1.5]{5}): the nodule sentences of these 1,313 training reports were converted into slot templates and merged with their frequencies, giving 82 nodule templates, and the 259 gland sentences were grouped by the gland size, parenchymal echo and parenchymal vascularity they state.\par

\textbf{Manual curation.} Manual review was applied to the 400 most frequent gland templates and the 40 most frequent measurement templates, and removed templates that are not typical of complete routine reports; all removed templates are listed in Appendix A3 (\responsetableref{tab_rep_excluded_templates}). Of the gland templates, 42 were removed: 9 describing postoperative or residual glands and 33 that do not form a complete gland description (9 measurement sentences, 13 describing only glandular vascularity and 11 describing focal lesions). Of the measurement templates, 15 were removed: 7 with a note on an emergency or bedside examination, 4 from postoperative examinations, 3 with missing values and 1 noting that the measurement was unsatisfactory. All other steps were carried out by the construction scripts with the thresholds given above.\par

\textbf{Slots and filling rules.} The SMU-HMC templates contain slots for nodule location, size, composition, echogenicity, margin, shape, aspect ratio, internal echo, echogenic foci and vascularity, and for the TI-RADS category in the impression. A query is built from the structured tool outputs, the best-matching template is retrieved by term matching, and its slots are filled with the corresponding values. Nodules are described by lobe; when both lobes contain nodules, a bilateral template is tried first. The impression is assembled from impression clauses retrieved with the nodule sentences as queries, penalising clauses whose cystic content or calcification disagrees with the nodule sentence, and the TI-RADS category is then filled in. If no clause is retrieved, the most similar whole impression is used; if neither is available, a rule-based impression states the TI-RADS category of each lobe or that no focal lesion was seen. The lymph-node sentence is a fixed sentence reporting enlarged lymph nodes when screening is positive and is otherwise retrieved from the library. When a tool gave no value, the generator used for the reports in the manuscript fell back on defaults: typical nodule features (for example, hypoechoic with a smooth margin), vascularity according to whether a Doppler image was available, and a random value within the typical range for a missing gland dimension; measured gland dimensions were also truncated to a fixed range. Note that our safety rules remove these fallbacks and instead mark the item as not assessed for the physician to complete (our response to Reviewer 1, Comment \hyperref[rev:R1.6]{6}). The KMVE templates contain slots for nodule location, number, echogenicity, margin, shape, echogenic foci and vascularity, which are filled with the same rules; nodule sizes are written as the masked size token of the KMVE reports, and the gland sentence is the most frequent KMVE sentence for the gland size, parenchymal echo and parenchymal vascularity stated by the gland-description tool.\par

\textbf{Centre-specific reporting styles.} Each library reproduces the wording of its source reports, and libraries can be exchanged without changing the image-analysis tools or the evidence schema. SMU-HMC and ZJH-TS reports were generated with the SMU-HMC library, and KMVE reports with the KMVE library, which follows the findings-only format of that dataset. A library for a new centre is built from its reports with the same procedure and added as a separate profile, which is selected when a case is started (Supplementary Fig.~\SuppFigTemplateBank{}). The dependence of the results on the template source is reported in our response to Reviewer 1, Comment \hyperref[rev:R1.5]{5}.\par

\textbf{Clinician assessment.} Two clinicians independently rated, blinded to the system, 270 reports of 90 cases (30 per test set; the ThyroidXAgent report and two baseline reports per case, with the baselines balanced across cases) against the reference report on five dimensions (1--5). Averaged over the 90 cases, the ThyroidXAgent report scored higher than the baseline reports of the same case in factual consistency (2.94 vs 2.46; difference 0.48, 95\% CI 0.24--0.72), readability (4.52 vs 4.08; 0.44, 0.28--0.60) and clinical usability (2.88 vs 2.41; 0.47, 0.22--0.71). By test set, the differences were clear on SMU-HMC (0.68 to 0.79) and ZJH-TS (0.40 to 0.57), mainly because in cases without nodules in the reference report, the baseline reports often described nodules that were not in the reference, whereas the ThyroidXAgent reports rarely did (for factual consistency, +1.43 in the 23 cases without nodules in the reference report, +0.28 in the 33 with one nodule and +0.03 in the 34 with two or more). On KMVE, the differences were small (0.09 to 0.22). We added these results to Results~\SecResReport{}. Agreement between the two clinicians is reported in our response to Reviewer 1, Comment \hyperref[rev:R1.5]{5}; it was low for readability.\par

\end{authorresponse}

\reviewersection{Reviewer 2, Comment 20 (R2.20)}\phantomsection\label{rev:R2.20}
\begin{reviewercomment}
Thresholds for connected-component analysis in post-processing are unspecified. The manuscript mentions removing isolated noise regions via connected-component analysis but does not specify: (a) the minimum connected-component area threshold; (b) whether small nodules (< 5 mm) may be erroneously removed; (c) the algorithm for separating multiple nodules. These details directly affect segmentation quality but are entirely absent from the Methods section.
\end{reviewercomment}
\begin{authorresponse}
We agree and clarify connected-component post-processing: the full rules are in this response below; Appendix~A2 §4 covers only how radiomics input differs from expert-mask connected-component steps (A2 focuses on PyRadiomics/AutoGluon, not a full segmentation post-processing spec). Reported segmentation Dice/HD95 and PyRadiomics use the same router-final nodule mask; we do not apply an additional connected-component refinement before radiomics.\par
Inside DINO segmentation experts, probability maps are binarized at 0.5; connected components smaller than 50 pixels are removed and a 5$\times$5 morphological closing is applied. For probability-ensemble masks, if multiple connected components remain, only the largest is retained. Each image is treated as a single ROI; we do not separately segment multiple nodules on the same image (reviewer (c): no per-lesion multi-nodule separation algorithm).\par
For (b) small nodules: the $\geq$50 px rule is not applied again on the router-final mask. On the GT equivalent-diameter $<$50 px subset (\responsetableref{tab_small_nodule_fail}; the ``$<$5.0 mm'' column is an illustrative 0.1 mm/px conversion), pipeline failures $M$ (empty final mask or Dice $<$0.1 vs.\ GT) are rare on most cohorts (e.g., DDTI 1/54, TN3K 2/68), with TN5K and ZJH-8K at 12/103 and 40/831; mean Dice on non-failed small-nodule cases is mostly 0.70--0.81 (ThyroidXL has no $<$50 px GT nodules).\par
\par
\end{authorresponse}
\begin{responsetable}
\caption{\responsetabletitle{tab_small_nodule_fail}{Nodules <5 mm --- Pipeline failure statistics}}
\begin{adjustbox}{max width=\textwidth,center}
\begin{tabular}{lccccc}
\toprule
Cohort & $n$ matched & <5.0 mm ($N$) & Failures ($M$) & Mean Dice (small-nodule subset) & GT equiv.\ diameter <50 px \\
\midrule
DDTI & 637 & 54 & 1 & 0.734 & 54 \\
PKTN & 150 & 15 & 1 & 0.741 & 15 \\
RJH-7K & 7288 & 366 & 1 & 0.766 & 366 \\
ThyroidXL & 2090 & 0 & 0 & --- & 0 \\
TN3K & 614 & 68 & 2 & 0.771 & 68 \\
TN5K & 1000 & 103 & 12 & 0.696 & 103 \\
ZJH-8K & 6807 & 831 & 40 & 0.808 & 831 \\
\bottomrule
\end{tabular}
\end{adjustbox}
\responsetablenote{Primary subset: GT equivalent diameter $<$50 px (same scale as the expert $\geq$50 px minimum-component rule); the ``$<$5.0 mm'' column is an illustrative mapping at 0.1 mm/px to relate to the reviewer's wording (public external cohorts lack reliable per-image spacing). Failures ($M$): empty router-final mask or Dice $<$0.1 vs.\ GT (full segmentation-and-routing pipeline, not radiomics-only connected-component removal). ZJH-8K matches \responsetableref{tab_dhash}: $n{=}6807$ after excluding 1151 ZJH images with dHash H=0 to TN3K train.}
\end{responsetable}


\reviewersection{Reviewer 2, Comment 21 (R2.21)}\phantomsection\label{rev:R2.21}
\begin{reviewercomment}
The AutoGluon-based radiomics branch is insufficiently described.
\end{reviewercomment}
\begin{authorresponse}
We thank the reviewer for this comment. Reproducibility details are in Appendix A2. At inference, two-dimensional PyRadiomics is extracted on the router-final nodule mask with no additional connected-component refinement before radiomics (the same mask underlies reported Dice/HD95; expert-internal $\geq$50 px rules are under \hyperref[rev:R2.20]{R2.20}). The AutoGluon tabular model was trained on the classification stacked dataset (\responsetableref{tab_baseline_protocol}) with reference masks; external tests use the frozen model only (no retuning on test). When independent classifiers disagree, the agent may combine AutoGluon malignancy probability with other signals; segmentation LLM prompts may include per-candidate radiomics-judge summaries (\hyperref[rev:R2.9]{R2.9} and Appendix A1).\par
\par
\end{authorresponse}

\reviewersection{Reviewer 2, Comment 22 (R2.22)}\phantomsection\label{rev:R2.22}
\begin{reviewercomment}
Please check the manuscript for grammatical errors, such as "ThyroidXAgent improved produced."
\end{reviewercomment}
\begin{authorresponse}
% Converted from S8_Numbers_and_language/reviewer_response_package_20261120_claude_v8/RESPONSE_R2_22_report_EN.md (2026-11-17)
We thank the reviewer. The sentence now reads ``ThyroidXAgent produced more clinically consistent reports''. We also corrected ``evidence-centerd'' to ``evidence-centred'', the heading ``Confilts'' to ``Competing interests'' and the statement ``The authors has no confilt of interests'' to ``The authors have no conflicts of interest'', and the figure labels ``Static Pipline'' (Fig.~\FigReportGen{}g), ``BlEU-4'' (Fig.~\FigThyClin{}) and ``Correlation with LLM a Judge score'' (Fig.~\FigThyClin{}c), and ``Lymph Node Metastatis'' (Fig. 3b).\par

\end{authorresponse}

\reviewersection{Reviewer 2, Comment 23 (R2.23)}\phantomsection\label{rev:R2.23}
\begin{reviewercomment}
Please specify whether TI-RADS refers to ACR TI-RADS, C-TIRADS, or another system, as categories such as "TI-RADS 4b" are system-specific.
\end{reviewercomment}
\begin{authorresponse}
% Converted from S8_Numbers_and_language/reviewer_response_package_20261120_claude_v8/RESPONSE_R2_23_EN.md (2026-11-17)
We thank the reviewer. ThyroidXAgent reports TI-RADS categories on the five levels of the ACR TI-RADS (TR1--TR5). The category is predicted from the nodule image by the TI-RADS branch of the multitask nodule model, which was trained on the ACR TI-RADS levels of the Stanford thyroid ultrasound cine-clip dataset and on DDTI. DDTI uses an earlier TI-RADS with subcategories 4a--4c, which we merged into level 4, so ThyroidXAgent does not output subcategories. The reference reports of SMU-HMC and ZJH-TS were written in routine practice at the two hospitals, which use the Chinese TI-RADS (C-TIRADS); most give a single category from 1 to 5, and 7 of the 400 SMU-HMC and 10 of the 150 ZJH-TS reports give subcategories 4a--4c. Our analyses do not require the two systems to agree: ThyClinScore checks only whether the impression states a TI-RADS category and does not compare categories, and in the reader study categories 4--5 were taken to indicate a malignant or suspicious nodule, which corresponds to the suspicious levels of both systems. The ``TI-RADS 4b'' in Supplementary Fig.~\SuppFigCaseReview{} was written by the Qwen3.5 Plus baseline, not by ThyroidXAgent.\par

\end{authorresponse}

\responsetoreviewer{Response to Reviewer 3}
\reviewersection{Reviewer 3, Comment 1 (R3.1)}\phantomsection\label{rev:R3.1}
\begin{reviewercomment}
The proposed system integrates a heterogeneous ensemble of expert models, radiomic features, metadata, and an LLM-based router. It is therefore unclear to what extent the performance obtained comes from the agent itself rather than from model ensembling or the use of stronger component models. Including controlled baseline comparisons with majority voting, probability averaging, rule-based routing, and a non-LLM learned router based on the same expert pool would better clarify the unique advantage provided by the agent.
\end{reviewercomment}
\begin{authorresponse}
We agree that gains must be separated from the fixed expert pool versus routing and fusion choices. On the same fixed expert pool and frozen candidate outputs as Reviewer 1, Comment \hyperref[rev:R1.2]{2} (\hyperref[rev:R1.2]{R1.2}), we tabulate majority vote, soft voting, and static agreement-based routing together with validation best single expert, max-confidence routing, non-LLM learned routing (6), deployable full-system LLM routing (7), and the non-deployable oracle (8) in \responsetablerefrange{tab_routing_ablation_seg}{tab_routing_ablation_cls}, so comparisons are not limited to weak heuristics or uncontrolled ensembling.\par
Strategy (6) is the requested non-LLM learned router: separate multinomial logistic routers for segmentation and classification ($C{=}0.1$, seed 42) on per-candidate quality and disagreement summaries (no device/data-source metadata; no LLM), with validation-defined supervision (segmentation: highest mask Dice vs.\ GT; classification: among correct experts when top classes disagree, highest confidence), then frozen weights on test for discrete expert selection. Strategy (7) is deployable ThyroidXAgent LLM routing under the same external-test protocol as \responsetableref{tab_baseline_protocol}, \responsetableref{tab_cls_ext}, and \responsetableref{tab_seg_ext}; it is reported in the same tables and protocol as static strategies (1)--(5).\par
Those tables show that no single static strategy dominates every cohort and task, and that routing via (6) or the full agent (7) versus static baselines varies by cohort, so performance should not be attributed to the expert pool or fixed averaging alone. In addition to TN3K, TN5K, and ThyroidXL, ZJH-8K, RJH-7K (segmentation only), DDTI, and PKTN (segmentation only) are included under the same eight-strategy protocol; point estimates and 95\% CIs follow Reviewer 2, Comment \hyperref[rev:R2.2]{2} (\hyperref[rev:R2.2]{R2.2}).\par
\par
\end{authorresponse}

\reviewersection{Reviewer 3, Comment 2 (R3.2)}\phantomsection\label{rev:R3.2}
\begin{reviewercomment}
The LLM router leverages metadata such as device and data source in addition to prediction quality signals. This may allow the system to learn center-specific expert preferences rather than case-level routing. Please provide ablation analysis to evaluate the effect from data-source, scanner metadata and report expert-selection patterns across centers.
\end{reviewercomment}
\begin{authorresponse}
We thank the reviewer. We agree that routing must be tested for over-reliance on data-source or scanner metadata versus case-level quality and disagreement signals. On the same fixed expert pool and frozen upstream candidate outputs as the full-system LLM router (strategy (7)) in \responsetablerefrange{tab_routing_ablation_seg}{tab_routing_ablation_cls}, we vary only metadata visible in the LLM router prompts (\responsetablerefrange{tab_metadata_ablation_seg}{tab_metadata_ablation_cls}: full metadata; data source withheld; device withheld; both withheld; quality and disagreement signals only; primary contrast full $\rightarrow$ quality/disagreement only). Strategy-6 learned routing (Reviewer 3, Comment \hyperref[rev:R3.1]{1}) does not use device/source metadata and is reported in \responsetablerefrange{tab_routing_ablation_seg}{tab_routing_ablation_cls}.\par
The router is not fine-tuned on hospital or centre IDs. On BM external cohorts (TN3K, TN5K, ThyroidXL, ZJH-8K, DDTI, etc.), device and data-source fields are set per dataset and are constant within each cohort's test set; we do not inject per-image centre or scanner-instance labels. Centre- and source-related effects in this work therefore operate at cohort-level metadata.\par
Segmentation and classification results under the metadata ablation are in \responsetablerefrange{tab_metadata_ablation_seg}{tab_metadata_ablation_cls}; per-cohort LLM expert-selection shares are in \responsetableref{tab_expert_selection_seg} and \responsetableref{tab_expert_selection_cls}.\par
\par
\end{authorresponse}
% LLM metadata ablation (tab_metadata_ablation_seg / tab_metadata_ablation_cls); TN5K
% [SYNC] Full metadata row = tab_routing_ablation_* TN5K strategy (7).
% [DATA-SOURCE] tab_metadata_ablation_seg / tab_metadata_ablation_cls | Real | B5 output/ablation/b5/tn5k/ (other four tiers)

\begin{responsetable}
\caption{\responsetabletitle{tab_metadata_ablation_seg}{LLM router metadata ablation --- nodule segmentation}}
\begin{adjustbox}{max width=\textwidth,center}
\begin{tabular}{@{}l l l@{}}
\toprule
Metadata condition & Dice $\uparrow$ & HD95 $\downarrow$ \\
\midrule
Device withheld & 0.8242 (0.8020--0.8446) & 7.8131 (5.9222--9.8453) \\
Data source withheld & 0.8179 (0.7950--0.8365) & 7.8435 (6.0382--9.7338) \\
Both withheld & 0.8179 (0.7945--0.8380) & 7.4440 (5.7012--9.2961) \\
Quality/disagreement only & 0.8228 (0.8010--0.8419) & 7.6001 (5.7902--9.5043) \\
Full metadata & \textbf{0.8388 (0.8260--0.8516)} & \textbf{7.1430 (5.9130--8.3730)} \\
\bottomrule
\end{tabular}
\end{adjustbox}
\responsetablenote{Fixed expert pool with TN5K external-test shared upstream expert outputs; segmentation uses LLM routing with only prompt-visible metadata changed (Appendix A1). Full-metadata row matches TN5K strategy (7) in \responsetableref{tab_routing_ablation_seg}. Device/source/both-withheld and quality/disagreement-only conditions use cluster-bootstrap 95\% CIs (1000 resamples, seed 42). Dice and HD95 use four decimal places (point estimates and CI bounds).}
\end{responsetable}

% LLM router metadata ablation — classification (label tab_metadata_ablation_cls)
\begin{responsetable}
\caption{\responsetabletitle{tab_metadata_ablation_cls}{LLM router metadata ablation --- malignancy classification}}
\begin{adjustbox}{max width=\textwidth,center}
\begin{tabular}{@{}l l l@{}}
\toprule
Metadata condition & AUROC $\uparrow$ & AUPRC $\uparrow$ \\
\midrule
Device withheld & 0.9403 (0.9048--0.9699) & 0.9733 (0.9549--0.9877) \\
Data source withheld & 0.9400 (0.9040--0.9690) & 0.9741 (0.9555--0.9878) \\
Both withheld & 0.9392 (0.9038--0.9679) & 0.9737 (0.9563--0.9873) \\
Quality/disagreement only & 0.9404 (0.9038--0.9700) & 0.9740 (0.9556--0.9878) \\
Full metadata & \textbf{0.9584 (0.9432--0.9736)} & \textbf{0.9763 (0.9674--0.9832)} \\
\bottomrule
\end{tabular}
\end{adjustbox}
\responsetablenote{Fixed expert pool with TN5K external-test shared upstream expert outputs; classification uses LLM routing with only prompt-visible metadata changed (Appendix A1). When all independent classification experts agree we use deterministic consensus without invoking the LLM; when they disagree, classification routing uses the LLM. Classification routing inputs when the LLM is invoked were not separately metadata-ablated (only prompt-visible fields varied).}
\end{responsetable}

% Source of truth: edit by hand
% [DATA-SOURCE] tab_expert_selection_seg
\begin{responsetable}
\caption{\responsetabletitle{tab_expert_selection_seg}{Per-case share of LLM-selected segmentation experts on each external test cohort}}
\begin{adjustbox}{max width=\textwidth,center}
\begin{tabular}{@{}l r r r r r@{}}
\toprule
Cohort & Multi-cohort joint & TN3K & TN5K & ThyroidXL & PKTN \\
\midrule
TN3K & 7.8\% & 9.4\% & 0.5\% & 66.4\% & 15.8\% \\
TN5K & 0.6\% & 5.7\% & 0.9\% & 83.9\% & 8.9\% \\
ThyroidXL & 40.0\% & 6.8\% & 0.6\% & 42.4\% & 10.2\% \\
ZJH-8K & 8.7\% & 9.8\% & 1.0\% & 65.5\% & 15.0\% \\
DDTI & 17.6\% & 9.1\% & 0.6\% & 59.7\% & 13.0\% \\
RJH-7K & 9.3\% & 8.6\% & 1.2\% & 71.6\% & 9.2\% \\
PKTN & 6.7\% & 13.3\% & 0.0\% & 68.7\% & 11.3\% \\
\bottomrule
\end{tabular}
\end{adjustbox}
\end{responsetable}

% Expert selection — classification (label tab_expert_selection_cls)
\begin{responsetable}
\caption{\responsetabletitle{tab_expert_selection_cls}{Per-case share of LLM-selected independent classifiers on each external test cohort with benign--malignant labels (five cohorts)}}
\begin{adjustbox}{max width=\textwidth,center}
\begin{tabular}{@{}l r r r r@{}}
\toprule
Cohort & Multi-cohort joint & TN3K & TN5K & ThyroidXL \\
\midrule
TN3K & 35.1\% & 16.5\% & 12.9\% & 35.5\% \\
TN5K & 56.4\% & 16.6\% & 11.1\% & 15.8\% \\
ThyroidXL & 55.0\% & 7.5\% & 30.0\% & 7.5\% \\
ZJH-8K & 74.6\% & 7.7\% & 5.4\% & 12.3\% \\
DDTI & 97.3\% & 0.0\% & 2.2\% & 0.4\% \\
\bottomrule
\end{tabular}
\end{adjustbox}
\end{responsetable}



\reviewersection{Reviewer 3, Comment 3 (R3.3)}\phantomsection\label{rev:R3.3}
\begin{reviewercomment}
The definition of "28,458 non-overlapping test cases" is unclear. Please clearly describe the data hierarchy and confirm that all train/validation/test splits are patient disjoint. Besides, a simple flowchart showing the relationship of patient, examination, image/video, lesion and task would greatly improve the clarity.
\end{reviewercomment}
\begin{authorresponse}
We agree that the patient / examination / image / task hierarchy must be explicit and that train/validation/test splits should be patient-disjoint where applicable. The revised abstract and Methods report image-benchmark test sizes without a single combined total (\hyperref[rev:R2.3]{R2.3}); stratified cohort roles and counts for the image-analysis arm are in \responsefigureref{fig_image_dataflow}.\par
\responsefigureref{fig_image_dataflow} summarizes the image-analysis arm: stacked public cohorts for segmentation/classification expert training, a pooled validation set (expert selection, learned routing, and abstention-threshold fitting; thresholds frozen for external testing), independent external test cohorts, and the ZJH malignant subset used only for LNM/PTC-subtype development (generic versus malignant-subtask external tests match Reviewer 1, Comment \hyperref[rev:R1.4]{4} (\hyperref[rev:R1.4]{R1.4})); cohort roles, splits, and counts are on the figure and in its caption.\par
Under this partitioning, public cohorts (TN3K, TN5K, ThyroidXL, PKTN, etc.) use their published train/val/test splits; cohorts designated external (e.g.\ DDTI, RJH-7K, ZJH-8K, NHC-MISD-TUS) were not used to train the stacked segmentation/classification experts. Where reproducible patient IDs exist (ZJH-8K, ThyroidXL), bootstrap and patient-level analyses use patient clusters (\hyperref[rev:R1.1]{R1.1}, \hyperref[rev:R2.2]{R2.2}); external cohorts without public IDs follow \hyperref[rev:R1.1]{R1.1} as one image per cluster (\responsetableref{tab_cls_ext}). dHash screening of each source train split against ZJH-8K is in \responsetableref{tab_dhash} (\hyperref[rev:R1.1]{R1.1}, \hyperref[rev:R2.4]{R2.4}).\par
That figure covers cohort--split--role assignments for segmentation/classification components. The patient--examination--image/video--lesion--task hierarchy is summarized in the paragraphs above and is distinct from the cohort-level split layout in \responsefigureref{fig_image_dataflow}.\par
\par
\end{authorresponse}

\reviewersection{Reviewer 3, Comment 4 (R3.4)}\phantomsection\label{rev:R3.4}
\begin{reviewercomment}
The reporting study includes only two physicians and 145 videos. The results are encouraging, but the conclusions regarding improved physician performance should be moderated. Please report confidence intervals, appropriate paired or mixed-effects statistics, and clearly distinguish annotation consistency from validated diagnostic accuracy.
\end{reviewercomment}
\begin{authorresponse}
% Converted from S10_Reader_study/reviewer_response_package_20261120_claude_v17/RESPONSE_R3_4_EN.md (2026-11-17)
We thank the reviewer and agree. We have moderated the conclusions and added the requested statistics.\par

\textbf{Confidence intervals and paired statistics.} We now report, for each physician and condition, the annotation consistency with 95\% CIs, the difference between conditions for each physician, the paired difference over the 145 cases with a case-bootstrap 95\% CI and McNemar's test, and a linear model with the condition and the physician as covariates (\responsefigureref{fig_rep_reader_by_physician}; our responses to Reviewer 1, Comment \hyperref[rev:R1.6]{6}, and Reviewer 2, Comment \hyperref[rev:R2.13]{13}). Results~\SecResReport{} now reports the annotation consistency by physician and the paired difference with its 95\% CI and McNemar's test, and Methods~\SecMethReader{} describes these analyses.\par

\textbf{Annotation consistency and diagnostic accuracy.} The outcome of the reporting study is consistency with the case labels derived from the video annotations, and we now describe it only as annotation consistency, not as diagnostic accuracy. In the safety study added in this revision (our response to Reviewer 1, Comment \hyperref[rev:R1.6]{6}), the reports were also compared with physician-confirmed findings: for each case, the reviewing physician confirmed the number of nodules and the benign or malignant tendency of each nodule in each lobe and the isthmus. These findings result from a physician's review of the video rather than from pathology, and we report the comparison as consistency with physician-confirmed findings. \responsefigureref{fig_rep_safety_study}b shows both references for the AI drafts, the manual and AI-assisted reports of the reader study and the final reports of the safety study in the same 70 cases.\par

\end{authorresponse}

\reviewersection{Reviewer 3, Comment 5 (R3.5)}\phantomsection\label{rev:R3.5}
\begin{reviewercomment}
ThyClinScore is mainly validated by correlation with an LLM judge. This is not sufficient as a clinical gold standard. To strengthen the validity of the metric, please validate the metric against blinded ratings from multiple radiologists and provide sensitivity analysis for the main weights and thresholds used. The manuscript should also acknowledge that ThyroidXAgent does not outperform all baselines across every semantic submetric.
\end{reviewercomment}
\begin{authorresponse}
% Converted from S1_ThyClinScore_sensitivity/reviewer_response_package_20261120_claude_v24/RESPONSE_R3_5_EN.md (2026-11-17)
We thank the reviewer and have addressed the three points.\par

\textbf{Validation against blinded ratings.} Two clinicians rated 270 reports blinded to the system. The Spearman correlation of ThyClinScore with their clinical-accuracy rating was 0.437 to 0.560 in the three test sets, with all 95\% CIs excluding zero. Inter-rater agreement (ICC(2,1)) was 0.54--0.79 for clinical accuracy, factual consistency and clinical usability, and lower for completeness and readability (our response to Reviewer 1, Comment \hyperref[rev:R1.5]{5}).\par

\textbf{Sensitivity analysis.} We recomputed ThyClinScore under 120 configurations of the weights and thresholds (our response to Reviewer 1, Comment \hyperref[rev:R1.5]{5}, and \responsefigureref{fig_rep_sensitivity_curves} and \responsefigureref{fig_rep_sensitivity_ranks}). ThyroidXAgent had the highest mean ThyClinScore in all 120 configurations in SMU-HMC, in 110 in KMVE and in 117 in ZJH-TS.\par

\textbf{Submetrics.} We now state in Results~\SecResReport{} that ThyroidXAgent does not lead on every semantic submetric. It had the highest ThyClinScore in all three test sets, but baselines led on the following submetrics (Supplementary Table~\SuppTabClinical{}). In SMU-HMC, LLaVA-Med had a lower FDR (0.1625 vs 0.2238). In KMVE (revised library), three baselines had a lower FDR (LLaVA-Med 0.0000, MedGemma 0.1877 and Gemini 2.5 Pro 0.3211 vs 0.3709); six had higher feature accuracy (GPT-5 0.6967, MedGemma 0.6820, GPT-4o 0.6608, Claude Sonnet 4.6 0.6605, Qwen3.5 Plus 0.6154 and Gemini 2.5 Pro 0.5093 vs 0.5087); Gemini 2.5 Pro had a higher lesion F1 (0.4648 vs 0.4223); five had higher completeness (Claude Sonnet 4.6 0.6707, MedGemma 0.6622, Qwen3.5 Plus 0.6525, LLaVA-Med 0.6508 and GPT-5 0.6506 vs 0.6329); and LLaVA-Med had higher consistency (0.5669 vs 0.5190). LLaVA-Med reported no abnormality in any KMVE case, which gives an FDR of 0. In ZJH-TS, LLaVA-Med had a lower FDR (0.2733 vs 0.2833), LLaVA-Med and the KMVE model had higher feature accuracy (0.6595 and 0.6288 vs 0.5820), and GPT-5 and Qwen3.5 Plus had higher consistency (0.4797 and 0.4767 vs 0.4564).\par

\end{authorresponse}

\reviewersection{Reviewer 3, Minor comment 1 (R3.M1)}\phantomsection\label{rev:R3.M1}
\begin{reviewercomment}
ThyroidXAgent receives structured outputs from multiple specialized tools, whereas several multimodal LLM baselines appear to work more directly from images. The manuscript should clarify this difference and avoid interpreting these comparisons as a pure test of language-generation ability.
\end{reviewercomment}
\begin{authorresponse}
% Converted from S7_Baseline_comparison/reviewer_response_package_20261119_claude_v6/RESPONSE_R3_m1_EN.md (2026-11-17)
We thank the reviewer and agree. The methods in the report-generation comparison received different inputs, so the comparison is between complete systems rather than a test of language generation. We now state this in the Results and describe the inputs of each method in the Methods.\par

\textbf{ThyroidXAgent.} ThyroidXAgent analyses all images of a case (median 13 images per case in SMU-HMC, 11.5 in ZJH-TS and 2 in KMVE) with specialised tools for segmentation, feature classification, measurement and lymph-node screening. Its language model plans and calls these tools and receives their structured outputs, but it does not write the report. The report is assembled from the structured findings by retrieving templates from a library built from training reports: the SMU-HMC library for SMU-HMC and ZJH-TS, and the KMVE library for KMVE.\par

\textbf{Baselines.} The five general-purpose multimodal models (GPT-4o, GPT-5, Gemini 2.5 Pro, Qwen3.5 Plus and Claude Sonnet 4.6) and MedGemma received a prompt that specified the report structure and gave three example reports in the reporting format of the dataset (few-shot prompting). They also received up to five images sampled at random from each case in SMU-HMC and ZJH-TS, or the two images of each case in KMVE. LLaVA-Med received two images and an English version of the prompt, and its output was translated into Chinese with GPT-4o-mini. The KMVE model, which was trained on the KMVE dataset to generate reports from images, received two images per case and was evaluated on SMU-HMC and ZJH-TS.\par

\textbf{What the comparisons measure.} The two comparisons address different questions. The comparison with the multimodal models and the KMVE model evaluates evidence-grounded report assembly as a whole: ThyroidXAgent derives structured evidence from all images with specialised tools and assembles the report from templates, whereas these models generate the report directly from a subset of the images (in KMVE, all methods received the same two images per case). Its result therefore reflects the benefit of grounding the report in tool-derived evidence, not the language-generation ability of the underlying model. The comparison with the static pipeline has matched inputs: the static pipeline has the same tools and templates as ThyroidXAgent, and the tools are called by fixed rules instead of being planned by the language model. On SMU-HMC and ZJH-TS, ThyroidXAgent had a higher ThyClinScore than the static pipeline (0.5189 vs 0.4129 and 0.4676 vs 0.3528; Fig.~\FigReportGen{}g and Supplementary Table~\SuppTabStatic{}), which reflects the contribution of adaptive tool routing. On KMVE, the two had the same ThyClinScore (0.4374), because KMVE reports focus on qualitative descriptions of the gland and nodules with measurements masked, so that a few tools suffice to generate them. We also corrected inconsistencies in the static-pipeline settings, and on KMVE the static pipeline now uses the revised KMVE library, as ThyroidXAgent does; Supplementary Table~\SuppTabStatic{} and Fig.~\FigReportGen{}g have been updated.\par

\end{authorresponse}
\responseappendixheading

\subsection*{Appendix A1. LLM router specification}
\label{app:A1}
\begin{responseappendixprose}
\noindent This appendix gives a plain-language flow summary first, then decoding parameters, English prompt templates matching current external-test deployment (compact JSON-only routing), required LLM output formats, and the contents of structured routing records. Flow summary: segmentation routing selects one mask among candidate experts; if \textbf{all} independent classifiers share the same top class, classification does not call the LLM; if they disagree, an LLM router selects one expert ($k{=}1$) or a multi-model fusion set ($k>1$, deployment-configured $k$); if independent classifiers still lack consensus, a further LLM arbitration step synthesizes independent models, AutoGluon on the selected mask, and the segmentation decision into a final benign/malignant call; \textbf{no free-text rationales are requested from the LLM}---only structured JSON fields; all steps are written to auditable structured records.
\par\noindent Verbatim blocks quote prompt wording; line breaks in this PDF are for readability only---the API receives program-assembled strings under Table A1-1 decoding settings.
\par
\begin{responseinlinetable}
\responseinlinetablecaption{Table A1-1. Shared LLM decoding parameters (segmentation router, classification router, and LLM arbitration when classifiers disagree)}
\label{tab_A1_decode}
\begin{adjustbox}{max width=\textwidth,center}
\begin{tabular}{ll}
\toprule
Parameter & Setting \\
\midrule
API type & OpenAI-compatible endpoint (DashScope) \\
Endpoint URL & \texttt{https://dashscope.aliyuncs.com/compatible-mode/v1} \\
Model & \texttt{qwen3.8-27b} \\
Temperature & 0.3 \\
Max output tokens & 768 \\
Extended thinking & Disabled \\
API key & Configured via environment variable; not written to routing records \\
\bottomrule
\end{tabular}
\end{adjustbox}
\end{responseinlinetable}

\noindent 1.~Segmentation LLM routing
\par
Runtime settings: LLM routing enabled; disagreement metrics included in the prompt; full device/source metadata (full-metadata condition in \responsetablerefrange{tab_metadata_ablation_seg}{tab_metadata_ablation_cls}); radiomics-judge summaries per candidate mask; single-mask selection per case.\par
System (base, always used):\par
\begin{promptverbatim}
You are a segmentation model selection agent. Choose the best mask from the JSON metrics in the user message.
Do not output chain-of-thought, Markdown, or any text outside JSON. The first character must be { and the last must be }.
Output one JSON line only: "selected_model", "confidence" in [0,1]; do not output a reasoning field. Example: {"selected_model":"ModelName", "confidence":0.9}
Priority: higher agreement_with_others; when experts disagree, lower mean_hd95_to_others; prefer radiomics_judge aligned with classification_anchor; morphology and device match secondary.
\end{promptverbatim}
Optional one-line system append at runtime: when all independent classifiers agree on the top class, append a classification-anchor line (e.g., anchor class and down-weight contradictory \texttt{radiomics\_judge}); otherwise append a no-anchor line prioritizing radiomics internal consistency and agreement.\par
User message: program-assembled JSON with device/source fields (metadata ablation may omit them), inter-candidate disagreement summaries, per-candidate quality and \texttt{radiomics\_judge} fields, and \texttt{classification\_anchor} (class string when unanimous, null otherwise).\par
Model output: one JSON line with \texttt{selected\_model} and \texttt{confidence} (no \texttt{reasoning} field).\par
Non-LLM fallbacks: if LLM routing is off or the API/parse fails, a deterministic rule-based static selector is used.\par
\noindent 2.~Classification --- standalone LLM router (no LLM when all independent classifiers share the same top class)
\par
When all independent classifiers share the same top class, deterministic consensus is used (highest top\_confidence) with no LLM call.\par
System (single-expert selection, $k{=}1$):\par
\begin{promptverbatim}
From the predictions JSON in the user message, select the most trustworthy model.
No chain-of-thought or Markdown. Output one JSON line: selected_model, selected_class, confidence in [0,1]. Do not output a reasoning field.
\end{promptverbatim}
System (multi-model probability fusion, $k>1$):\par
\begin{promptverbatim}
Select the k most trustworthy models (k set at runtime, descending trust) for class-probability averaging.
Output pure JSON only: selected_models (string array of length exactly k). Do not output reasoning, selected_class, or confidence.
\end{promptverbatim}
User message: device information (may be omitted under metadata ablation) and per-model prediction JSON (assembled by the program, including the routing $k$).\par
Model output: for $k{=}1$, \texttt{selected\_model}, \texttt{selected\_class}, and \texttt{confidence}; for $k>1$, a length-$k$ \texttt{selected\_models} array only (no \texttt{reasoning} in either mode).\par
\noindent 3.~Classification --- LLM arbitration when classifiers disagree
\par
System:\par
\begin{promptverbatim}
Thyroid ultrasound arbitration when independent classifiers did not reach consensus.
Use independent_models, autogluon_based_on_selected_mask, and segmentation_decision in the user JSON.
No chain-of-thought or Markdown. Output one JSON line only:
{"selected_class":"<benign/malignant label>", "confidence":0.0-1.0}
Do not output a reasoning field. If all signals are ambiguous, bias slightly toward malignant but lower confidence.
\end{promptverbatim}
User message: one JSON object with scenario text, \texttt{independent\_models} (all independent classifiers), AutoGluon on the selected mask (optional), and \texttt{segmentation\_decision} (selected segmentation expert name only---no LLM free-text rationale).\par
Model output: one JSON line with \texttt{selected\_class} and \texttt{confidence} (no \texttt{reasoning}); the routing record marks ``LLM arbitration'' and ``LLM invoked.''\par
\noindent 4.~Classification --- deterministic consensus (all independent classifiers share the same top class): no LLM; the routing record marks ``classification consensus.''
\par
\noindent 5.~Structured routing records (field definitions in Table A1-2; ten de-identified exemplars Rec-01--Rec-10 below)
\par
\begin{responseinlinetable}
\responseinlinetablecaption{Table A1-2. Information stored per structured routing record (de-identified archive)}
\label{tab_A1_log_fields}
\begin{adjustbox}{max width=\textwidth,center}
\begin{tabular}{ll}
\toprule
Item & Description \\
\midrule
Case ID & De-identified image/case identifier \\
Task type & Segmentation routing / classification routing / LLM arbitration when classifiers disagree \\
Candidate similarity tier & ``Similar'' or ``easy'' (quality gap between top two candidates $\leq 0.10$ $\Rightarrow$ similar; $>0.10$ $\Rightarrow$ easy) \\
LLM invoked & Yes / no \\
Decision mode & LLM routing, classification consensus, rule fallback, etc. \\
Selected segmentation expert & Multi-cohort joint vs.\ cohort-specific expert (segmentation) \\
Final benign/malignant class & Classification or arbitration tasks \\
Confidence & 0--1 \\
Decision note & Auditable note (not free-text LLM output); programmatic text when consensus applies; exemplars may summarize key candidate inputs \\
Decoding settings & Model name, temperature, max tokens (Table A1-1) \\
Timestamp & ISO 8601 \\
\bottomrule
\end{tabular}
\end{adjustbox}
\end{responseinlinetable}
\par
% [DATA-SOURCE] Rec-01..Rec-10 | TN3K routing_log.jsonl de-identified exemplars
\noindent\textbf{De-identified routing-record exemplars (Rec-01--Rec-10).}
\label{sec_A1_routing_exemplars}
\par\smallskip
\noindent\textbf{Rec-01.}\par
\begin{itemize}[leftmargin=1.5em,topsep=0.1em,itemsep=0.05em]
\item Case ID: Rec-01
\item Task type: Classification routing (LLM arbitration when classifiers disagree)
\item Candidate similarity tier: Similar
\item LLM invoked: Yes
\item Decision mode: LLM arbitration
\item Final benign/malignant class: Benign
\item Confidence: 0.720
\item Decision note: Independent classifiers: 4 independent classifiers: Malignant (0.66), Benign (0.78), Benign (0.99), Benign (1.00); AutoGluon (selected mask): Malignant (0.587)
\end{itemize}
\noindent\textbf{Rec-02.}\par
\begin{itemize}[leftmargin=1.5em,topsep=0.1em,itemsep=0.05em]
\item Case ID: Rec-02
\item Task type: Segmentation routing
\item Candidate similarity tier: Similar
\item LLM invoked: Yes
\item Decision mode: LLM mask selection
\item Selected segmentation expert: TN3K-specific
\item Confidence: 0.840
\item Decision note: Candidates: Multi-cohort joint (agreement 0.820); TN3K-specific [selected] (agreement 0.819); ThyroidXL-specific (agreement 0.733); PKTN-specific (agreement 0.772); TN5K-specific (agreement 0.664); Classification anchor: Benign
\end{itemize}
\noindent\textbf{Rec-03.}\par
\begin{itemize}[leftmargin=1.5em,topsep=0.1em,itemsep=0.05em]
\item Case ID: Rec-03
\item Task type: Classification routing (LLM arbitration when classifiers disagree)
\item Candidate similarity tier: Similar
\item LLM invoked: Yes
\item Decision mode: LLM arbitration
\item Final benign/malignant class: Malignant
\item Confidence: 0.650
\item Decision note: Independent classifiers: 4 independent classifiers: Malignant (0.64), Malignant (1.00), Benign (0.96), Benign (1.00); AutoGluon (selected mask): Malignant (0.766)
\end{itemize}
\noindent\textbf{Rec-04.}\par
\begin{itemize}[leftmargin=1.5em,topsep=0.1em,itemsep=0.05em]
\item Case ID: Rec-04
\item Task type: Segmentation routing
\item Candidate similarity tier: Similar
\item LLM invoked: Yes
\item Decision mode: LLM mask selection
\item Selected segmentation expert: TN3K-specific
\item Confidence: 0.680
\item Decision note: Candidates: Multi-cohort joint (agreement 0.355); TN3K-specific [selected] (agreement 0.369); ThyroidXL-specific (agreement 0.317); PKTN-specific (agreement 0.301); TN5K-specific (agreement 0.090)
\end{itemize}
\noindent\textbf{Rec-05.}\par
\begin{itemize}[leftmargin=1.5em,topsep=0.1em,itemsep=0.05em]
\item Case ID: Rec-05
\item Task type: Classification routing (LLM arbitration when classifiers disagree)
\item Candidate similarity tier: Similar
\item LLM invoked: Yes
\item Decision mode: LLM arbitration
\item Final benign/malignant class: Malignant
\item Confidence: 0.600
\item Decision note: Independent classifiers: 4 independent classifiers: Malignant (1.00), Malignant (1.00), Benign (0.81), Benign (1.00); AutoGluon (selected mask): Malignant (0.661)
\end{itemize}
\noindent\textbf{Rec-06.}\par
\begin{itemize}[leftmargin=1.5em,topsep=0.1em,itemsep=0.05em]
\item Case ID: Rec-06
\item Task type: Segmentation routing
\item Candidate similarity tier: Similar
\item LLM invoked: Yes
\item Decision mode: LLM mask selection
\item Selected segmentation expert: PKTN-specific
\item Confidence: 0.860
\item Decision note: Candidates: Multi-cohort joint (agreement 0.886); TN3K-specific (agreement 0.916); ThyroidXL-specific (agreement 0.893); PKTN-specific [selected] (agreement 0.917); TN5K-specific (agreement 0.856); Classification anchor: Malignant
\end{itemize}
\noindent\textbf{\textbf{Rec-07}.}\par
\begin{itemize}[leftmargin=1.5em,topsep=0.1em,itemsep=0.05em]
\item Case ID: Rec-07
\item Task type: Segmentation routing
\item Candidate similarity tier: Similar
\item LLM invoked: Yes
\item Decision mode: LLM mask selection
\item Selected segmentation expert: Multi-cohort joint
\item Confidence: 0.860
\item Decision note: Candidates: Multi-cohort joint [selected] (agreement 0.795); TN3K-specific (agreement 0.796); ThyroidXL-specific (agreement 0.776); PKTN-specific (agreement 0.803); TN5K-specific (agreement 0.499); Classification anchor: Benign
\end{itemize}
\noindent\textbf{Rec-08.}\par
\begin{itemize}[leftmargin=1.5em,topsep=0.1em,itemsep=0.05em]
\item Case ID: Rec-08
\item Task type: Classification routing (LLM arbitration when classifiers disagree)
\item Candidate similarity tier: Similar
\item LLM invoked: Yes
\item Decision mode: LLM arbitration
\item Final benign/malignant class: Malignant
\item Confidence: 0.600
\item Decision note: Independent classifiers: 4 independent classifiers: Benign (0.97), Malignant (0.83), Benign (0.93), Benign (1.00); AutoGluon (selected mask): Malignant (0.768)
\end{itemize}
\noindent\textbf{Rec-09.}\par
\begin{itemize}[leftmargin=1.5em,topsep=0.1em,itemsep=0.05em]
\item Case ID: Rec-09
\item Task type: Classification routing (LLM arbitration when classifiers disagree)
\item Candidate similarity tier: Similar
\item LLM invoked: Yes
\item Decision mode: LLM arbitration
\item Final benign/malignant class: Malignant
\item Confidence: 0.830
\item Decision note: Independent classifiers: 4 independent classifiers: Malignant (1.00), Malignant (1.00), Malignant (0.93), Benign (1.00); AutoGluon (selected mask): Malignant (0.779)
\end{itemize}
\noindent\textbf{Rec-10.}\par
\begin{itemize}[leftmargin=1.5em,topsep=0.1em,itemsep=0.05em]
\item Case ID: Rec-10
\item Task type: Classification routing (deterministic consensus)
\item Candidate similarity tier: Easy
\item LLM invoked: No
\item Decision mode: Deterministic consensus
\item Final benign/malignant class: Benign (consensus)
\item Confidence: 0.998
\item Decision note: Programmatic consensus (Benign (consensus)); adopted without LLM.
\end{itemize}
\par\smallskip\noindent\footnotesize The ten de-identified exemplars above (Rec-01--Rec-10) balance segmentation LLM routing and classification routing (LLM-invoked routes and deterministic consensus, among others). Listed segmentation LLM routes were all ``similar'' tier (top-two candidate agreement gap $\leq 0.10$). Exemplar fields follow Table~A1-2.

\par
\end{responseappendixprose}
\FloatBarrier
\responseproseappendixstart
\subsection*{Appendix A2. Radiomics / AutoGluon specification}
\label{app:A2}
\begin{responseappendixprose}
\noindent This appendix describes where radiomics sits in the agent pipeline, PyRadiomics and AutoGluon-Tabular versions and hyperparameters, and how radiomics input differs from connected-component post-processing on DINO expert/ensemble masks.
\par
\noindent 1.~Pipeline placement
\begin{enumerate}[leftmargin=1.6em,topsep=0.2em,itemsep=0.15em]
\item Segmentation routing yields the router-final nodule mask.
\item No radiomics-specific connected-component post-processing; inference and training share the same 2D PyRadiomics settings (Table A2-1).
\item PyRadiomics features → AutoGluon-Tabular inference; malignancy probability is stored in the evidence bank and used for routing and classification fusion.
\item Optional SHAP feature attribution for the tabular model.
\end{enumerate}
\noindent 2.~PyRadiomics
\par
\begin{responseinlinetable}
\responseinlinetablecaption{Table A2-1. PyRadiomics feature extraction settings}
\label{tab:A2-pyradiomics}
\begin{adjustbox}{max width=\textwidth,center}
\begin{tabular}{ll}
\toprule
Item & Setting \\
\midrule
Package version & \texttt{pyradiomics 3.1.0} \\
Parameter set & Shared 2D PyRadiomics configuration for training and inference \\
Geometry & Forced 2D; mask correction; resampling spacing (1.0, 1.0) px (isotropic assumption) \\
Intensity & Normalization scale 100; histogram binWidth 25; symmetrical GLCM \\
Image filters & Original, Wavelet, LoG ($\sigma=2$, 4), Square, SquareRoot, Logarithm, Exponential, Gradient \\
Feature classes & Full shape2D; GLCM and preset texture families \\
Input & Grayscale ultrasound; mask foreground = positive pixels \\
Feature dimension & 100 radiomics columns at extraction; 86 numeric features used by AutoGluon (14 sparse glszm columns dropped internally) \\
\bottomrule
\end{tabular}
\end{adjustbox}
\end{responseinlinetable}

\noindent 3.~AutoGluon-Tabular
\par
\begin{responseinlinetable}
\responseinlinetablecaption{Table A2-2. AutoGluon-Tabular training and deployment settings}
\label{tab:A2-autogluon}
\begin{adjustbox}{max width=\textwidth,center}
\begin{tabular}{ll}
\toprule
Item & Setting \\
\midrule
Package version & \texttt{1.4.0} \\
Preset & best\_quality \\
Bagging & 8-fold, 1 stack level \\
Optimisation metric & Binary ROC-AUC \\
Training data & Classification stacked dataset (\responsetableref{tab_baseline_protocol}; benign/malignant labels) \\
Training masks & Reference masks on the same stacked set \\
Deployed model & AutoGluon weighted ensemble (bagged LightGBM, CatBoost, neural nets; auto-selected on development data) \\
External inference & Frozen ensemble from development training; no updating on external tests \\
\bottomrule
\end{tabular}
\end{adjustbox}
\end{responseinlinetable}

\noindent 4.~Distinction from expert-mask connected-component post-processing
\par
\noindent Connected-component rules apply to DINO expert/ensemble outputs only (0.5 binarization; retain components $\geq$50 px; 5$\times$5 closing; largest component on ensemble masks; single ROI, no per-lesion multi-nodule separation), not as an extra step before radiomics.
\par
\end{responseappendixprose}

\responseproseappendixstart
\subsection*{Appendix A3. Templates removed by manual review}
\label{app:A3}
% Converted from S4_Template_construction/reviewer_response_package_20261119_claude_v13/tables/Table_R2_19_excluded_templates.md (2026-11-17)
\begin{responsepaginatedtable}
\begin{longtable}{@{}>{\raggedright\arraybackslash}p{0.110\textwidth}>{\raggedright\arraybackslash}p{0.100\textwidth}>{\raggedright\arraybackslash}p{0.340\textwidth}>{\raggedright\arraybackslash}p{0.390\textwidth}@{}}
\caption*{\responsetabletitle{tab_rep_excluded_templates}{Templates removed by manual review from the 400 most frequent gland templates and the 40 most frequent measurement templates of the SMU-HMC library}}\\
\toprule
\textbf{Library} & \textbf{Frequency in source reports} & \textbf{Template (original)} & \textbf{Template (English translation)} \\
\midrule
\endfirsthead
\multicolumn{4}{@{}r@{}}{\footnotesize\emph{\ResponseTableName~\ref{tab_rep_excluded_templates}\responsecontinuedsuffix}} \\
\toprule
\textbf{Library} & \textbf{Frequency in source reports} & \textbf{Template (original)} & \textbf{Template (English translation)} \\
\midrule
\endhead
Gland & 32 & \zh{甲状腺术后。} & Post-thyroidectomy. \\
Gland & 16 & \zh{甲状腺部分切除术后。} & Post partial thyroidectomy. \\
Gland & 5 & \zh{甲状腺术后：左侧残留甲状腺大小约<SIZE3>，右侧残留甲状腺大小约<SIZE3>。} & Post-thyroidectomy: residual left thyroid about <SIZE3>, residual right thyroid about <SIZE3>. \\
Gland & 4 & \zh{甲状腺术后：甲状腺左叶大小:<SIZE3>，右叶大小:<SIZE3>。} & Post-thyroidectomy: thyroid left lobe: <SIZE3>, right lobe: <SIZE3>. \\
Gland & 3 & \zh{左侧残留甲状腺大小:<SIZE3>，右侧残留甲状腺大小:<SIZE3>。} & Residual left thyroid: <SIZE3>, residual right thyroid: <SIZE3>. \\
Gland & 2 & \zh{甲状腺术后：左侧残留甲状腺大小:<SIZE3>，右侧残留甲状腺大小:<SIZE3>，峡部厚<NUM>mm。} & Post-thyroidectomy: residual left thyroid: <SIZE3>, residual right thyroid: <SIZE3>, isthmus <NUM> mm thick. \\
Gland & 2 & \zh{甲状腺部分切除术后：甲状腺左叶大小:<SIZE3>，右叶大小:<SIZE3>。} & Post partial thyroidectomy: thyroid left lobe: <SIZE3>, right lobe: <SIZE3>. \\
Gland & 2 & \zh{左侧残留甲状腺大小约<SIZE3>，右侧残留甲状腺大小约<SIZE3>。} & Residual left thyroid about <SIZE3>, residual right thyroid about <SIZE3>. \\
Gland & 2 & \zh{甲状腺CA术后。} & Post-thyroidectomy for thyroid cancer. \\
Gland & 252 & \zh{甲状腺左叶<SIZE3>，右叶<SIZE3>，峡部<NUM>mm。} & Thyroid left lobe <SIZE3>, right lobe <SIZE3>, isthmus <NUM> mm. \\
Gland & 40 & \zh{甲状腺左叶大小约}<NUM>mm$\times$<NUM>mm$\times$<NUM>mm\zh{，右叶大小约}<NUM>mm$\times$<NUM>mm$\times$<NUM>mm\zh{，峡部厚约<NUM>mm。} & Thyroid left lobe about <NUM> mm $\times$ <NUM> mm $\times$ <NUM> mm, right lobe about <NUM> mm $\times$ <NUM> mm $\times$ <NUM> mm, isthmus about <NUM> mm thick. \\
Gland & 37 & \zh{甲状腺左叶大小:<SIZE2>，右叶大小:<SIZE2>，峡部厚<NUM>mm。} & Thyroid left lobe: <SIZE2>, right lobe: <SIZE2>, isthmus <NUM> mm thick. \\
Gland & 19 & \zh{甲状腺左叶大小:<SIZE3>，右叶大小:<SIZE3>。} & Thyroid left lobe: <SIZE3>, right lobe: <SIZE3>. \\
Gland & 9 & \zh{甲状腺左叶约<SIZE3>，右叶大小约<SIZE3>，峡部约<NUM>mm。} & Thyroid left lobe about <SIZE3>, right lobe about <SIZE3>, isthmus about <NUM> mm. \\
Gland & 8 & \zh{甲状腺左侧叶大小:<SIZE3>，右侧叶大小:<SIZE3>，峡部厚<NUM>mm。} & Thyroid left lobe: <SIZE3>, right lobe: <SIZE3>, isthmus <NUM> mm thick. \\
Gland & 4 & \zh{甲状腺左叶<SIZE3>，右叶<SIZE3>，峡部mm。} & Thyroid left lobe <SIZE3>, right lobe <SIZE3>, isthmus mm. \\
Gland & 4 & \zh{甲状腺右叶大小:<SIZE3>，左叶大小:<SIZE3>。} & Thyroid right lobe: <SIZE3>, left lobe: <SIZE3>. \\
Gland & 2 & \zh{甲状腺左叶大小约<SIZE3>，右叶大小约<SIZE3>。} & Thyroid left lobe about <SIZE3>, right lobe about <SIZE3>. \\
Gland & 107 & CDFI\zh{：腺体内见少量纤细分枝状彩色血流信号。} & CDFI: a small amount of fine branching colour flow signal in the gland. \\
Gland & 26 & CDFI\zh{：腺体内可探及较丰富血流信号。} & CDFI: relatively abundant flow signal in the gland. \\
Gland & 13 & CDFI\zh{：腺体内见血流信号轻度增多，呈点状或分枝状。} & CDFI: mildly increased flow signal in the gland, dotted or branching. \\
Gland & 7 & CDFI\zh{：腺体内见血管轻度增多，呈点状或分枝状。} & CDFI: mildly increased vessels in the gland, dotted or branching. \\
Gland & 3 & CDFI\zh{：腺体内血管弥漫性明显增多。} & CDFI: diffusely and markedly increased vessels in the gland. \\
Gland & 3 & CDFI\zh{于甲状腺腺体内可见极丰富血流信号显示，呈}``\zh{火海样}''\zh{。} & CDFI: very abundant flow signal in the thyroid gland, with a ``thyroid inferno'' pattern. \\
Gland & 3 & \zh{腺体内血管弥漫性明显增多。} & Diffusely and markedly increased vessels in the gland. \\
Gland & 2 & CDFI\zh{：腺体内可见较丰富血流信号显示，呈}``\zh{火海样}''\zh{。} & CDFI: relatively abundant flow signal in the gland, with a ``thyroid inferno'' pattern. \\
Gland & 2 & CDFI\zh{：腺体内血流弥漫性明显增多，呈火海样。} & CDFI: diffusely and markedly increased flow in the gland, with a thyroid inferno pattern. \\
Gland & 2 & CDFI\zh{：腺体内血流信号弥漫性明显增多。} & CDFI: diffusely and markedly increased flow signal in the gland. \\
Gland & 2 & CDFI\zh{：腺体内血流信号稍丰富。} & CDFI: slightly abundant flow signal in the gland. \\
Gland & 2 & CDFI\zh{：腺体内可探及丰富血流信号。} & CDFI: abundant flow signal in the gland. \\
Gland & 2 & CDFI\zh{：腺体内血流信号弥漫性增多。} & CDFI: diffusely increased flow signal in the gland. \\
Gland & 10 & \zh{双侧甲状腺内可见多个无回声，大小约<SIZE2>、<SIZE2>、<SIZE2>等，内均可见一强回声光点，CDFI：其内未见明确血流信号。} & Multiple anechoic lesions in both lobes, about <SIZE2>, <SIZE2>, <SIZE2> and others, each containing an echogenic dot; CDFI: no definite flow signal within. \\
Gland & 6 & \zh{甲状腺双侧叶见多个无回声，较大约<SIZE2>，边界清，暗区内见点状强回声，CDFI：其内未见血流信号。} & Multiple anechoic lesions in both lobes, the larger about <SIZE2>, well defined, with punctate echogenic foci in the anechoic area; CDFI: no flow signal within. \\
Gland & 4 & \zh{双侧甲状腺内可见多个无回声，大小约<SIZE2>、<SIZE2>、<SIZE2>，内均可见一强回声光点，CDFI：其内未见明确血流信号。} & Multiple anechoic lesions in both lobes, about <SIZE2>, <SIZE2> and <SIZE2>, each containing an echogenic dot; CDFI: no definite flow signal within. \\
Gland & 4 & \zh{甲状腺右侧叶见多个无回声，较大约<SIZE2>，边界清，暗区内见点状强回声，CDFI：其内未见血流信号。} & Multiple anechoic lesions in the right lobe, the larger about <SIZE2>, well defined, with punctate echogenic foci in the anechoic area; CDFI: no flow signal within. \\
Gland & 4 & \zh{双侧甲状腺内可见多个无回声，大小约<SIZE2>、<SIZE2>，内均可见一强回声光点，CDFI：其内未见明确血流信号。} & Multiple anechoic lesions in both lobes, about <SIZE2> and <SIZE2>, each containing an echogenic dot; CDFI: no definite flow signal within. \\
Gland & 2 & \zh{甲状腺左叶另见强回声，大小约<SIZE2>。} & Another echogenic focus in the left lobe, about <SIZE2>. \\
Gland & 2 & \zh{双侧甲状腺见无回声，左侧较大者约<SIZE2>，右侧较大者约<SIZE2>，边界清，CDFI：其内未见血流信号。} & Anechoic lesions in both lobes, the larger on the left about <SIZE2> and on the right about <SIZE2>, well defined; CDFI: no flow signal within. \\
Gland & 2 & \zh{甲状腺右叶见实质性光团，大小约<SIZE2>，CDFI：其内未见血流信号。} & A solid mass in the right lobe, about <SIZE2>; CDFI: no flow signal within. \\
Gland & 2 & \zh{甲状腺右叶探及实质性光团，大小约<SIZE2>、<SIZE2>，CDFI：其内未探及血流信号。} & Solid masses in the right lobe, about <SIZE2> and <SIZE2>; CDFI: no flow signal detected within. \\
Gland & 2 & \zh{甲状腺左叶中部见强回声，大小约<SIZE2>。} & An echogenic focus in the middle of the left lobe, about <SIZE2>. \\
Gland & 2 & \zh{甲状腺右叶中部见强回声后伴声影，大小约<SIZE2>。} & An echogenic focus with posterior acoustic shadowing in the middle of the right lobe, about <SIZE2>. \\
Measurement & 17 & \zh{急诊超声检查：甲状腺左叶大小约<SIZE3>，右叶大小约<SIZE3>，峡部厚约<NUM>mm。} & Emergency ultrasound: thyroid left lobe about <SIZE3>, right lobe about <SIZE3>, isthmus about <NUM> mm thick. \\
Measurement & 14 & \zh{急诊超声检查：甲状腺左叶大小:<SIZE3>，右叶大小:<SIZE3>，峡部厚<NUM>mm。} & Emergency ultrasound: thyroid left lobe: <SIZE3>, right lobe: <SIZE3>, isthmus <NUM> mm thick. \\
Measurement & 11 & \zh{床旁超声检查，图像未能留存，结果仅供参考：甲状腺左叶大小约<SIZE3>，右叶大小约<SIZE3>，峡部厚约<NUM>mm。} & Bedside ultrasound, images not stored, results for reference only: thyroid left lobe about <SIZE3>, right lobe about <SIZE3>, isthmus about <NUM> mm thick. \\
Measurement & 6 & \zh{急诊彩超：甲状腺左叶大小:<SIZE3>，右叶大小:<SIZE3>，峡部厚<NUM>mm。} & Emergency colour Doppler ultrasound: thyroid left lobe: <SIZE3>, right lobe: <SIZE3>, isthmus <NUM> mm thick. \\
Measurement & 4 & \zh{床旁超声检查，图像未能留存，结果仅供参考：甲状腺左叶大小:<SIZE3>，右叶大小:<SIZE3>，峡部厚<NUM>mm。} & Bedside ultrasound, images not stored, results for reference only: thyroid left lobe: <SIZE3>, right lobe: <SIZE3>, isthmus <NUM> mm thick. \\
Measurement & 3 & \zh{急诊超声检查：甲状腺右叶大小:<SIZE3>，左叶大小:<SIZE3>，峡部厚<NUM>mm。} & Emergency ultrasound: thyroid right lobe: <SIZE3>, left lobe: <SIZE3>, isthmus <NUM> mm thick. \\
Measurement & 2 & \zh{急诊超声检查：甲状腺左叶大小约<SIZE3>，右叶大小约<SIZE3>，峡部<NUM>mm。} & Emergency ultrasound: thyroid left lobe about <SIZE3>, right lobe about <SIZE3>, isthmus <NUM> mm. \\
Measurement & 8 & \zh{甲状腺部分切除术后：甲状腺左叶大小:<SIZE3>，右叶大小:<SIZE3>，峡部厚<NUM>mm。} & Post partial thyroidectomy: thyroid left lobe: <SIZE3>, right lobe: <SIZE3>, isthmus <NUM> mm thick. \\
Measurement & 4 & \zh{甲状腺术后：甲状腺左叶大小:<SIZE3>，右叶大小:<SIZE3>，峡部厚<NUM>mm。} & Post-thyroidectomy: thyroid left lobe: <SIZE3>, right lobe: <SIZE3>, isthmus <NUM> mm thick. \\
Measurement & 3 & \zh{甲状腺术后：甲状腺左叶大小约<SIZE3>，右叶大小约<SIZE3>，峡部厚约<NUM>mm。} & Post-thyroidectomy: thyroid left lobe about <SIZE3>, right lobe about <SIZE3>, isthmus about <NUM> mm thick. \\
Measurement & 2 & \zh{甲状腺部分切除术后：甲状腺右叶大小:<SIZE3>，左叶大小:<SIZE3>，峡部厚<NUM>mm。} & Post partial thyroidectomy: thyroid right lobe: <SIZE3>, left lobe: <SIZE3>, isthmus <NUM> mm thick. \\
Measurement & 16 & \zh{甲状腺左叶大小:<SIZE3>，右叶大小:<SIZE3>，峡部厚mm。} & Thyroid left lobe: <SIZE3>, right lobe: <SIZE3>, isthmus mm thick. \\
Measurement & 3 & \zh{甲状腺左叶大小约<SIZE3>，右叶大小约<SIZE3>，峡部厚mm。} & Thyroid left lobe about <SIZE3>, right lobe about <SIZE3>, isthmus mm thick. \\
Measurement & 2 & \zh{甲状腺左叶大小}:$\times$<SIZE2>\zh{，右叶大小:<SIZE3>，峡部厚<NUM>mm。} & Thyroid left lobe: $\times$ <SIZE2>, right lobe: <SIZE3>, isthmus <NUM> mm thick. \\
Measurement & 1 & \zh{（甲状腺体积过大，测量欠满意）甲状腺左叶大小:<SIZE3>，右叶大小:<SIZE3>，峡部厚<NUM>mm。} & (Thyroid too large; measurement unsatisfactory) Thyroid left lobe: <SIZE3>, right lobe: <SIZE3>, isthmus <NUM> mm thick. \\
\bottomrule
\end{longtable}
\responsetablenote{Templates keep their slot placeholders: <SIZE2>, two-dimensional size; <SIZE3>, three-dimensional size; <NUM>, number. Frequency in source reports: the number of clauses in the source reports mapped to each template.}
\end{responsepaginatedtable}


\responseproseappendixstart
\subsection*{Appendix A4. Data flow}
\label{app:A4}
% Converted from S9_Data_flow/reviewer_response_package_20261120_claude_v7/figures/FIGURE_CAPTION_EN.md and S9_Data_flow/reviewer_response_package_20261120_claude_v7/figures/Fig_dataflow_image_analysis.pdf (2026-11-17)
\begin{responsefigure}
\includegraphics[width=\textwidth,height=0.78\textheight,keepaspectratio]{report_rep_flow_image.pdf}
\caption{\responsefiguretitle{fig_rep_flow_image}{Data flow of the image-analysis arm. All counts are images. The training, validation and internal-test sets were drawn from TN3K, TN5K, ThyroidXL and PKTN; the validation set was used to select checkpoints, fit the router and set the referral thresholds, which were then frozen for testing. The external-test cohorts were not used for model development. ZJH-8K served as an external test set for classification and segmentation, whereas its malignant primary lesions were split at the patient level to develop (with five-fold cross-validation) and test the LNM and PTC subtype models.}}
\label{fig_rep_flow_image}
\end{responsefigure}

% Converted from S9_Data_flow/reviewer_response_package_20261120_claude_v7/figures/FIGURE_CAPTION_EN.md and S9_Data_flow/reviewer_response_package_20261120_claude_v7/figures/Fig_dataflow_report_generation.pdf (2026-11-17)
\begin{responsefigure}
\includegraphics[width=\textwidth,height=0.78\textheight,keepaspectratio]{report_rep_flow_report.pdf}
\caption{\responsefiguretitle{fig_rep_flow_report}{Data flow of the report-generation arm. For the auxiliary tools, SMU-HMC data are counted in reports and the other datasets in images (frames for TGVideo); the training, validation and test splits of each tool are given in Supplementary Table S13. The SMU-HMC template library was built from reports with a completeness score of at least 0.9. The KMVE library excluded 406 training reports that were near-duplicates of test reports, and its rules were checked on the validation split. ZJH-TS reports were generated with the SMU-HMC library. The yellow safety rules were selected on the agent's reports for the three test sets and evaluated on TNVideo.}}
\label{fig_rep_flow_report}
\end{responsefigure}

